\chapter{Technische Umsetzung}
\label{chap:technische_umsetzung}

Die Grundlagen und Erklärungen dieses Kapitels wurden in den vorherigen Kapiteln bereits behandelt. 
Dabei wurden die Datenobjekte Makroprofile, Anfrageprofile, Anfrage-Kandidat-Paare sowie die Vergleichsmerkmale und die Merkmalsvektoren in Kapitel~\ref{chap:konzeption_des_empfehlungssystems} erläutert.
Die Grundlagen der Datenaufbereitung und des Datensplits wurden im Kapitel~\ref{chap:datenaufbereitung_und_datensplit} behandelt.
% \section{Implementierung der Datenobjekte}
% \label{sec:implementierung_datenobjekte}

% Die fachlichen Grundlagen der Datenobjekte: Makroprofile, Anfrageprofile und Anfrage-Kandidat-Paare wurden bereits in Kapitel~\ref{chap:konzeption_des_empfehlungssystems} beschrieben. 
% In der technischen Umsetzung werden diese Datenobjekte aus den aufbereiteten Artikeln des jeweiligen Datensplits und den zugeordneten Makros erzeugt.
% Es werden ausschließlich die technischen Artikel-und Makroeigenschaften verwendet, die aus einem technischen Datenblatt gewonnen werden können. Hierzu gehören beispielsweise Hersteller, Spannung, Strom, Leistung, Anschlüsse, und technische Hinweise. 
% EPLAN-interne Pfade, Bearbeiterangaben und Änderungszeitpunkte werden ausgeschlossen. Diese Informationen beschreiben die Ablage oder Bearbeitung der Daten, nicht jedoch  selbst. Würden sie verwendet, 
% könnte das Modell das passende Makro an einem Dateipfad oder einer internen Kennung erkennen, obwohl diese Angaben bei einer späteren Datenblattanfrage nicht vorhanden sind.

% Für das Motor-Makro \texttt{M1\_F1\_V1} werden alle in den zugeordneten Trainingsartikeln beobachteten Werte in das Makroprofil übernommen. 
% Besitzen die Motorartikel beispielsweise die Spannungen \(\{380,400,420\}\,\mathrm{V}\), die Ströme \(\{10,12,15\}\,\mathrm{A}\) und die Leistungen \(\{5,7\}\,\mathrm{kW}\), bleiben alle Ausprägungen erhalten. 
% Das Makroprofil kann damit wie ein gemeinsames technisches Sammelblatt des Motor-Makros betrachtet werden. Es zeigt nicht nur einen einzelnen Motor, sondern den gesamten Wertebereich der Motorartikel, die mit diesem Makro dargestellt werden können.

% Hierfür wird die Konfiguration \texttt{all\_observed\_train\_values} verwendet. Sie bedeutet, dass jeder in den Trainingsdaten beobachtete Wert übernommen wird. Eine Top-N-Begrenzung findet nicht statt. Dies ist notwendig, 
% weil ein selten vorkommender Wert wie 420~V trotzdem eine gültige Motorvariante darstellen kann. Würden nur die häufigsten Werte gespeichert, 
% könnte diese Variante aus dem Makroprofil verschwinden und bei einer späteren Anfrage nicht mehr vollständig verglichen werden.

% Die Konfiguration wurde außerdem gewählt, weil die vorhandenen Makros unterschiedlich viele Trainingsartikel besitzen. Eine feste Grenze von beispielsweise 8 oder 16 Werten würde große Makroprofile stärker beschneiden als kleine Makroprofile. 
% Beim Motor-Makro könnte dadurch ein selten beobachteter Spannungs- oder Leistungswert entfallen, obwohl er fachlich gültig ist. Die vollständige Übernahme vermeidet diese ungleiche Behandlung. Die Häufigkeit eines Wertes bleibt für Auswertungen erhalten, 
% entscheidet jedoch nicht darüber, ob der Wert Teil des Makroprofils wird. Die Werte werden deterministisch sortiert. Das bedeutet, dass dieselben Eingangsdaten immer dieselbe Reihenfolge im gespeicherten Makroprofil erzeugen. 
% Dies erleichtert den Vergleich versionierter Artefakte und den reproduzierbaren nächtlichen Pipeline-Lauf.

% Das Anfrageprofil beschreibt dagegen nur den Motorartikel der jeweiligen Anfrage. Es enthält dessen technische Artikeleigenschaften sowie solche technischen Makroeigenschaften, die ebenfalls in einem Datenblatt vorkommen können. 
% Für den Beispielmotor verbleiben damit unter anderem SEW, 400~V, 10~A, 4000~W, 50~Hz und der technische Hinweis \textit{Mit Kühlung}. Die Makrokennung \texttt{M1\_F1\_V1} und die Artikel-Makro-Zuordnung werden entfernt. 
% Das Anfrageprofil kann somit wie ein Motorsteckbrief ohne aufgedruckte Lösung verstanden werden: Es beschreibt, welcher Motor vorliegt, verrät dem Modell jedoch nicht, welches Makro gesucht wird.

% Die datenblattkompatiblen Informationen aus dem zugeordneten Makro werden auch im Anfrageprofil verwendet, weil sie fachliche Angaben ergänzen können, die im Artikelbestand uneinheitlich verteilt sind. Übernommen werden jedoch nur Angaben, 
% die mit derselben Extraktionslogik auch aus einem Datenblatt gebildet werden können. Beim Motor betrifft dies beispielsweise eine technische Beschreibung zur Kühlung, zur Anschlussart oder zur elektrischen Ausführung. 
% Nicht übernommen werden die Anzahl der EPLAN-Hauptobjekte, Projektschritte, Varianten, Symbolbibliotheken oder Ablagepfade. Diese Angaben beschreiben den Aufbau der Makrodatei und nicht den Motor. Die Trennung stellt sicher, 
% dass Training und spätere Datenblattanfrage dieselbe fachliche Informationsgrundlage verwenden.

% Anschließend wird das Anfrageprofil mit jedem Makroprofil kombiniert. Für das Motorbeispiel entstehen unter anderem die Anfrage-Kandidat-Paare \((q_{Motor},M1\_F1\_V1)\), \((q_{Motor},M1\_F1\_V2)\) und \((q_{Motor},M1\_F2\_V1)\). 
% Ein \textit{Anfrage-Kandidat-Paar} ist damit ein einzelner Vergleich zwischen dem Motorsteckbrief und einem möglichen Motor-Makro.

% Für jedes Paar werden Vergleichsmerkmale berechnet. Ein \textit{Vergleichsmerkmal} beschreibt eine einzelne Gemeinsamkeit oder Abweichung. 
% Stimmen in Anfrage und Makroprofil beispielsweise 400~V überein, entsteht eine hohe Spannungsübereinstimmung. Enthält ein anderes Makroprofil dagegen nur 230~V, entsteht ein größerer numerischer Abstand. 
% Entsprechend werden auch Hersteller, Strom, Leistung, Frequenz und technische Hinweise verglichen.

% Die Vergleichsergebnisse werden zu einem Merkmalsvektor zusammengeführt. Ein \textit{Merkmalsvektor} ist in diesem Zusammenhang eine geordnete Zahlenreihe. Jede Stelle beschreibt ein festgelegtes Vergleichsmerkmal. 
% Bildlich entspricht der Merkmalsvektor einer ausgefüllten Prüfliste: Ein Zahlenwert beschreibt die Spannungsübereinstimmung, ein weiterer die Stromübereinstimmung und ein weiterer die Übereinstimmung des technischen Hinweises. 
% Diese Zahlenreihe bildet die Eingabe des LambdaMART-Modells.

\section{Implementierung des LambdaMART-Modells}
\label{sec:implementierung_des_lambdamart_modells}

\subsection{Merkmalsmatrix und Labelvektoren}
\label{sec:merkmalsmatrix_und_labelvektoren}

Die in Abschnitt~\ref{sec:sec:anfrage_kandidat_paar_vergleichsmerkmale_und_merkmalsvektoren} erklärten Anfrage-Kandidat-Paare lassen sich aus Abbildung\ref{fig:anfrage_kandidat_paare_motor} in Form einer so genannten Merkmalsmatrix darstellen.
Diese Merkmalsmatrix wird für das Training, die Validierung und den Test des LambdaMART-Modells verwendet und jeweils als Trainingsmatrix \(X_{Train}\), Validierungsmatrix \(X_{Val}\) und Testmatrix \(X_{Test}\) bezeichnet.
Die Trainingsmatrix dient dem Aufbau des Modells. Die Validierungsmatrix dient der Auswahl der Modellkonfiguration. Die Testmatrix dient ausschließlich der abschließenden Bewertung.
Eine Merkmalsmatrix wird in Zeilen und Spalten organisiert. Jede Zeile entspricht dabei dem Merkmalsvektoreines Anfrage-Kandidat-Paares. In jeder Spalte stehen die berechneten Vergleichsmerkmale. 
Die Trainingsmatrix enthält die Merkmalsvektoren der Trainingsanfragen. Entsprechend enthalten die Validierungsmatrix und die Testmatrix die Merkmalsvektoren der Validierungs- und Testanfragen. 
Dadurch können beliebig viele Merkmalsvektoren mit gleichem Aufbau der Vergleichsmerkmale der Merkmalsmatrix hinzugefügt werden, und das LambdaMART-Modell verarbeitet die Vergleichsmerkmale in allen Datensplits mit derselben Bedeutung. 
Eine Matrixzeile lässt sich damit wie folgt darstellen: \[ \text{Anfrageprofil} + \text{Makroprofil} \rightarrow \text{Anfrage-Kandidat-Paar} \rightarrow \text{Vergleichsmerkmale} \rightarrow \text{Merkmalsvektor}. \] 
Für das Training werden jeder Matrixzeile zusätzlich ein Labelvektor zugeordnet, welcher die Relevanz des Anfrage-Kandidat-Paares dieser Zeile beschreibt.
Die Labelvektoren werden entsprechend ihres Verwendungszwecks als Trainingslabelvektor \(y_{Train}\), 
Validierungslabelvektor \(y_{Val}\) sowie Testlabelvektor \(y_{Test}\) bezeichnet. 

Die Relevanz eines Anfrage-Kandidat-Paares im Labelvektor wird dabei durch eine numerische Bewertung ausgedrückt. Diese Bewertung erfolgt in 2 Varianten: binär und hierarchisch.
Diese Bewertung ergibt sich aus der Makrobezeichnung und soll anahnd des Beispiels M1\_F1\_V1 erläutert werden. Sie wird in folgender Skala dargestellt: 

Beispiel: M1\_F1\_V1
\begin{itemize}
    \item M1: Makro-Klasse
    \item F1: Makro-Familie
    \item V1: Makro-Variante
\end{itemize}

Binäre Bewertung:

\begin{itemize}
    \item 0: keine Übereinstimmung
    \item 1: exakte Übereinstimmung zwischen Makroprofil und Makrobezeichnung
\end{itemize}

Hierarchische Bewertung:

\begin{itemize}
    \item 0: keine Übereinstimmung
    \item 1: Übereinstimmung der Makro-Klasse zwischen Makroprofil und Makrobezeichnung
    \item 2: Übereinstimmung der Makro-Familie zwischen Makroprofil und Makrobezeichnung
    \item 4: exakte Übereinstimmung zwischen Makroprofil und Makrobezeichnung
\end{itemize}
 
Im Gegensatz zur binären Bewertung wird bei der hierarchischen Bewertung berücksichtigt, dass ein Motor-Makro derselben Familie fachlich näher am Makroprofil liegt als ein nicht verwandtes Makro.
Welche Variante für den vorhandenen Datenbestand geeigneter ist, wird im Validierungsprozess bestimmt.

% Bild

\subsection{Anfrage an das LambdaMART-Modell}
\label{sec:anfrage_an_das_lambdamart_modell}

Um eine Anfrage \(q\) an das LambdaMART-Modell zu übergeben zu können, wird das Anfrageprofil von \(q\) mit allen Makroprofilen aus dem Suchraum\_Makroprofile kombiniert, 
um dadurch für jedes Anfrage-Kandidat-Paar von \(q\) einen Merkmalsvektor für die Merkmalsmatrix zu berechnen.
Somit entsteht für \(q\) eine Menge von Anfrage-Kandidat-Paaren, welche die die Übereinstimmung aller Makroprofile mit dem Anfrageprofil von \(q\) beschreibt.
Damit die Makroprofile durch das LambdaMART-Modell der Anfrage \(q\) zugeordnet werden können, werden alle Anfrage-Kandidat-Paare von \(q\) in einer sogenannten \textit{Rankinggruppe} zusammengefasst.
Damit weiß das LambdaMART-Modell später welche Anfrage-Kandidat-Paare für das Ranking innerhalb von \(q\) berücksichtigt werden müssen. Daraus ergibt sich eine eindeutige Zuordnung einer Matrixzeile zu der Anfrage \(q\). 

% Bild





\subsection{Training des LambdaMART-Modells} 
\label{sec:training_lambdamart_modells} 

Der \texttt{LGBMRanker} lernt anhand der Trainingsmatrix, welche Vergleichsmerkmale für die Reihenfolge der Makrokandidaten relevant sind. Hierzu verwendet das LambdaMART-Modell mehrere aufeinanderfolgende Entscheidungsbäume. 
Ein \textit{Entscheidungsbaum} prüft die Werte eines Merkmalsvektors schrittweise anhand einfacher Bedingungen. Für den Beispielmotor kann dabei zunächst geprüft werden, ob die Spannungsübereinstimmung einen festgelegten Wert überschreitet. 
Anschließend kann eine weitere Bedingung die Übereinstimmung der Leistung oder des technischen Hinweises berücksichtigen. Am Ende des Entscheidungsweges entsteht ein Zahlenwert für das zugehörige Anfrage-Kandidat-Paar. 
Beim \textit{Gradient-Boosting} werden mehrere Entscheidungsbäume nacheinander aufgebaut. Der erste Baum erzeugt eine erste Bewertung der Makrokandidaten. 
Die folgenden Bäume berücksichtigen die noch vorhandenen Fehlplatzierungen und ergänzen schrittweise Korrekturen. Bildlich erstellt der erste Baum einen ersten Entwurf der Makro-Rangliste. 
Jeder weitere Baum prüft diesen Entwurf und verbessert die Reihenfolge an den Stellen, an denen relevante Kandidaten noch zu niedrig eingeordnet sind. Die Beiträge aller Entscheidungsbäume werden zu einem Relevanzscore zusammengeführt. 
Der \textit{Relevanzscore} ist der vom Modell berechnete Zahlenwert eines Anfrage-Kandidat-Paares. Für jede Zeile der Trainings-, Validierungs-, Test- oder Anwendungsmatrix entsteht genau ein Relevanzscore. Er besitzt keine technische Einheit, 
sondern dient ausschließlich dem Vergleich der Makrokandidaten innerhalb derselben Rankinggruppe. Für den Beispielmotor kann \texttt{M1\_F1\_V1} einen höheren Relevanzscore als \texttt{M1\_F1\_V2} erhalten, 
weil dessen Makroprofil stärker mit der Spannung, dem Strom, der Leistung, der Frequenz und dem technischen Hinweis des Anfrageprofils übereinstimmt. 
Die Makro-Rangliste entsteht anschließend durch die absteigende Sortierung der Relevanzscores innerhalb der Rankinggruppe. Der Kandidat mit dem höchsten Relevanzscore steht auf Rang 1. Eine nachträgliche regelbasierte Sortierung findet nicht statt. 
Insbesondere werden Kandidaten innerhalb der bestbewerteten Makrofamilie nicht zusätzlich anhand einzelner Typ-, Suffix- oder Beschreibungstreffer verschoben. Eine solche Feinordnung würde die vom Modell erzeugte Rangfolge verändern. 
Recall@k, MRR und nDCG würden dann nicht mehr ausschließlich die Leistung des LambdaMART-Modells bewerten. Die unveränderten Relevanzscores ermöglichen dagegen einen eindeutigen Vergleich der Modellkonfigurationen. 
Bei identischen Relevanzscores wird lediglich stabil sortiert. Eine \textit{stabile Sortierung} behält bei gleichen Werten die bereits vorhandene Reihenfolge der Kandidaten bei. Sie verändert weder den Relevanzscore noch die fachliche Bewertung, 
sorgt jedoch dafür, dass ein wiederholter Pipeline-Lauf dieselbe Makro-Rangliste erzeugt.

Der vollständige Trainingseingang besteht damit aus: \[ X_{Train} = \text{Merkmalsvektoren der Anfrage-Kandidat-Paare}, \] \[ y_{Train} = \text{zugehörige Relevanzlabels}, \] \[ group_{Train} = \text{Anzahl der Makrokandidaten je Anfrage}. \] 
Für Validation und Test wird dieselbe Struktur mit \((X_{Val},y_{Val},group_{Val})\) und \((X_{Test},y_{Test},group_{Test})\) verwendet.  

\section{Hyperparameteroptimierung}
\label{sec:hyperparameteroptimierung}

Die Leistungsfähigkeit des LambdaMART-Modells wird durch Hyperparameter beeinflusst. \textit{Hyperparameter} sind Einstellungen, die den Lernprozess steuern und vor dem Training festgelegt werden. 
Sie können mit den Einstellmöglichkeiten einer Maschine verglichen werden: Das Werkstück bleibt gleich, aber Geschwindigkeit, Anzahl der Arbeitsschritte und Genauigkeit der Bearbeitung können unterschiedlich eingestellt werden.

Zur Bestimmung einer geeigneten Modellkonfiguration wird eine Grid Search verwendet. Eine \textit{Grid Search} ist die systematische Erprobung mehrerer festgelegter Kombinationen von Hyperparametern. 
Bildlich entspricht sie einer Reihe gleich aufgebauter Versuche. In jedem Versuch werden dieselben Trainings- und Validierungsdaten verwendet, jedoch andere Einstellungen des LambdaMART-Modells.

Das Grid enthält beispielsweise folgende Konfiguration:

Die Werte des Grids wurden als begrenzter Suchraum gewählt, weil der vorhandene Datenbestand aus einer vergleichsweise kleinen Zahl von Artikel-Makro-Zuordnungen und zahlreichen unterschiedlich großen Makrogruppen besteht. 
Ein sehr großer Suchraum würde viele nahezu gleichartige Versuche erzeugen, die Rechenzeit erhöhen und die Gefahr einer zufälligen Anpassung an die Validierungsdaten vergrößern. Die gewählten Stufen vergleichen deshalb bewusst kleinere, 
mittlere und komplexere Modelle.

\begin{quote}
\texttt{n\_estimators = 300}, \texttt{learning\_rate = 0.05},
\texttt{num\_leaves = 31}, \texttt{min\_child\_samples = 20} und
\texttt{lambdarank\_truncation\_level = 10}.
\end{quote}

\texttt{n\_estimators} gibt die Anzahl der nacheinander erzeugten Entscheidungsbäume an. Bei 300 werden somit 300 einzelne Bäume zu einem gemeinsamen Modell verbunden. Mehr Bäume ermöglichen weitere Korrekturschritte, 
erhöhen jedoch auch den Trainingsaufwand und das Risiko, Besonderheiten der Trainingsdaten auswendig zu lernen.

Die Konfigurationen mit 300, 500, 700 und 900 Bäumen decken dabei schrittweise unterschiedliche Modellgrößen ab. 300 Bäume bilden die kompakte Ausgangskonfiguration. 
500 und 700 Bäume prüfen, ob zusätzliche Korrekturschritte die Abgrenzung ähnlicher Motor-Makros verbessern. 900 Bäume bilden die obere Grenze des untersuchten Bereichs. Eine weitere Erhöhung wurde nicht vorgesehen, 
weil der zusätzliche Rechen- und Speicherbedarf steigt, während bei begrenzten Trainingsdaten das Overfitting-Risiko zunimmt.

Die \texttt{learning\_rate} ist die Lernrate. Sie bestimmt, wie stark jeder neue Baum den bisherigen Relevanzscore verändern darf. Eine Lernrate von 0,05 bedeutet, dass jeder einzelne Korrekturbeitrag nur in begrenzter Stärke übernommen wird. 
Eine zweite Konfiguration verwendet beispielsweise 500 Bäume bei einer Lernrate von 0,04. Bildlich werden damit mehr kleinere Korrekturen durchgeführt, anstatt weniger große Änderungen vorzunehmen. Dies kann stabilere Zusammenhänge zwischen Spannung, 
Strom, Leistung und technischen Hinweisen erzeugen.

Im Grid werden die Lernraten 0,05, 0,04, 0,03 und 0,025 mit einer steigenden Anzahl von Bäumen verbunden. Diese Abstufung hält die gesamte Veränderung des Modells in einem vergleichbaren Bereich: Wenige Bäume dürfen etwas größere Schritte ausführen, 
während viele Bäume kleinere Schritte verwenden. Für den Motor bedeutet dies, dass eine einzelne Spannungseigenschaft das Ranking nicht sofort stark verändert, sondern gemeinsam mit Strom, Leistung, 
Frequenz und technischen Hinweisen über mehrere Korrekturschritte bewertet wird.

\texttt{num\_leaves} bezeichnet die maximale Anzahl der Blätter eines Entscheidungsbaums. Ein Blatt ist das Ende eines Entscheidungsweges. Mehr Blätter ermöglichen feinere Unterscheidungen. 
Mit 63 Blättern kann das Modell beispielsweise mehr Kombinationen aus Spannung, Leistung und Kühlung abbilden als mit 31 Blättern. Bei dem begrenzten Datenbestand steigt dadurch jedoch das Risiko des Overfittings.

Die verwendeten Werte 15, 31 und 63 verdoppeln die mögliche Baumkomplexität annähernd von Stufe zu Stufe. 15 Blätter prüfen, ob wenige robuste Regeln ausreichen. 31 Blätter bilden die mittlere Konfiguration. 
63 Blätter erlauben feinere Abgrenzungen zwischen ähnlichen Motorvarianten. Größere Werte werden nicht verwendet, weil kleine Makrogruppen dann durch sehr spezifische Pfade getrennt werden könnten, die sich nicht auf neue Datenblätter übertragen lassen.

\textit{Overfitting} bedeutet, dass das Modell die Trainingsbeispiele zu genau auswendig lernt. Es verhält sich dann wie ein Prüfer, der nur die bereits bekannten Motoren wiedererkennt, 
aber bei einem neuen Motor mit leicht abweichenden Werten unsicher wird. Deshalb werden sowohl einfachere als auch komplexere Konfigurationen untersucht.

\texttt{min\_child\_samples} legt fest, wie viele Trainingsbeispiele mindestens in einem Blatt enthalten sein müssen. Ein höherer Wert verhindert Entscheidungen, die nur auf sehr wenigen Motorartikeln beruhen. 
Dies ist für den vorhandenen Datenbestand besonders wichtig, da viele Makros nur wenige zugeordnete Artikel besitzen.

Die Stufen 10, 20 und 30 prüfen, wie stark solche Einzelentscheidungen begrenzt werden müssen. Der Wert 10 lässt auch Regeln für kleinere Gruppen zu. Der Wert 20 bildet die ausgeglichene Ausgangseinstellung. 
Der Wert 30 fordert eine breitere Datengrundlage je Blatt und wirkt dadurch stärker gegen Overfitting. Die Werte wurden so gewählt, dass seltene Motorvarianten nicht grundsätzlich ausgeschlossen werden, 
einzelne Artikel aber möglichst keine eigene Entscheidungsregel bilden.

\texttt{lambdarank\_truncation\_level} bestimmt, bis zu welcher Rangposition Fehlplatzierungen während des Trainings besonders berücksichtigt werden. Der Wert 10 richtet den Lernprozess auf den betrachteten Top-10-Bereich aus. 
Dies entspricht dem Ziel, geeignete Makros möglichst innerhalb der ersten zehn Positionen bereitzustellen, und stimmt mit Recall@10 und nDCG@10 überein.

Zusätzlich werden die Werte 5 und 20 untersucht. Der Wert 5 konzentriert das Lernen stärker auf die ersten Plätze, die für eine direkte Empfehlung besonders wichtig sind. 
Der Wert 20 berücksichtigt einen breiteren Ausschnitt der Rangliste und kann dadurch die Reihenfolge fachlich verwandter Makros stabilisieren. 
Der Wert 10 bildet die fachliche Mitte und entspricht unmittelbar dem in der Arbeit ausgewerteten Top-10-Bereich.

Zusätzlich werden binäre und hierarchische Labels sowie moderate und starke Regularisierung verglichen. Eine \textit{Regularisierung} begrenzt unnötig komplexe Entscheidungen des Modells. Sie kann mit einer Vorgabe verglichen werden, 
nur solche Regeln zu verwenden, die sich bei mehreren Motorartikeln wiederholen. Die moderat regularisierte Konfiguration verwendet beispielsweise \texttt{reg\_alpha = 0.10} und \texttt{reg\_lambda = 0.50}. 
Die starke Regularisierung verwendet \texttt{reg\_alpha = 0.50}, \texttt{reg\_lambda = 1.00} und \texttt{min\_split\_gain = 0.01}. Dadurch soll verhindert werden, dass einzelne Ausnahmen zu einer unnötig komplizierten Entscheidungsstruktur führen.

Als Referenz werden zunächst \texttt{reg\_alpha = 0.05}, \texttt{reg\_lambda = 0.20} und \texttt{min\_split\_gain = 0.0} verwendet. Diese Werte wirken bereits leicht begrenzend, ohne mögliche technische Zusammenhänge früh zu unterdrücken. 
Die moderate und die starke Stufe erhöhen die Begrenzung anschließend gezielt. Dadurch lässt sich prüfen, ob ein Leistungsgewinn aus wiederkehrenden Zusammenhängen oder nur aus sehr speziellen Kombinationen einzelner Artikel entsteht. 
Beim Motor soll beispielsweise die gemeinsame Bedeutung von 400~V, 10~A und 4000~W gelernt werden können; eine nur einmal vorkommende Schreibweise soll dagegen keine eigene, stark verzweigte Regel erzeugen.

\section{Validierungsprozess}
\label{sec:validierungsprozess}

Jede Modellkonfiguration erzeugt für jeden Validierungsartikel eine Makro-Rangliste. Die Validierungsdaten wurden während des Trainings nicht zum Lernen der Entscheidungsbäume verwendet. 
Sie entsprechen damit einer Zwischenprüfung: Das Modell bearbeitet neue Aufgaben, während verschiedene Modellkonfigurationen noch miteinander verglichen werden dürfen.

Für eine Motoranfrage wird geprüft, auf welchem Rang sich das bekannte Zielmakro \texttt{M1\_F1\_V1} befindet und wie die fachlich verwandten Motor-Makros angeordnet sind. 
Hierzu werden Recall@1, Recall@3, Recall@5, Recall@10, MRR und hierarchisches nDCG@10 berechnet.

Recall@k beantwortet die Frage, ob das Zielmakro innerhalb der ersten \(k\) Positionen vorkommt. Bildlich wird nur ein Sichtfenster über die ersten Plätze der Rangliste gelegt. Bei Recall@3 werden die ersten drei Plätze betrachtet. 
Steht \texttt{M1\_F1\_V1} auf Rang 2, ist die Anfrage für Recall@3 erfolgreich. Ob das Makro auf Rang 1, 2 oder 3 steht, wird innerhalb dieser Metrik nicht weiter unterschieden.

Der Mean Reciprocal Rank (MRR) berücksichtigt dagegen die genaue Rangposition des ersten relevanten Makros. Hierzu wird der Kehrwert des Rangs verwendet. Steht \texttt{M1\_F1\_V1} auf Rang 1, beträgt der Wert 1. Steht es auf Rang 2, 
beträgt er \(1/2\). Auf Rang 5 beträgt er \(1/5\). Frühere Treffer erhalten dadurch einen höheren Wert.

Der Normalized Discounted Cumulative Gain (nDCG) bewertet die Reihenfolge mehrerer Kandidaten. Vereinfacht prüft nDCG, ob die wichtigen Makros weit oben und die weniger wichtigen Makros weiter unten stehen. 
Im Motorbeispiel ist eine Rangliste günstiger, wenn nach dem exakten Zielmakro zunächst Makros derselben Familie folgen und erst danach fachlich nicht verwandte Makros erscheinen. Der Zusatz \textit{normalized} bedeutet, 
dass die erzeugte Rangliste mit der bestmöglichen Reihenfolge verglichen wird. Ein Wert von 1 entspricht der idealen Reihenfolge.

Recall@k, MRR und nDCG betrachten somit unterschiedliche Seiten derselben Makro-Rangliste. Recall@k prüft, ob das Ziel gefunden wird. MRR prüft, wie früh es gefunden wird. nDCG prüft, ob die gesamte Reihenfolge fachlich sinnvoll aufgebaut ist.

Für die Auswahl wird zunächst der Recall-Teilscore gebildet:

\[
S_{Recall}=\frac{Recall@1+Recall@3+Recall@5+Recall@10}{4}.
\]

Die vier Recall-Werte werden zuerst zusammengefasst, damit die Recall-Gruppe insgesamt nicht stärker gewichtet wird als MRR oder nDCG. Der Validierungsscore ergibt sich anschließend zu

\[
S_{Validierung}=\frac{S_{Recall}+nDCG@10+MRR}{3}.
\]

Damit besitzen Recall, nDCG und MRR jeweils ein Gewicht von einem Drittel. Bildlich erhält jede Metrikgruppe eine gleich große Stimme. Eine Modellkonfiguration kann daher nicht allein deshalb ausgewählt werden, weil sie bei Recall@1 besonders gut ist, 
während die übrige Rangliste ungeeignet sortiert wird.

Die vier Recall-Werte werden innerhalb ihrer Gruppe gemittelt, weil sie dieselbe Frage mit unterschiedlich großem Sichtfenster beantworten. Würden Recall@1, Recall@3, Recall@5 und Recall@10 jeweils einzeln neben MRR und nDCG gewichtet, 
würde Recall vier Stimmen erhalten. Die zweistufige Berechnung stellt deshalb sicher, dass Finden, frühe Platzierung und fachliche Reihenfolge gleichberechtigt bleiben. 
Diese Auswahlregel wird sowohl innerhalb der LightGBM-Grid-Search als auch beim Vergleich der vollständigen Modellkonfigurationen verwendet. Dadurch erfolgt die Auswahl auf beiden Ebenen nach derselben fachlichen Zielsetzung.

Eine Modellkonfiguration wird nur bevorzugt, wenn sie die Mindestwerte \(Recall@10 \geq 0{,}50\) und \(nDCG@10 \geq 0{,}70\) erfüllt. Recall@10 von 0,50 bedeutet, 
dass sich das Zielmakro bei mindestens der Hälfte der Validierungsanfragen unter den ersten zehn Positionen befindet. nDCG@10 von 0,70 fordert zusätzlich eine ausreichende Qualität der Reihenfolge. 
Diese Grenzen verhindern die Auswahl einer Konfiguration mit einem guten Einzelwert, aber einer insgesamt unzureichenden Makro-Rangliste.

Die Schwellen wurden als Mindestanforderungen und nicht als Zielwerte festgelegt. Recall@10 von 0,50 verhindert, dass eine Konfiguration ausgewählt wird, 
die das korrekte Makro bei mehr als der Hälfte der Anfragen vollständig außerhalb des vorgesehenen Auswahlbereichs einordnet. nDCG@10 von 0,70 verhindert gleichzeitig eine Rangliste, in der zwar gelegentlich das Zielmakro enthalten ist, 
fachlich verwandte Kandidaten jedoch unsystematisch verteilt sind. Erst unter den Konfigurationen, die beide Anforderungen erfüllen, entscheidet der gleichgewichtete Validierungsscore.

Der Vergleich zwischen Trainings- und Validierungsergebnis dient zusätzlich der Erkennung von Overfitting. Erzielt eine Modellkonfiguration für die bekannten Motorartikel nahezu perfekte Werte, 
ordnet das Zielmakro bei den Validierungsartikeln jedoch deutlich schlechter ein, hat sie wahrscheinlich Besonderheiten der Trainingsdaten gelernt, die sich nicht auf neue Motorartikel übertragen lassen.

\section{Modellauswahl}
\label{sec:modellauswahl_technische_umsetzung}

Die Modellauswahl erfolgt ausschließlich anhand des Validierungsscores. Der Validierungsscore kann mit einer Gesamtnote verglichen werden, die aus drei gleichberechtigten Teilnoten besteht: Finden des Zielmakros, 
frühe Platzierung und sinnvolle Reihenfolge der gesamten Makro-Rangliste.

Angenommen eine Modellkonfiguration erreicht für die Motor- und übrigen Validierungsanfragen \(Recall@1=0{,}60\), \(Recall@3=0{,}70\), \(Recall@5=0{,}80\), \(Recall@10=0{,}90\), \(nDCG@10=0{,}80\) und \(MRR=0{,}70\). Zunächst ergibt sich

\[
S_{Recall}=\frac{0{,}60+0{,}70+0{,}80+0{,}90}{4}=0{,}75.
\]

Anschließend werden die drei Metrikgruppen gleichmäßig zusammengeführt:

\[
S_{Validierung}=\frac{0{,}75+0{,}80+0{,}70}{3}=0{,}75.
\]

Der Wert 0,75 beschreibt die gemeinsame Qualität der Konfiguration. Er ist keine zusätzliche Metrik für eine einzelne Rangliste, sondern eine festgelegte Auswahlregel für den Vergleich der Modellkonfigurationen.

Die Konfiguration mit dem höchsten gültigen Validierungsscore wird ausgewählt. Bei gleichen Scores dienen Recall@1, nDCG@10, MRR und Recall@10 ausschließlich als deterministische Tie-Breaker. 
Ein \textit{Tie-Breaker} ist eine feste Entscheidungsregel für einen Gleichstand. Er kommt nur zum Einsatz, wenn zwei Konfigurationen denselben Validierungsscore besitzen.

Die optionale Leave-one-out-Auswertung ist standardmäßig deaktiviert und beeinflusst weder Training noch Modellauswahl. 
Bei einer \textit{Leave-one-out-Auswertung} wird für die Diagnose ein Trainingsartikel vorübergehend aus seinem eigenen Makroprofil entfernt. Anschließend wird geprüft, ob das Modell das Zielmakro weiterhin findet. 
Bildlich wird beim Motorprofil kurzzeitig die Karte des gerade geprüften Motors aus dem gemeinsamen Stapel genommen. Dadurch kann untersucht werden, wie stark der Artikel von seinen eigenen Angaben im Makroprofil profitiert. 
Für das reguläre Training wird jedoch das vollständige Profil von \texttt{M1\_F1\_V1} verwendet.

Die Standardkonfiguration lautet daher \texttt{KI\_ENABLE\_TRAIN\_LEAVE\_ONE\_OUT\_EVALUATION = 0}. Nur der gesonderte Diagnoselauf aktiviert den Wert 1. Diese Trennung wurde gewählt, 
weil alle Makros im Trainingsanteil vorkommen müssen und Makros mit nur einem Artikel oder einer unteilbaren Zusammenhangskomponente dem Training zugeordnet werden. Würde Leave-one-out als reguläres Trainingsverfahren eingesetzt, 
könnte bei einem solchen Motor-Makro kein vollständiges Zielprofil mehr verbleiben. Als Diagnose zeigt die Auswertung dagegen den Eigenbeitrag eines Artikels, ohne das trainierte Modell oder die Modellauswahl zu verändern.

\section{Testprozess}
\label{sec:testprozess_technische_umsetzung}

Nach der Modellauswahl wird genau diese Konfiguration unverändert auf den Testdaten bewertet. Die Testdaten entsprechen der Abschlussprüfung. Sie wurden weder zum Training noch zur Auswahl der Hyperparameter verwendet. 
Dadurch zeigen sie, wie gut sich die erlernten Zusammenhänge auf bisher unbekannte Artikel übertragen lassen.

Für einen unbekannten Motorartikel wird erneut eine Makro-Rangliste erzeugt. Befindet sich \texttt{M1\_F1\_V1} auf Rang 2, ergibt sich folgendes Bild: Recall@1 ist nicht erfolgreich, weil nur der erste Platz betrachtet wird. 
Recall@3, Recall@5 und Recall@10 sind erfolgreich, weil Rang 2 innerhalb dieser Bereiche liegt. Der reziproke Rang beträgt \(1/2\). nDCG@10 bewertet zusätzlich, ob die übrigen Motor-Makros sinnvoll angeordnet sind.

Die Testmetriken werden nicht zur Änderung der Label-Policy, der Hyperparameter oder der verwendeten technischen Eigenschaften eingesetzt. Andernfalls würde die Abschlussprüfung nachträglich zur Übungsaufgabe werden. 
Die Testwerte dienen deshalb ausschließlich der abschließenden Bewertung des bereits anhand der Validierungsdaten ausgewählten Modells.

Diese Festlegung ist besonders für die Konfigurationsserie relevant. Binäre und hierarchische Labels sowie Referenz-, moderate und starke Regularisierung werden ausschließlich anhand der Validierungsdaten verglichen. 
Der optionale Leave-one-out-Lauf nimmt nicht an der Modellauswahl teil. Erst nachdem eine Konfiguration feststeht, werden ihre Testwerte betrachtet. 
Dadurch bleibt der Testanteil ein unabhängiger Nachweis für die Übertragbarkeit auf unbekannte Motorartikel.

\section{Anwendungsprozess}
\label{sec:anwendungsprozess_technische_umsetzung}

Im Anwendungsprozess bildet ein neues Motordatenblatt den Ausgangspunkt. Auf diesem Datenblatt stehen beispielsweise Hersteller SEW, 400~V, 10~A, 4000~W, 50~Hz und der Hinweis \textit{Mit Kühlung}. 
Diese Angaben werden mit derselben Normalisierungs- und Extraktionslogik verarbeitet, die bereits für die Trainings-, Validierungs- und Testdaten verwendet wurde.

Die Normalisierung sorgt dafür, dass unterschiedliche Schreibweisen derselben technischen Eigenschaft gemeinsam verglichen werden können. Steht auf einem Datenblatt beispielsweise 4~kW und in den Trainingsdaten 4000~W, 
werden beide Angaben auf eine gemeinsame Darstellung überführt. Dadurch erkennt das System, dass kein technischer Unterschied vorliegt.

Aus den technischen Eigenschaften wird ein Anfrageprofil erzeugt. Dieses Anfrageprofil kann wie ein ausgefüllter Motorsteckbrief betrachtet werden, der an alle vorhandenen Makroprofile gehalten wird. 
Für jedes Anfrage-Kandidat-Paar werden dieselben Vergleichsmerkmale wie während des Trainings berechnet. Das LambdaMART-Modell erzeugt anschließend für jeden Kandidaten einen Relevanzscore.

Angenommen \texttt{M1\_F1\_V1} stimmt bei Spannung, Strom, Leistung, Frequenz und dem Hinweis \textit{Mit Kühlung} weitgehend mit der Anfrage überein. \texttt{M1\_F1\_V2} besitzt ebenfalls passende elektrische Werte, 
beschreibt jedoch eine andere Ausführung. \texttt{M1\_F2\_V1} gehört zwar zur Klasse der Motor-Makros, weist jedoch eine weniger passende Familie auf. Das Modell kann daraus beispielsweise folgende Makro-Rangliste erzeugen:

\begin{enumerate}
    \item \texttt{M1\_F1\_V1},
    \item \texttt{M1\_F1\_V2},
    \item \texttt{M1\_F2\_V1}.
\end{enumerate}

Die Kandidaten werden ausschließlich nach ihrem unveränderten LambdaMART-Relevanzscore absteigend sortiert. Es wird keine nachträgliche Motorregel angewendet. 
Die Platzierung von \texttt{M1\_F1\_V1} entsteht somit aus den während des Trainings gelernten technischen Zusammenhängen.

Die gleiche Normalisierung, Merkmalsreihenfolge und Featuredefinition wird in Training und Anwendung verwendet, weil bereits eine abweichende Einheitenumrechnung oder eine anders angeordnete Zahlenreihe die Modellbewertung verändern würde. 
Die für das Modell gespeicherten Feature-Namen und Normalisierungsparameter werden deshalb gemeinsam mit dem Modell bereitgestellt. So wird beispielsweise 4~kW auch im Anwendungsprozess in derselben Form wie 4000~W verarbeitet.

\section{Kapitelergebnis}
\label{sec:kapitelergebnis_technische_umsetzung}

Die technische Umsetzung erzeugt aus den aufbereiteten Artikel- und Makrodaten vollständige Makroprofile, datenblattkompatible Anfrageprofile sowie Anfrage-Kandidat-Paare mit numerischen Merkmalsvektoren. 
Am Motorbeispiel lässt sich dieser Ablauf wie eine schrittweise Prüfung darstellen: Zunächst werden die technischen Eigenschaften des Motors geordnet. Anschließend wird der Motorsteckbrief mit jedem Makroprofil verglichen. 
Das LambdaMART-Modell bewertet die entstehenden Vergleiche und sortiert die Makros nach ihrer geschätzten Relevanz.

Das Modell wird ohne manuelle Beispielgewichtung und ohne regelbasierte Nachsortierung trainiert. Die Hyperparameter und die Label-Policy werden anhand eines gleichgewichteten Validierungsscores aus Recall@k, nDCG@10 und MRR ausgewählt. 
Die Validierungsdaten dienen der Auswahl der Modellkonfiguration. Die Testdaten dienen ausschließlich der abschließenden Bewertung.

Das Ergebnis dieses Kapitels bildet das folgende Artefakt:

\begin{quote}
\textbf{Artefakt A4: Trainiertes Empfehlungssystem}
\end{quote}

\begin{itemize}
    \item erzeugte Makroprofile mit allen beobachteten Trainingswerten,
    \item erzeugte Anfrageprofile mit datenblattkompatiblen technischen Eigenschaften,
    \item Anfrage-Kandidat-Paare mit den zugehörigen Vergleichsmerkmalen,
    \item Merkmalsvektoren für Training, Validierung, Test und Anwendung,
    \item trainiertes LambdaMART-Modell,
    \item anhand der Validierungsdaten ausgewählte Modellkonfiguration,
    \item validiertes und getestetes Empfehlungssystem,
    \item Anwendungsprozess zur Erzeugung einer nach Relevanz sortierten Makro-Rangliste.
\end{itemize}

Artefakt A4 bildet die Grundlage für die anschließende Evaluation der Empfehlungsqualität. Im folgenden Kapitel werden die erzeugten Makro-Ranglisten anhand der in Abschnitt~\ref{sec:validierungsansätze} eingeführten Metriken Recall@k, MRR und nDCG bewertet.


% \chapter{Technische Umsetzung}
% \label{chap:technische_umsetzung}


% \section{Implementierung der Datenobjekte} 

% Die fachlichen Grundlagen der Makroprofile, Anfrageprofile und Anfrage-Kandidat-Paare wurden bereits in Kapitel \ref{chap:konzeption_des_empfehlungssystems} beschrieben. 
% Ziel dieses Abschnitts ist daher die technische Umsetzung dieser Datenobjekte auf Basis der Trainings- und Evaluationsbasis aus Artefakt A1. Hierzu werden zunächst die aufbereiteten Artikel-, 
% Makro- und Zuordnungsdaten eingelesen und entsprechend ihrer Datensplit-Zuordnung den Trainings-, Validierungs- und Testdaten zugeordnet. 
% Anschließend werden die Makroprofile aus den Artikeln der Trainingsdaten erzeugt und als Suchraum der Makrokandidaten bereitgestellt. 
% Parallel dazu erfolgt die Bildung der Anfrageprofile aus den Artikeln der jeweiligen Datenbestände. Makroidentifizierende Informationen werden dabei entfernt, 
% sodass ausschließlich die für eine spätere Anfrage verfügbaren technischen Eigenschaften erhalten bleiben. Im letzten Schritt werden die Anfrageprofile mit den Makroprofilen kombiniert und daraus Anfrage-Kandidat-Paare erzeugt. 
% Für jedes Anfrage-Kandidat-Paar werden die in Abschnitt \ref{sec:anfrage_kandidat_paar_vergleichsmerkmale_und_merkmalsvektor} beschriebenen Vergleichsmerkmale berechnet und zu Merkmalsvektoren zusammengeführt. 
% Das Ergebnis dieses Verarbeitungsschrittes sind die für Training, Validierung und Test benötigten Anfrage-Kandidat-Paare einschließlich der zugehörigen Merkmalsvektoren und Relevanzinformationen. 
% Diese bilden die Eingabe des nachfolgenden Trainingsprozesses des LambdaMART-Modells.


% \section{Training des LambdaMART-Modells} 

% Nach der Erzeugung der Anfrage-Kandidat-Paare und der zugehörigen Merkmalsvektoren liegen für jede Anfrage die Eingabedaten für das Learning-to-Rank-Verfahren vor. Ziel des Trainingsprozesses besteht darin, 
% die Zusammenhänge zwischen den Vergleichsmerkmalen eines Anfrage-Kandidat-Paares und dessen Relevanz für die Anfrage zu erlernen. Hierzu wird das in Kapitel \ref{chap:auswahl_des_modellansatzes} ausgewählte LambdaMART-Verfahren verwendet. 
% Die fachlichen Grundlagen von Learning-to-Rank, LambdaRank und LambdaMART wurden bereits in den Abschnitten \ref{sec:learning_to_rank}, \ref{sec:lambdarank} und \ref{sec:lambdamart} erläutert. 
% Die Eingabe des Trainingsprozesses besteht aus den Merkmalsvektoren der Anfrage-Kandidat-Paare. Jeder Merkmalsvektor beschreibt dabei die technische Übereinstimmung zwischen einem Anfrageprofil und einem Makroprofil. 
% Zusätzlich besitzt jedes Anfrage-Kandidat-Paar ein Relevanzlabel, das aus der bekannten Artikel-Makro-Zuordnung der Trainingsdaten abgeleitet wird. Während des Trainings werden die Anfrage-Kandidat-Paare anfrageweise zu Rankinggruppen zusammengefasst. 
% Innerhalb einer Rankinggruppe lernt das Modell, relevante Makros höher zu bewerten als nicht relevante Makros. Hierzu erzeugt LambdaMART für jedes Anfrage-Kandidat-Paar einen numerischen Relevanzscore. 
% Die Berechnung dieser Relevanzscores erfolgt durch mehrere aufeinanderfolgende Entscheidungsbäume. Jeder neu erzeugte Baum versucht dabei, die Fehler der bereits trainierten Bäume zu verringern. 
% Die Ergebnisse aller Bäume werden anschließend zu einem gemeinsamen Relevanzscore zusammengeführt. Dadurch entsteht ein Modell, das technische Muster innerhalb der Vergleichsmerkmale erkennen und zur Bewertung neuer Anfragen verwenden kann. 
% Das Ergebnis des Trainingsprozesses ist ein trainiertes LambdaMART-Modell, welches für jedes Anfrage-Kandidat-Paar einen Relevanzscore berechnen kann. Die endgültige Auswahl der Modellkonfiguration erfolgt anschließend anhand des Validierungsprozesses.

% \section{Hyperparameteroptimierung}

% Die Leistungsfähigkeit eines LambdaMART-Modells wird wesentlich durch dessen Hyperparameter beeinflusst. Hyperparameter beschreiben dabei Eigenschaften des Lernprozesses und werden vor dem eigentlichen Training festgelegt. 
% Im Gegensatz zu den Modellparametern werden sie nicht aus den Trainingsdaten gelernt, sondern müssen geeignet gewählt werden. Um eine geeignete Modellkonfiguration zu bestimmen, wird eine Grid Search durchgeführt. 
% Hierbei werden mehrere Kombinationen von Hyperparametern systematisch erzeugt und nacheinander trainiert. Jede Modellkonfiguration durchläuft anschließend den in Abschnitt 7.4 beschriebenen Validierungsprozess. 
% Die untersuchten Modellkonfigurationen unterscheiden sich unter anderem in der Anzahl der Entscheidungsbäume, der Lernrate sowie der Komplexität der einzelnen Entscheidungsbäume. 
% Dadurch können sowohl einfachere als auch komplexere Modelle untersucht werden. Für jede Modellkonfiguration wird ein separates LambdaMART-Modell trainiert. Nach Abschluss des Trainings erfolgt die Bewertung anhand der Validierungsdaten. 
% Die erzielten Ergebnisse werden gespeichert und mit den Ergebnissen der übrigen Modellkonfigurationen verglichen. Die Grid Search verfolgt das Ziel, eine Modellkonfiguration zu finden, 
% die die Ranglistenqualität maximiert und gleichzeitig eine ausreichende Generalisierungsfähigkeit auf unbekannte Daten besitzt. 
% Dadurch soll verhindert werden, dass das Modell die Trainingsdaten zu stark auswendig lernt und auf neuen Anfragen schlechtere Ergebnisse liefert. 
% Das Ergebnis der Hyperparameteroptimierung besteht aus einer ausgewählten Modellkonfiguration, welche als Grundlage für die anschließende Validierung und den späteren Testprozess dient.


% \section{Validierungsprozess} 

% Ziel des Validierungsprozesses ist die Bewertung der während der Hyperparameteroptimierung erzeugten Modellkonfigurationen. Während die Trainingsdaten zum Erlernen der Zusammenhänge zwischen den Vergleichsmerkmalen und den Relevanzwerten verwendet werden, 
% dienen die Validierungsdaten zur unabhängigen Beurteilung der Leistungsfähigkeit des Modells auf bisher nicht verwendeten Anfrageprofilen (vergleiche Abschnitt \ref{sec:splitregeln}). 
% Hierzu wird jede während der Grid Search erzeugte Modellkonfiguration nach Abschluss des Trainings mit den Validierungsanfragen bewertet. 
% Für jede Anfrage wird eine Rangliste aus Makrokandidaten erzeugt und anschließend mit den bekannten Artikel-Makro-Zuordnungen der Validierungsdaten verglichen. 
% Die Bewertung erfolgt anhand der in Abschnitt \ref{sec:validierungsansätze} eingeführten Metriken Recall@k, Mean Reciprocal Rank (MRR) und Normalized Discounted Cumulative Gain (nDCG). Recall@k bewertet dabei die Wahrscheinlichkeit, 
% innerhalb der ersten \(k\) Positionen mindestens ein relevantes Makro zu finden. Der MRR berücksichtigt zusätzlich die Position des ersten relevanten Makros innerhalb der Rangliste. 
% nDCG bewertet die Qualität der gesamten Rangliste und berücksichtigt dabei die Position der relevanten Makros. Für jede Modellkonfiguration werden die erzielten Kennzahlen gespeichert und mit den Ergebnissen der übrigen Modellkonfigurationen verglichen. 
% Dadurch entsteht für jede Konfiguration ein Qualitätsprofil, welches die Empfehlungsqualität des Modells auf den Validierungsdaten beschreibt. Der Validierungsprozess dient gleichzeitig der Erkennung von Overfitting. 
% Erzielt ein Modell sehr gute Ergebnisse auf den Trainingsdaten, jedoch deutlich schlechtere Ergebnisse auf den Validierungsdaten, 
% deutet dies darauf hin, dass das Modell Trainingsdaten auswendig gelernt hat und die erlernten Zusammenhänge nur eingeschränkt auf neue Anfragen übertragen kann. Das Ergebnis des Validierungsprozesses ist die Auswahl derjenigen Modellkonfiguration, 
% die die höchste Empfehlungsqualität auf den Validierungsdaten erzielt. Diese Modellkonfiguration wird anschließend als Grundlage für den Testprozess verwendet.

% \section{Modellauswahl} 

% Nach Abschluss der Hyperparameteroptimierung und des Validierungsprozesses liegen mehrere bewertete Modellkonfigurationen vor. Ziel der Modellauswahl besteht darin, aus diesen Konfigurationen das Modell mit der höchsten Empfehlungsqualität auszuwählen. 
% Hierzu werden die im Validierungsprozess ermittelten Kennzahlen Recall@k, MRR und nDCG gemeinsam betrachtet. Die Auswahl erfolgt dabei nicht ausschließlich anhand einer einzelnen Kennzahl, sondern anhand der Gesamtqualität der erzeugten Ranglisten. 
% Besondere Bedeutung kommt den Kennzahlen Recall@k und nDCG zu. Während Recall@k bewertet, ob das relevante Makro innerhalb der vorderen Rangpositionen gefunden wird, 
% bewertet nDCG zusätzlich die Qualität der gesamten Rangfolge der Makrokandidaten (vergleiche Abschnitt \ref{sec:validierungsansätze}). 
% Die Modellkonfiguration mit den besten Ergebnissen auf den Validierungsdaten wird als finale Modellkonfiguration ausgewählt. Anschließend erfolgt ein erneutes Training dieser Konfiguration auf Grundlage der vorgesehenen Trainingsdaten. 
% Dadurch wird sichergestellt, dass für den nachfolgenden Testprozess die leistungsfähigste Modellvariante verwendet wird. 
% Das Ergebnis der Modellauswahl ist ein trainiertes LambdaMART-Modell mit festgelegten Hyperparametern und einer validierten Rankingfunktion. 
% Dieses Modell bildet die Grundlage für den anschließenden Testprozess und die spätere Anwendung des Empfehlungssystems.


% \section{Testprozess} 

% Nach der Auswahl der finalen Modellkonfiguration erfolgt die abschließende Bewertung des trainierten LambdaMART-Modells auf den Testdaten. 
% Die Testdaten wurden bereits während der Erstellung des Datensplits von den Trainings- und Validierungsdaten getrennt und wurden bisher weder für das Training noch für die Modellauswahl verwendet (vergleiche Abschnitt \ref{sec:splitregeln}). 
% Ziel des Testprozesses ist die Untersuchung der Generalisierungsfähigkeit des Modells. Hierbei wird überprüft, inwieweit die während des Trainings erlernten Zusammenhänge auf bisher unbekannte Anfrageprofile übertragen werden können. 
% Für jede Testanfrage wird analog zum Validierungsprozess eine Makro-Rangliste erzeugt. Die Bewertung dieser Ranglisten erfolgt anhand der in Abschnitt \ref{sec:validierungsansätze} beschriebenen Kennzahlen Recall@k, MRR und nDCG. 
% Im Gegensatz zum Validierungsprozess werden während des Testprozesses keine Hyperparameter mehr angepasst und keine Modellentscheidungen mehr getroffen. Das bereits ausgewählte LambdaMART-Modell wird unverändert auf den Testdaten angewendet. 
% Dadurch wird sichergestellt, dass die erzielten Ergebnisse eine unverzerrte Einschätzung der zu erwartenden Leistungsfähigkeit des Empfehlungssystems darstellen. 
% Die Ergebnisse des Testprozesses bilden die Grundlage für die spätere Bewertung der Empfehlungsqualität. 
% Gleichzeitig liefern sie die Kennzahlen für den Vergleich mit alternativen Konfigurationen und dienen als Nachweis der Leistungsfähigkeit des entwickelten Empfehlungssystems. 
% Das Ergebnis des Testprozesses ist somit die finale Bewertung des trainierten LambdaMART-Modells auf bisher unbekannten Anfragen und bildet die Grundlage für die anschließende Anwendung des Empfehlungssystems.


% \section{Anwendungsprozess} 

% Nach Abschluss des Testprozesses steht ein trainiertes und validiertes LambdaMART-Modell für die Verarbeitung neuer Anfragen zur Verfügung. 
% Der Anwendungsprozess entspricht dabei dem in Kapitel \ref{chap:konzeption_des_empfehlungssystems} beschriebenen Datenfluss, wird jedoch auf bisher unbekannte Artikel angewendet. 
% Ausgangspunkt des Anwendungsprozesses ist ein neuer Artikel beziehungsweise ein neues Datenblatt. 
% Aus den verfügbaren technischen Eigenschaften werden zunächst die für das Empfehlungssystem erforderlichen Informationen extrahiert und in ein Anfrageprofil überführt. 
% Analog zum Trainingsprozess werden dabei ausschließlich solche Informationen verwendet, die zum Zeitpunkt der Anfrage bekannt sind. Anschließend wird das Anfrageprofil mit allen Makroprofilen des \texttt{Suchraum\_Makroprofile} kombiniert. 
% Für jedes Anfrage-Kandidat-Paar werden die Vergleichsmerkmale berechnet und zu einem Merkmalsvektor zusammengeführt. Diese Merkmalsvektoren bilden die Eingabe des trainierten LambdaMART-Modells. 
% Das Modell berechnet anschließend für jeden Makrokandidaten einen Relevanzscore. Je höher dieser Relevanzscore ist, desto besser passt das jeweilige Makro zum Anfrageprofil. 
% Die Makrokandidaten werden daraufhin entsprechend ihrer Relevanzscores in absteigender Reihenfolge sortiert. Das Ergebnis des Anwendungsprozesses ist eine Makro-Rangliste, 
% welche die bewerteten Makrokandidaten entsprechend ihrer geschätzten Relevanz für die Anfrage enthält. Die vorderen Positionen dieser Rangliste stellen die vom Empfehlungssystem empfohlenen Makros dar. Durch die Trennung von Anfrageprofil, 
% Makroprofil und Vergleichsmerkmalen bleibt die Architektur des Anwendungspfads identisch zur Architektur des Trainings- und Testprozesses. 
% Dadurch können die während des Trainings erlernten Zusammenhänge unmittelbar auf neue und bisher unbekannte Anfragen übertragen werden.

% \section{Kapitelergebnis} 
% Ziel dieses Kapitels war die technische Umsetzung der in Kapitel \ref{chap:konzeption_des_empfehlungssystems} entwickelten Architektur des Empfehlungssystems. 
% Hierzu wurden die einzelnen Verarbeitungsschritte von der Erzeugung der Datenobjekte bis zur Bereitstellung eines trainierten LambdaMART-Modells beschrieben. Zunächst wurden Makroprofile, Anfrageprofile sowie die Anfrage-Kandidat-Paare technisch erzeugt. 
% Anschließend wurden für jedes Anfrage-Kandidat-Paar die Vergleichsmerkmale berechnet und zu Merkmalsvektoren zusammengeführt, welche die Eingabe des LambdaMART-Modells bilden. Darauf aufbauend erfolgte das Training des Modells. 
% Durch die Hyperparameteroptimierung und den Validierungsprozess wurde eine geeignete Modellkonfiguration bestimmt. 
% Die ausgewählte Konfiguration wurde anschließend im Testprozess auf bisher unbekannten Anfragen bewertet und für die spätere Anwendung bereitgestellt. Der abschließende Anwendungsprozess zeigt, 
% wie neue Anfrageprofile durch das trainierte Modell bewertet und in eine nach Relevanz sortierte Makro-Rangliste überführt werden können. 

% Das Ergebnis dieses Kapitels bildet das folgende Artefakt: 

% \begin{quote} \textbf{Artefakt A4: Trainiertes Empfehlungssystem} 
% \end{quote} 
% \begin{itemize} 
%     \item erzeugte Makroprofile, 
%     \item erzeugte Anfrageprofile, 
%     \item Anfrage-Kandidat-Paare mit den zugehörigen Vergleichsmerkmalen, 
%     \item Merkmalsvektoren für Training, Validierung, Test und Anwendung, 
%     \item trainiertes LambdaMART-Modell, 
%     \item ausgewählte Modellkonfiguration, 
%     \item validiertes und getestetes Empfehlungssystem, 
%     \item Anwendungsprozess zur Erzeugung einer nach Relevanz sortierten Makro-Rangliste. 
% \end{itemize} 

% Artefakt A4 bildet die Grundlage für die anschließende Evaluation der Empfehlungsqualität. 
% Im folgenden Kapitel werden die erzeugten Makro-Ranglisten anhand der in Abschnitt \ref{sec:validierungsansätze} eingeführten Metriken Recall@k, MRR und nDCG bewertet.




% \section{Implementierungsgrundlage}


% \begin{itemize} \item Umsetzung der in Kapitel \ref{chap:konzeption_des_empfehlungssystems} entwickelten Architektur. \item Verwendung der Trainings- und Evaluationsbasis aus Artefakt A1. \item Verwendung des ausgewählten LambdaMART-Modells aus Artefakt A2. \item Ziel ist die Erstellung eines trainierten Empfehlungssystems als Artefakt A4. \item Umsetzung als mehrstufige Pipeline entsprechend der Komponentenarchitektur aus Kapitel \ref{chap:konzeption_des_empfehlungssystems}. \end{itemize}


% \section{Erzeugung der Makroprofile}

% \begin{itemize} \item Einlesen der Trainingsartikel. \item Gruppierung aller Artikel anhand der Artikel-Makro-Zuordnung. \item Bildung eines Makroprofils je Zielmakro. \item Übernahme der technischen Makroeigenschaften. \item Zusammenführung der technischen Artikeleigenschaften aller zugeordneten Trainingsartikel. \item Speicherung der Merkmalswerte als Mengen. \item Speicherung der Makroprofile im Suchraum\_Makroprofile. \item Ermittlung zusätzlicher Informationen: \begin{itemize} \item Makrofamilie, \item Varianten, \item Support, \item Häufigkeiten, \item Klassifikationen. \end{itemize} \end{itemize}

% \section{Erzeugung der Anfrageprofile}

% \begin{itemize} \item Erstellung eines Anfrageprofils für jeden Artikel. \item Übernahme der technischen Artikeleigenschaften. \item Übernahme der nicht identitätsstiftenden technischen Makroeigenschaften. \item Entfernung makroidentifizierender Informationen. \item Ausschluss von: \begin{itemize} \item P\_ARTICLE\_MACRO, \item Makrobezeichnungen, \item Zielmakrokennungen, \item Dokumentpfaden, \item Bearbeiterinformationen. \end{itemize} \item Speicherung der Anfrageprofile in: \begin{itemize} \item Anfrageprofile\_Training, \item Anfrageprofile\_Validierung, \item Anfrageprofile\_Test. \end{itemize} \end{itemize}


% \section{Erzeugung der Anfrage-Kandidat-Paare}
% \begin{itemize} \item Kombination jedes Anfrageprofils mit allen Makroprofilen des Suchraums. \item Bildung der vollständigen Kandidatenmenge. \item Erzeugung einer Rankinggruppe pro Anfrage. \item Training: \begin{itemize} \item positives Zielmakro, \item negative Makrokandidaten. \end{itemize} \item Anwendung: \begin{itemize} \item unbekannter Anfrageartikel, \item Bewertung aller Makrokandidaten. \end{itemize} \item Ergebnis: \begin{itemize} \item Anfrage-Kandidat-Paar, \item Rankinggruppe. \end{itemize} \end{itemize}

% \section{Berechnung der Vergleichsmerkmale}
% \begin{itemize} \item Technische Vergleichsmerkmale bilden die Eingabe von LambdaMART. \item Vergleich zwischen Anfrageprofil und Makrokandidatenprofil. \item Berechnung strukturierter Merkmale. \item Beispiele: \begin{itemize} \item Herstellerübereinstimmung, \item Produkttypübereinstimmung, \item Spannungsklassenübereinstimmung, \item Stromklassenübereinstimmung, \item Leistungsklassenübereinstimmung, \item Frequenzübereinstimmung, \item Klassifikationsübereinstimmung. \end{itemize} \item Berechnung numerischer Abweichungen. \item Berechnung kategorialer Übereinstimmungen. \item Berechnung lexikalischer Übereinstimmungen. \item Bildung des Merkmalsvektors. \end{itemize}

% \section{Training des LambdaMART-Modells}
% \begin{itemize} \item Verwendung der Anfrage-Kandidat-Paare der Trainingsdaten. \item Verwendung des LightGBM Rankers. \item Ranking-Objective: \begin{itemize} \item \texttt{lambdarank} \end{itemize} \item Bildung von Rankinggruppen. \item Nutzung der Relevanzlabels. \item Gradient Boosting Decision Trees. \item Erzeugung von Relevanzscores. \item Optimierung der Kandidatenreihenfolge. \item Speicherung der trainierten Modelle. \end{itemize}


% \section{Hyperparameteroptimierung}
% \begin{itemize} \item Durchführung einer Grid Search. \item Training mehrerer Modellkonfigurationen. \item Variation der Hyperparameter: \begin{itemize} \item Anzahl Bäume, \item Lernrate, \item Anzahl Blätter, \item Min Child Samples, \item Truncation Level. \end{itemize} \item Bewertung auf den Validierungsdaten. \item Vergleich aller Modellkonfigurationen. \item Auswahl der Konfiguration mit der höchsten Validierungsleistung. \end{itemize}

% \section{Validierungsprozess}
% \begin{itemize} \item Bewertung jeder Modellkonfiguration. \item Verwendung der Validierungsanfragen. \item Berechnung von: \begin{itemize} \item Recall@1, \item Recall@3, \item Recall@5, \item Recall@10. \end{itemize} \item Berechnung von: \begin{itemize} \item MRR, \item nDCG@1, \item nDCG@3, \item nDCG@5, \item nDCG@10. \end{itemize} \item Erkennung von Overfitting. \item Auswahl des besten Modells. \end{itemize}

% \section{Testprozess}
% \begin{itemize} \item Verwendung des ausgewählten LambdaMART-Modells. \item Verwendung der Testanfragen. \item Keine weitere Parameteranpassung. \item Berechnung der finalen Kennzahlen. \item Bewertung der Generalisierungsfähigkeit. \item Erstellung des Testergebnisses. \end{itemize}

% \section{Anwendungsprozess}
% \begin{itemize} \item Eingang eines neuen Artikels. \item Erstellung eines Anfrageprofils. \item Bildung der Anfrage-Kandidat-Paare. \item Berechnung der Vergleichsmerkmale. \item Bewertung aller Makrokandidaten durch das trainierte LambdaMART-Modell. \item Berechnung der Relevanzscores. \item Sortierung nach absteigender Relevanz. \item Erzeugung der Makro-Rangliste. \item Ausgabe der Top-k Empfehlungen. \end{itemize}

% \section{Kapitelergebnis}
% \begin{itemize} \item Erzeugung der Makroprofile. \item Erzeugung der Anfrageprofile. \item Erzeugung der Anfrage-Kandidat-Paare. \item Berechnung der Vergleichsmerkmale. \item Training des LambdaMART-Modells. \item Hyperparameteroptimierung. \item Validierung. \item Test. \item Bereitstellung des trainierten Empfehlungssystems. \end{itemize} \begin{quote} \textbf{Artefakt A4: Trainiertes Empfehlungssystem} \end{quot

    
% \textbf{DIE NACHFOLGENDEN ABSCHNITTE DIENEN NUR DER INFORMATION FÜR DIE AUSFORMULIERUNG DES KAPITELS DAR !!!!!!!!!!!!!!!!!!!!!!!!}


% \section{Überführung der Artefakte in die Implementierung} 
% \label{sec:ueberfuehrung_der_artefakte} 

% Kapitel~7 beschreibt die technische Zusammenführung der in den vorherigen Kapiteln entwickelten Artefakte. 
% Die aufbereitete Trainings- und Evaluationsbasis A1, der ausgewählte Modellansatz A2 und die Systemarchitektur A3 bilden die Eingaben der Implementierung. 
% Aus ihnen wird das trainierte Empfehlungssystem A4 erzeugt. 

% \subsection{Technische Schwerpunkte} 
% \label{subsec:technische_schwerpunkte_artefaktueberfuehrung} 

% Die Implementierung umfasst folgende Schwerpunkte: 
% \begin{itemize} 
%     \item Übernahme der aufbereiteten Daten und Datensplits aus A1, 
%     \item Umsetzung des in A2 ausgewählten LambdaMART-Verfahrens, 
%     \item Abbildung der in A3 definierten Komponenten auf ausführbare Python-Module, 
%     \item Bildung der Kandidatenprofile und Anfrage-Kandidat-Matrizen, 
%     \item Training und Auswahl des Rankingmodells, 
%     \item Speicherung der Modell-, Profil-, Feature- und Laufzeitinformationen, 
%     \item Erzeugung des Artefakts A4. 
% \end{itemize} 

% \subsection{Technische Umsetzung} 

% \label{subsec:technische_umsetzung_artefaktueberfuehrung} 

% Die Implementierung ist als sequenzielle Python-Pipeline aufgebaut. Sie gliedert sich in die Bereiche \texttt{stage\_910\_data\_preparation}, \texttt{stage\_920\_data\_splitting}, \texttt{stage\_930\_model\_training} und \texttt{stage\_940\_evaluation}. 
% Die Stufe 910 erzeugt die standardisierten Datenrepräsentationen. Stufe 920 erstellt die gruppierten Trainings-, Validierungs- und Testanteile. Stufe 930 bildet die Kandidatenprofile und Anfrage-Kandidat-Matrizen und führt das LambdaMART-Training durch. 
% Stufe 940 erzeugt Vorhersagen, Rankingmetriken und technische Prüfberichte. Die vollständige Ausführung wird durch \texttt{step\_936\_run\_full\_training\_pipeline.py} gesteuert. 
% Die zusammengefasste technische Evaluation erfolgt über \texttt{step\_949\_evaluate\_all\_suite.py}. Das Ergebnis der Pipeline ist nicht nur die serialisierte Modelldatei. 

% Für eine konsistente Modellanwendung werden zusätzlich folgende Bestandteile benötigt: 

% \begin{itemize} 
%     \item das trainierte LambdaMART-Modell, 
%     \item die aus dem Trainingsanteil erzeugten Kandidatenprofile, 
%     \item die verwendeten Featurebezeichnungen, 
%     \item die verbindliche Feature-Reihenfolge, 
%     \item die aktive Modell- und Laufzeitkonfiguration, 
%     \item die Metadaten des Trainingslaufs. 
% \end{itemize} 

% Diese Bestandteile bilden gemeinsam das Artefakt A4. 

% \subsection{Abgleich zwischen bisheriger Arbeit und Code} 

% \label{subsec:abgleich_artefaktueberfuehrung} 
% \subsubsection*{Übereinstimmungen} 

% Die im bisherigen Text vorgesehene Abfolge aus Datenaufbereitung, Datenaufteilung, Modelltraining und Evaluation wird durch die Pipeline-Stufen 910 bis 940 umgesetzt. 
% Auch die Auswahl von LambdaMART und die Trennung von Trainings-, Validierungs- und Testdaten entsprechen der bisherigen Konzeption. 

% \subsubsection*{Technische Konkretisierungen} 

% Die bisherige Arbeit beschreibt die Artefakte und Datenflüsse auf fachlicher Ebene. Der Code konkretisiert diese durch Python-Module, Konfigurationsdateien, laufbezogene Ergebnisverzeichnisse und serialisierte Zwischen- und Ergebnisdateien. 
% Artefakt A4 umfasst technisch mehr als das trainierte Modell. Für die Anwendung sind auch die Kandidatenprofile, die Featuredefinition und die aktive Laufzeitkonfiguration erforderlich. 

% \subsubsection*{Abweichungen} 

% A1, A2 und A3 werden nicht in Kapitel~7 erzeugt. Sie bilden die Eingaben der Implementierung. Erst deren technische Zusammenführung erzeugt A4: 
% \begin{equation} A1 + A2 + A3 \xrightarrow{\text{Implementierung und Training}} A4. 
%     \label{eq:artefaktueberfuehrung_a4} 
% \end{equation} 

% Der Code berechnet bereits während der Pipelineausführung Evaluationsmetriken und Prüfberichte. Dies stellt keinen methodischen Widerspruch zur Kapitelstruktur dar. Kapitel~7 beschreibt deren technische Erzeugung, 
% während Kapitel~8 die Ergebnisse bewertet und interpretiert. \subsubsection*{Konsequenz für die Dokumentation} A4 wird als vollständiges Modellartefakt dokumentiert. 
% Es besteht aus zwei Teilen: 
% \begin{enumerate} 
%     \item einem Runtime-Bundle aus Modell, Kandidatenprofilen, Featuredefinition und Laufzeitkonfiguration, 
%     \item einem Nachweispaket aus Run-Manifest, Trainingszusammenfassung, Metriken, Prüfberichten und Vorhersagedateien. 
% \end{enumerate} 

% Nicht verwendete Konfigurationsalternativen und ältere Modellkomponenten sind nicht Bestandteil von A4. 

% \subsection{Zusammenhang mit der bisherigen Arbeit} 
% \label{subsec:zusammenhang_artefaktueberfuehrung} 
% Artefakt A1 stellt die aufbereiteten Daten und Datensplits bereit. Artefakt A2 legt LambdaMART als Rankingverfahren fest. Artefakt A3 definiert die erforderlichen Datenrepräsentationen, 
% Komponenten und Datenflüsse. Kapitel~7 ordnet diese Vorgaben den implementierten Modulen, Funktionen, Konfigurationen und Dateien zu. Das daraus erzeugte Artefakt A4 bildet die technische Grundlage für die Evaluation in Kapitel~8. 

% \subsection{Darstellung der Artefaktüberführung} 
% \label{subsec:darstellung_artefaktueberfuehrung} 
% Abbildung~\ref{fig:artefaktueberfuehrung_implementierung} zeigt die Eingangsartefakte, ihre Verarbeitung durch die Pipeline und die Bestandteile des erzeugten Artefakts A4. 
% \begin{figure}[H] \centering \resizebox{\textwidth}{!}{
%     \begin{tikzpicture}[ >=Stealth, font=\small, input/.style={ draw=blue!65!black, fill=blue!8, rounded corners=2mm, text width=4.0cm, minimum height=1.35cm, align=center, inner sep=6pt }, 
%         process/.style={ draw=green!55!black, fill=green!9, rounded corners=2mm, text width=3.5cm, minimum height=1.35cm, align=center, inner sep=6pt }, 
%         output/.style={ draw=orange!75!black, fill=orange!12, rounded corners=2mm, text width=5.2cm, minimum height=1.5cm, align=center, inner sep=7pt }, 
%         next/.style={ draw=gray!70!black, fill=gray!8, rounded corners=2mm, text width=5.0cm, minimum height=1.25cm, align=center, inner sep=6pt }, 
%         flow/.style={ ->, thick }, nextflow/.style={ ->, thick, dashed } ] 
%         % Eingangsartefakte 
%         \node[input] (a1) at (-5.5,5.5) { \textbf{A1}\\ Datenbasis und Datensplits }; \node[input] (a2) at (0,5.5) { \textbf{A2}\\ Modellansatz LambdaMART }; \node[input] (a3) at (5.5,5.5) { \textbf{A3}\\ Systemarchitektur }; 
%         % Gemeinsamer Eintrittspunkt 
%         \coordinate (join) at (0,4.1); \draw[thick] (a1.south) |- (join); \draw[thick] (a2.south) -- (join); \draw[thick] (a3.south) |- (join); 
%         % Pipeline 
%         \node[process] (s910) at (-5.7,2.4) { \textbf{Stufe 910}\\ Datenaufbereitung }; \node[process] (s920) at (-1.9,2.4) { \textbf{Stufe 920}\\ Datenaufteilung }; \node[process] (s930) at (1.9,2.4) { \textbf{Stufe 930}\\ Modelltraining }; \node[process] (s940) at (5.7,2.4) { \textbf{Stufe 940}\\ technische Evaluation }; \coordinate (pipelineEntry) at (-7.5,4.1); \draw[flow] (join) -- (pipelineEntry) |- (s910.north); \draw[flow] (s910.east) -- (s920.west); \draw[flow] (s920.east) -- (s930.west); \draw[flow] (s930.east) -- (s940.west); 
%         % A4 
%         \coordinate (resultJoin) at (0,1.0); \draw[thick] (s930.south) |- (resultJoin); \draw[thick] (s940.south) |- (resultJoin); \node[output] (runtime) at (-3.0,-0.5) { \textbf{Runtime-Bundle}\\ Modell, Kandidatenprofile, Featuredefinition und Laufzeitkonfiguration }; \node[output] (evidence) at (3.0,-0.5) { \textbf{Nachweispaket}\\ Manifest, Zusammenfassung, Metriken, Prüfberichte und Vorhersagen }; \draw[flow] (resultJoin) -| (runtime.north); \draw[flow] (resultJoin) -| (evidence.north); \node[ draw=orange!75!black, rounded corners=3mm, fit=(runtime)(evidence), inner sep=11pt, label={ [font=\bfseries, text=orange!75!black] below:Artefakt A4 } ] (a4) {}; 
%         % Kapitel 8 
%         \node[next] (chapter8) at (0,-3.6) { \textbf{Kapitel 8}\\ Evaluation und Interpretation }; \draw[nextflow] (a4.south) -- (chapter8.north); 
%     \end{tikzpicture}} 
%     \caption{ Überführung der Artefakte A1, A2 und A3 durch die Pipeline-Stufen 910 bis 940 in das Artefakt A4. } \label{fig:artefaktueberfuehrung_implementierung} 
% \end{figure}

% \section{Implementierungsumgebung und Softwarestruktur} 
% \label{sec:implementierungsumgebung_und_softwarestruktur} 

% Die Modellimplementierung ist als modular aufgebaute Python-Anwendung realisiert. Die Softwarestruktur trennt Konfiguration, Datenmodelle, Datenzugriff, Pipelineverarbeitung, Benutzeranwendungen und Hilfswerkzeuge. 
% Dadurch können die einzelnen Verarbeitungsschritte unabhängig ausgeführt, geprüft und erweitert werden. 

% \subsection{Technische Schwerpunkte} 
% \label{subsec:technische_schwerpunkte_softwarestruktur} 

% Die Implementierungsumgebung umfasst folgende Bestandteile: 
% \begin{itemize} 
%     \item paketbasierte Gliederung der Python-Anwendung, 
%     \item zentrale Verwaltung der Projekt- und Datenpfade, 
%     \item YAML-basierte Feature-, Modell-, Pfad- und Trainingskonfiguration, 
%     \item Datenmodelle zur einheitlichen Repräsentation der Artikelinformationen, 
%     \item getrennte Module für Datenzugriff, Pipeline und Anwendungen, 
%     \item Auswahl unterschiedlicher Datenbestände über Umgebungsvariablen, 
%     \item dateibasierte Persistierung der Zwischen- und Ergebnisartefakte. 
% \end{itemize} 

% \subsection{Technische Umsetzung} 
% \label{subsec:technische_umsetzung_softwarestruktur} 

% Die Implementierung gliedert sich in mehrere Python-Pakete mit getrennten Verantwortlichkeiten: 
% \begin{description} 
%     \item[\texttt{pkg\_100\_config}] enthält die zentrale Konfigurations- und Pfadverwaltung. 
%     \item[\texttt{pkg\_200\_schemas}] definiert die Datenmodelle der verarbeiteten Objekte. 
%     \item[\texttt{pkg\_300\_data}] stellt Funktionen zum Laden und Speichern der Daten bereit. 
%     \item[\texttt{pkg\_400\_repository}] kapselt den Zugriff auf persistierte Daten- und Modellartefakte. 
%     \item[\texttt{pkg\_900\_pipeline}] enthält die ausführbaren Stufen der Datenaufbereitung, Datenaufteilung, des Modelltrainings und der Evaluation. 
%     \item[\texttt{pkg\_1000\_apps}] stellt Anwendungen zur Ausführung, Überwachung und Auswertung der Pipeline bereit. 
%     \item[\texttt{pkg\_1100\_tools}] enthält ergänzende Werkzeuge, beispielsweise für den Import neuer Makros. 
% \end{description} 

% Die grundlegenden Projektpfade werden in \texttt{cfg\_100\_settings.py} definiert. Dazu gehören die Verzeichnisse für Rohdaten, aufbereitete Daten, Modelle, Lookup-Dateien, Trainingsläufe und Evaluationsergebnisse. 
% Die Pipeline verwendet dadurch keine dezentral im Programmcode hinterlegten absoluten Pfade. Der Zugriff auf die Konfiguration erfolgt über \texttt{cfg\_120\_config\_loader.py}. 
% Der dort implementierte Konfigurationslader liest die YAML-Dateien ein und stellt ihre Inhalte den Pipelinekomponenten zur Verfügung. Bereits geladene Konfigurationen werden zwischengespeichert, 
% sodass sie während eines Pipeline-Laufs konsistent verwendet werden. Die Konfiguration ist insbesondere auf folgende Dateien verteilt: 
% \begin{itemize} 
%     \item \texttt{cfg\_130\_feature\_config.yaml} für die Merkmalsverarbeitung, 
%     \item \texttt{cfg\_140\_model\_config.yaml} für modellbezogene Einstellungen, 
%     \item \texttt{cfg\_150\_paths.yaml} für die Verzeichnisstruktur, 
%     \item weitere Trainingskonfigurationen für die Ausführung und Parametrisierung der Pipeline. 
% \end{itemize} 

% Über die Umgebungsvariable \texttt{KI\_DATASET\_MODE} kann zwischen den Datenbeständen \texttt{current}, \texttt{ideal} und \texttt{synthetic} unterschieden werden. Die Auswahl beeinflusst die verwendeten Eingabepfade, 
% ohne die ausführenden Pipeline-Module zu verändern. Die Datenübergabe zwischen den Verarbeitungsschritten erfolgt überwiegend dateibasiert. Strukturierte Einzelobjekte und gesammelte Datenbestände werden als JSON- beziehungsweise JSONL-Dateien gespeichert.
% Modell-, Konfigurations- und Evaluationsartefakte werden in laufbezogenen Ergebnisverzeichnissen abgelegt. Diese Softwarestruktur ermöglicht eine sequenzielle Ausführung der Pipeline, ohne die einzelnen Stufen fest miteinander zu koppeln. 
% Jede Stufe liest definierte Eingabeartefakte, verarbeitet diese und stellt ihre Ergebnisse für die nachfolgende Stufe bereit. 

% \subsection{Abgleich zwischen bisheriger Arbeit und Code} 
% \label{subsec:abgleich_softwarestruktur} 

% \subsubsection*{Übereinstimmungen} 
% Die in Kapitel~6 vorgesehene Trennung von Datenverarbeitung, Modelltraining, Evaluation und Anwendung ist in der Paketstruktur umgesetzt. 
% Auch die definierten Datenflüsse werden durch standardisierte Dateien und getrennte Pipeline-Stufen abgebildet. Die modulare Struktur unterstützt die bereits formulierten Anforderungen an Wartbarkeit, 
% Erweiterbarkeit und reproduzierbare Ausführung. 

% \subsubsection*{Technische Konkretisierungen} 
% Die bisherige Arbeit beschreibt die Systemkomponenten auf architektonischer Ebene. Der Code konkretisiert diese durch Konfigurations-, Schema-, Datenzugriffs-, Repository-, Pipeline-, Anwendungs- und Werkzeugpakete. 
% Die Konfiguration wird nicht durch eine einzelne Datei bereitgestellt. Sie setzt sich aus Python-Einstellungen, YAML-Dateien, Umgebungsvariablen und einzelnen pipelineinternen Konstanten zusammen. Welche Parameter verwendet werden, 
% ergibt sich daher aus der für den jeweiligen Lauf wirksamen Gesamtkonfiguration. Die Implementierung ergänzt die konzipierte Architektur zudem um Anwendungs- und Werkzeugmodule. Diese übernehmen die Pipelineausführung, 
% die Anzeige von Trainingsläufen sowie den Import neuer Makros. \subsubsection*{Abweichungen} Die bisherige Darstellung der Systemarchitektur vermittelt eine weitgehend zentrale Konfigurationsverwaltung. 

% Im Code ist die Konfiguration dagegen auf mehrere Ebenen verteilt: 

% \begin{itemize} 
%     \item Python-Konstanten und Pfade, 
%     \item YAML-Konfigurationsdateien, 
%     \item Umgebungsvariablen, 
%     \item pipelineinterne Modellparameter. 
% \end{itemize} 

% Dadurch können sich unterschiedliche Parameterquellen überlagern. Für die Reproduzierbarkeit ist deshalb nicht nur eine einzelne Konfigurationsdatei, sondern die vollständige Laufkonfiguration maßgeblich. 
% Die Datei \texttt{cfg\_140\_model\_config.yaml} enthält neben LambdaMART-bezogenen Einstellungen auch Parameter früherer oder alternativer Modellansätze. Dazu zählen Konfigurationen für neuronale Repräsentationen. 
% Diese Parameter sind nicht Bestandteil des finalen LambdaMART-Modells, sofern sie während des konkreten Trainingslaufs nicht verwendet werden. Darüber hinaus unterstützt der Code mehrere Datenmodi. 
% Die bisherige Arbeit beschreibt dagegen überwiegend einen einheitlichen Datenbestand. Ergebnisse verschiedener Datenmodi dürfen daher nicht ohne Kennzeichnung miteinander verglichen werden. 

% \subsubsection*{Konsequenz für die Dokumentation} 

% Für Kapitel~7 werden ausschließlich die Module und Konfigurationsparameter beschrieben, die im finalen LambdaMART-Trainingslauf tatsächlich verwendet werden. 
% Legacy-Konfigurationen und alternative Modellansätze werden nicht als Bestandteil des Artefakts A4 dargestellt. Jeder reproduzierbare Trainingslauf muss mindestens folgende Informationen eindeutig festhalten: 
% \begin{itemize} 
%     \item ausgewählter Datenmodus, 
%     \item verwendete Eingabe- und Ausgabepfade, 
%     \item aktive Featurekonfiguration, 
%     \item aktive LambdaMART-Konfiguration, 
%     \item Zufallsseed und Splitkonfiguration, 
%     \item Version beziehungsweise Kennung des Trainingslaufs. 
% \end{itemize} 

% \subsection{Zusammenhang mit der bisherigen Arbeit} 
% \label{subsec:zusammenhang_softwarestruktur} 

% Die Softwarestruktur überführt die in Kapitel~6 definierten Systemkomponenten in ausführbare Module. Die dort fachlich beschriebenen Datenobjekte werden durch das Schemapaket abgebildet. 
% Die konzipierten Verarbeitungsschritte werden durch die Pipelinepakete umgesetzt. Die vorgesehenen Datenablagen entsprechen den dateibasierten Zwischen- und Ergebnisartefakten. Artefakt A1 wird über die Daten- und Repositorymodule eingelesen. 
% Die Konfiguration des in A2 ausgewählten LambdaMART-Modells erfolgt über die Modell- und Trainingsparameter. Die in A3 definierte Architektur bestimmt die Gliederung der Pipeline. 
% Die Softwarestruktur bildet damit den technischen Rahmen für die nachfolgenden Abschnitte. Abschnitt~7.3 beschreibt innerhalb dieses Rahmens die konkrete Implementierung der Datenrepräsentationen. 

% \subsection{Darstellung der Softwarestruktur} 

% \label{subsec:darstellung_softwarestruktur} 
% Abbildung~\ref{fig:implementierungsumgebung_softwarestruktur} zeigt die Beziehung zwischen Konfiguration, Programmpaketen, Pipeline-Stufen und erzeugten Artefakten. Die gestrichelte Verbindung kennzeichnet, 
% dass die aktive Laufkonfiguration aus mehreren Konfigurationsquellen gebildet wird. 

% \begin{figure}[H] 
%     \centering \resizebox{\textwidth}{!}{ 
%     \begin{tikzpicture}[ >=Stealth, font=\small, config/.style={ draw=blue!65!black, fill=blue!8, rounded corners=2mm, text width=4.2cm, minimum height=1.35cm, align=center, inner sep=6pt }, 
%         package/.style={ draw=gray!70!black, fill=gray!8, rounded corners=2mm, text width=3.3cm, minimum height=1.3cm, align=center, inner sep=6pt }, 
%         pipeline/.style={ draw=green!55!black, fill=green!9, rounded corners=2mm, text width=3.5cm, minimum height=1.35cm, align=center, inner sep=6pt }, 
%         artifact/.style={ draw=orange!75!black, fill=orange!12, rounded corners=2mm, text width=5.0cm, minimum height=1.4cm, align=center, inner sep=7pt }, 
%         warning/.style={ draw=red!65!black, dashed, fill=red!4, rounded corners=2mm, text width=4.5cm, minimum height=1.1cm, align=center, inner sep=6pt }, 
%         flow/.style={ ->, thick }, configflow/.style={ ->, thick, dashed, blue!65!black } ] 
%         % Konfigurationsquellen 
%         \node[config] (pythonConfig) at (-5.5,6.0) { \textbf{Python-Konfiguration}\\ \texttt{cfg\_100\_settings.py} }; \node[config] (yamlConfig) at (0,6.0) { \textbf{YAML-Konfiguration}\\ Features, Modelle, Pfade\\ und Training }; \node[config] (environment) at (5.5,6.0) { \textbf{Umgebungsvariablen}\\ \texttt{KI\_DATASET\_MODE} }; 
%         % Konfigurationslader 
%         \node[package] (loader) at (0,4.2) { \textbf{Konfigurationszugriff}\\ \texttt{cfg\_120\_config\_loader.py} }; \draw[configflow] (pythonConfig.south) |- (loader.west); \draw[configflow] (yamlConfig.south) -- (loader.north); \draw[configflow] (environment.south) |- (loader.east); 
%         % Kernpakete 
%         \node[package] (schemas) at (-5.7,2.2) { \textbf{Schemas}\\ \texttt{pkg\_200\_schemas} }; \node[package] (data) at (-1.9,2.2) { \textbf{Datenzugriff}\\ \texttt{pkg\_300\_data}\\ \texttt{pkg\_400\_repository} }; \node[package] (pipelinePackage) at (1.9,2.2) { \textbf{Pipeline}\\ \texttt{pkg\_900\_pipeline} }; \node[package] (applications) at (5.7,2.2) { \textbf{Anwendungen und Werkzeuge}\\ \texttt{pkg\_1000\_apps}\\ \texttt{pkg\_1100\_tools} }; \coordinate (packageJoin) at (0,3.2); \draw[flow] (loader.south) -- (packageJoin); \draw[flow] (packageJoin) -| (schemas.north); \draw[flow] (packageJoin) -| (data.north); \draw[flow] (packageJoin) -| (pipelinePackage.north); \draw[flow] (packageJoin) -| (applications.north); 
%         % Pipeline-Stufen 
%         \node[pipeline] (stages) at (0,0.1) { \textbf{Ausführbare Pipeline}\\ Stufen 910, 920, 930 und 940 }; \coordinate (coreJoin) at (0,1.1); \draw[thick] (schemas.south) |- (coreJoin); \draw[thick] (data.south) |- (coreJoin); \draw[thick] (pipelinePackage.south) |- (coreJoin); \draw[thick] (applications.south) |- (coreJoin); \draw[flow] (coreJoin) -- (stages.north); 
%         % Ergebnis 
%         \node[artifact] (artifacts) at (0,-2.0) { \textbf{Laufbezogene Artefakte}\\ JSON-/JSONL-Daten, Kandidatenprofile, Modell, Manifest und Evaluationsdateien }; \draw[flow] (stages.south) -- (artifacts.north); 
%         % Abweichung 
%         \node[warning] (difference) at (6.2,-2.0) { Die aktive Konfiguration entsteht aus mehreren Quellen; vorhandene Legacy-Parameter sind davon abzugrenzen. }; \draw[ ->, thick, dashed, red!65!black ] (difference.west) -- (artifacts.east); \end{tikzpicture} } 
%         \caption{ Aufbau der Implementierungsumgebung aus Konfigurationsquellen, Programmpaketen, Pipeline-Stufen und laufbezogenen Artefakten. } \label{fig:implementierungsumgebung_softwarestruktur} \end{figure}



% \section{Technische Umsetzung der Datenrepräsentationen} 
% \label{sec:technische_umsetzung_der_datenrepraesentationen} 

% Die Pipeline überführt die aufbereiteten Artikelinformationen in einheitliche, maschinenlesbare Datenobjekte. Dabei werden allgemeine Artikeldaten, Makrostrukturmerkmale, Artikelmerkmale und lexikalische Merkmale getrennt verarbeitet. Diese Repräsentationen bilden die Ausgangsbasis für die Kandidatenprofile und die spätere Berechnung der Anfrage-Kandidat-Merkmale. 

% \subsection{Technische Schwerpunkte} 
% \label{subsec:technische_schwerpunkte_datenrepraesentationen} 

% Die technische Datenrepräsentation umfasst: \begin{itemize} \item ein einheitliches Artikeldatenmodell, \item dimensionsbezogene Merkmalsvektoren, \item getrennte Makrostruktur-, Artikel- und Lexikalrepräsentationen, \item Fingerprints zur Erkennung zusammengehöriger Datensätze, \item Merkmals- und Dimensionsstatistiken, \item Ausschluss prozessbezogener und identifizierender Metadaten, \item Bereinigung redundanter Merkmalsdimensionen, \item Speicherung als JSON- und JSONL-Dateien. \end{itemize} 

% \subsection{Technische Umsetzung} 
% \label{subsec:technische_umsetzung_datenrepraesentationen} 

% Die zentrale Repräsentation eines Artikels wird durch das in \texttt{schema\_210\_article.py} definierte Datenmodell \texttt{Article} bereitgestellt. Dieses enthält insbesondere: \begin{itemize} \item eine eindeutige Artikel-ID, \item den optionalen Pfad zur zugehörigen EMA-Datei, \item einen dimensionsbezogenen Merkmalsvektor, \item die Anzahl der enthaltenen Merkmale und Dimensionen, \item einen Artikel-Fingerprint, \item ergänzende Metadaten. \end{itemize} Der dimensionsbezogene Merkmalsvektor ordnet jede Information einer definierten Merkmalsdimension zu. Dadurch bleiben Merkmalsart und Merkmalswert voneinander getrennt. Die Pipeline kann diese Informationen dimensionsabhängig bereinigen, aggregieren und vergleichen. Für die weitere Verarbeitung werden die relevanten Informationen nach ihrer Funktion in drei Rollen gegliedert: \begin{description} \item[Makrostrukturrepräsentation] enthält technische Funktions-, Anschluss- und Strukturmerkmale. \item[Artikelrepräsentation] enthält die technischen und physikalischen Eigenschaften des Artikels. \item[Lexikalrepräsentation] enthält kontrollierte Produkt-, Typ-, Kommunikations- und Technologiemerkmale. \end{description} Diese Trennung wird im Modelltraining durch die Datenobjekte \texttt{macro\_shape\_vector}, \texttt{article\_vector} und \texttt{lexical\_vector} abgebildet. Die Rollenvektoren enthalten keine unmittelbar verwendbaren Rangwerte. Sie stellen strukturierte Ausgangsdaten für die spätere Bildung der Vergleichsmerkmale dar. Vor der Speicherung werden redundante Dimensionen durch die entsprechenden Bereinigungsfunktionen entfernt. Dies verhindert, dass dieselbe Information gleichzeitig im allgemeinen Vektor und in einem spezialisierten Rollenvektor enthalten ist. Zusätzlich werden Prozess-, Bearbeiter-, Dokumentpfad- und Identitätsinformationen von den Modellmerkmalen getrennt. Die Persistierung erfolgt über die Datenzugriffskomponenten des Pakets \texttt{pkg\_300\_data}. Der dort verwendete \texttt{DataLoader} stellt Funktionen zum Lesen und Schreiben von JSON- und JSONL-Dateien bereit. Einzelne Artikel können dadurch als separate Dateien gespeichert werden, während zusammengefasste Exportdateien den gebündelten Zugriff auf einen Datenbestand ermöglichen. Die gesammelten Dateien dienen ausschließlich dem Export und der technischen Prüfung. Die Datenaufbereitung und das Modelltraining verwenden die dafür vorgesehenen Einzeldateien beziehungsweise die explizit erzeugten Pipelineartefakte. Dadurch wird vermieden, dass unterschiedliche Datenstände aus Einzel- und Sammeldateien gleichzeitig in einen Verarbeitungsschritt eingehen. Der Fingerprint eines Artikels dient der Erkennung technisch zusammengehöriger Datensätze. Er wird insbesondere bei der gruppierten Datenaufteilung verwendet, damit gleichartige oder inhaltlich abhängige Artikel nicht auf unterschiedliche Datensplits verteilt werden. 
% \subsection{Abgleich zwischen bisheriger Arbeit und Code} 
% \label{subsec:abgleich_datenrepraesentationen} 
% \subsubsection*{Übereinstimmungen} Die bisherige Arbeit und der Code unterscheiden übereinstimmend zwischen Artikelmerkmalen, Makrostrukturmerkmalen und lexikalischen Merkmalen. Auch die dimensionsbezogene Speicherung der technischen Eigenschaften sowie der Ausschluss nicht modellgeeigneter Metadaten entsprechen der bisherigen Konzeption. Die Fingerprints konkretisieren die bereits formulierte Anforderung, zusammengehörige Artikel, Varianten und technisch abhängige Datensätze innerhalb desselben Splits zu halten. \subsubsection*{Technische Konkretisierungen} Die bisherige Arbeit beschreibt die Merkmalsrepräsentationen überwiegend fachlich. Der Code konkretisiert sie durch das \texttt{Article}-Datenmodell, rollenbezogene Vektoren, dimensionsbezogene Wertelisten und explizite Lade- und Speicherfunktionen. Zusätzlich enthält das Artikeldatenmodell technische Hilfsfelder wie Merkmalsanzahlen, Dimensionsanzahlen, Fingerprints und Metadaten. Diese Felder unterstützen die Datenprüfung und Pipelineausführung, sind jedoch nicht automatisch Eingabemerkmale des Rankingmodells. Der Code trennt außerdem zwischen der serialisierbaren Datenrepräsentation und den für das Training vorbereiteten Laufzeitstrukturen. Die gespeicherten JSON-Objekte werden beim Laden in die für Profilbildung und Merkmalsvergleich benötigten Datenstrukturen überführt. \subsubsection*{Abweichungen} Die bisherige Arbeit bezeichnet Artikel-, Makro- und Lexikalvektoren teilweise als unmittelbare numerische Eingabe des LambdaMART-Modells. Dies entspricht nicht der tatsächlichen Implementierung. LambdaMART verarbeitet keine direkte Verkettung dieser Repräsentationen. Die Rollenvektoren dienen zunächst als strukturierte Ausgangsdaten. Die eigentliche Modelleingabe entsteht erst durch die Berechnung numerischer Vergleichsmerkmale für ein Anfrage-Kandidat-Paar: \begin{equation} \left( \text{Anfragerepräsentationen}, \text{Kandidatenrepräsentationen} \right) \xrightarrow{\text{Merkmalsvergleich}} \boldsymbol{x}_{q,c}. \label{eq:bildung_paarmerkmalsvektor} \end{equation} Dabei bezeichnet $\boldsymbol{x}_{q,c}$ den numerischen Merkmalsvektor für die Anfrage $q$ und den Kandidaten $c$. Erst dieser Paarmerkmalsvektor wird als Eingabe an LambdaMART übergeben. Eine weitere Abweichung besteht in der technischen Behandlung der Metadaten. Das \texttt{Article}-Objekt kann Dokumentpfade, Bearbeitungsinformationen und weitere Prozessangaben enthalten. Deren Vorhandensein im Datenobjekt bedeutet jedoch nicht, dass sie als Modellmerkmale verwendet werden. Die bisherige Arbeit muss deshalb zwischen gespeicherten Begleitinformationen und trainierbaren Merkmalen unterscheiden. Auch die Bezeichnung \textit{Vektor} ist technisch zu präzisieren. Bei den gespeicherten Artikel-, Makrostruktur- und Lexikalrepräsentationen handelt es sich überwiegend um dimensionsbezogene Mengen oder Listen von Merkmalswerten. Ein vollständig numerischer Vektor mit fester Reihenfolge entsteht erst bei der Erzeugung der Anfrage-Kandidat-Matrix. \subsubsection*{Konsequenz für die Dokumentation} Die Datenrepräsentationen werden in Kapitel~7 als strukturierte Zwischenrepräsentationen bezeichnet. Folgende Begriffe sind zu unterscheiden: \begin{description} \item[Rollenrepräsentation] Dimensionsbezogene Artikel-, Makrostruktur- oder Lexikalinformationen. \item[Profilrepräsentation] Aggregierte Rollenrepräsentationen eines Kandidaten. \item[Paarmerkmalsvektor] Numerischer Vektor mit fester Feature-Reihenfolge für ein Anfrage-Kandidat-Paar. \item[Modelleingabe] Matrix der numerischen Paarmerkmalsvektoren einschließlich der zugehörigen Anfragegruppen und Relevanzlabels. \end{description} Nur die Paarmerkmalsvektoren sind direkte Eingaben des LambdaMART-Modells. Gespeicherte Metadaten, Identitätsinformationen und Prozessfelder bleiben von der Modelleingabe ausgeschlossen. \subsection{Zusammenhang mit der bisherigen Arbeit} \label{subsec:zusammenhang_datenrepraesentationen} Die Datenrepräsentationen setzen die in den vorherigen Kapiteln definierten Merkmalsgruppen technisch um. Artefakt A1 stellt dafür die bereinigten Artikelinformationen und die erforderlichen Fingerprints bereit. Artefakt A3 legt die Trennung zwischen Artikel-, Makrostruktur- und Lexikalmerkmalen fest. Die Implementierung überführt diese Vorgaben in das \texttt{Article}-Schema und die zugehörigen Rollenrepräsentationen. Damit entsteht die technische Ausgangsbasis für die in Abschnitt~\ref{sec:technische_bildung_der_kandidatenprofile} beschriebene Bildung der Kandidatenprofile. Die Datenrepräsentationen sind noch keine unmittelbaren LambdaMART-Eingaben. Ihre Überführung in numerische Anfrage-Kandidat-Merkmale erfolgt erst beim Aufbau der Rankingmatrizen. \subsection{Darstellung der Datenrepräsentationen} \label{subsec:darstellung_datenrepraesentationen} Abbildung~\ref{fig:technische_datenrepraesentationen} zeigt die Überführung eines aufbereiteten Artikels in das technische Artikeldatenmodell. Aus diesem werden die drei Rollenrepräsentationen gebildet. Die dargestellte Trennung verdeutlicht, dass die Rollenrepräsentationen zunächst gespeichert und erst in einem nachfolgenden Schritt in numerische Paarmerkmale überführt werden. 
% \begin{figure}[H] \centering \resizebox{\textwidth}{!}{ \begin{tikzpicture}[ >=Stealth, font=\small, input/.style={ draw=blue!65!black, fill=blue!8, rounded corners=2mm, text width=4.4cm, minimum height=1.35cm, align=center, inner sep=6pt }, schema/.style={ draw=gray!70!black, fill=gray!8, rounded corners=2mm, text width=5.1cm, minimum height=1.55cm, align=center, inner sep=7pt }, role/.style={ draw=green!55!black, fill=green!9, rounded corners=2mm, text width=4.1cm, minimum height=1.45cm, align=center, inner sep=6pt }, metadata/.style={ draw=red!65!black, fill=red!4, dashed, rounded corners=2mm, text width=4.3cm, minimum height=1.25cm, align=center, inner sep=6pt }, storage/.style={ draw=orange!75!black, fill=orange!12, rounded corners=2mm, text width=4.7cm, minimum height=1.4cm, align=center, inner sep=7pt }, next/.style={ draw=violet!65!black, fill=violet!7, rounded corners=2mm, text width=5.0cm, minimum height=1.4cm, align=center, inner sep=7pt }, flow/.style={ ->, thick }, exclusion/.style={ ->, thick, dashed, red!65!black } ] 
%     % Eingang 
%     \node[input] (preparedArticle) at (0,6.1) { \textbf{Aufbereiteter Artikel}\\ kanonisierte Eigenschaften, Fingerprint und Begleitinformationen }; 
%     % Schema 
%     \node[schema] (articleSchema) at (0,4.0) { \textbf{\texttt{Article}-Datenmodell}\\ Artikel-ID, dimensionsbezogener Vektor, Merkmalsanzahlen, Fingerprint und Metadaten }; \draw[flow] (preparedArticle.south) -- (articleSchema.north); 
%     % Rollenrepräsentationen 
%     \node[role] (macroShape) at (-5.2,1.7) { \textbf{Makrostrukturrepräsentation}\\ Funktions-, Anschluss- und Strukturmerkmale }; \node[role] (articleVector) at (0,1.7) { \textbf{Artikelrepräsentation}\\ technische und physikalische Artikeleigenschaften }; \node[role] (lexicalVector) at (5.2,1.7) { \textbf{Lexikalrepräsentation}\\ Produkt-, Typ-, Kommunikations- und Technologiemerkmale }; \coordinate (roleJoin) at (0,2.9); \draw[thick] (articleSchema.south) -- (roleJoin); \draw[flow] (roleJoin) -| (macroShape.north); \draw[flow] (roleJoin) -- (articleVector.north); \draw[flow] (roleJoin) -| (lexicalVector.north); 
%     % Ausgeschlossene Metadaten 
%     \node[metadata] (excludedMetadata) at (7.2,4.0) { \textbf{Nicht trainierbare Begleitdaten}\\ Identitäten, Dokumentpfade, Bearbeiter- und Prozessinformationen }; \draw[exclusion] (articleSchema.east) -- (excludedMetadata.west); 
%     % Persistierung 
%     \node[storage] (storage) at (-3.0,-0.8) { \textbf{Persistierung}\\ Einzeldateien und definierte Pipelineartefakte in JSON/JSONL }; 
%     % Nachfolgende Verarbeitung 
%     \node[next] (pairFeatures) at (3.0,-0.8) { \textbf{Nachfolgende Verarbeitung}\\ Profilaggregation und Berechnung numerischer Anfrage-Kandidat-Merkmale }; \coordinate (representationJoin) at (0,0.4); \draw[thick] (macroShape.south) |- (representationJoin); \draw[thick] (articleVector.south) -- (representationJoin); \draw[thick] (lexicalVector.south) |- (representationJoin); \draw[flow] (representationJoin) -| (storage.north); \draw[flow] (representationJoin) -| (pairFeatures.north); 
%     % Direkte Modelleingabe explizit ausschließen 
%     \node[metadata] (notDirectInput) at (0,-2.8) { \textbf{Abgrenzung}\\ Rollenrepräsentationen sind keine direkte LambdaMART-Eingabe. }; \draw[exclusion] (pairFeatures.south) -- (notDirectInput.north); \end{tikzpicture} } 
%     \caption{ Technische Überführung der aufbereiteten Artikeldaten in das \texttt{Article}-Datenmodell und die getrennten Rollenrepräsentationen. Die Rollenrepräsentationen werden persistiert und für die Profil- und Paarmerkmalsbildung verwendet, jedoch nicht direkt an LambdaMART übergeben. } \label{fig:technische_datenrepraesentationen} \end{figure}


% \section{Technische Bildung der Kandidatenprofile} \label{sec:technische_bildung_der_kandidatenprofile} Die Kandidatenprofile bilden die zur Laufzeit verfügbaren Repräsentationen der Makrokandidaten. Sie werden ausschließlich aus den dem jeweiligen Zielmakro zugeordneten Trainingsartikeln und den zulässigen technischen Makroinformationen erzeugt. Die resultierenden Profile werden serialisiert und anschließend für die Bildung der Anfrage-Kandidat-Merkmale verwendet. \subsection{Technische Schwerpunkte} \label{subsec:technische_schwerpunkte_kandidatenprofile} Die Kandidatenprofilbildung umfasst: \begin{itemize} \item Gruppierung der Trainingsartikel nach dem Zielmakro, \item ausschließliche Verwendung des Trainingsanteils, \item getrennte Aggregation der Makrostruktur-, Artikel- und Lexikalrepräsentationen, \item Ermittlung der Makrofamilie und des Familienpräfixes, \item Berechnung von Support- und Fingerprint-Informationen, \item Begrenzung der gespeicherten Werte je Merkmalsdimension, \item Ausschluss unzulässiger und redundanter Dimensionen, \item Vorbereitung zusätzlicher Laufzeitstrukturen, \item Serialisierung als \texttt{candidate\_profiles.json}. \end{itemize} \subsection{Technische Umsetzung} \label{subsec:technische_umsetzung_kandidatenprofile} Die Profilbildung erfolgt nach der gruppierten Datenaufteilung. Als Eingabe werden ausschließlich die dem Trainingsanteil zugeordneten Artikel verwendet. Validierungs- und Testartikel fließen nicht in die Kandidatenprofile ein. Dadurch bleiben deren Eigenschaften bei der späteren Modellbewertung unbekannt. Die Trainingsartikel werden anhand ihres Zielmakros gruppiert. Für jedes Zielmakro wird ein separates Kandidatenprofil erzeugt. Innerhalb dieser Gruppe werden die zulässigen Merkmale der zugeordneten Artikel dimensionsbezogen gezählt und aggregiert. Die Aggregation wird getrennt für folgende Repräsentationen durchgeführt: \begin{itemize} \item Makrostrukturmerkmale, \item technische Artikelmerkmale, \item lexikalische Merkmale. \end{itemize} Die Funktion \texttt{counter\_profile\_to\_vector()} überführt die während der Aggregation gebildeten dimensionsbezogenen Zähler in die serialisierbare Profilrepräsentation. Dabei werden die Werte einer Dimension entsprechend ihrer Auftretenshäufigkeit geordnet. Die Anzahl der übernommenen Werte wird durch \texttt{PROFILE\_TOP\_VALUES\_PER\_DIMENSION\_DEFAULT} und \texttt{PROFILE\_TOP\_VALUES\_PER\_DIMENSION\_BY\_DIMENSION} begrenzt. Der allgemeine Grenzwert gilt für alle Dimensionen, für die keine spezifische Obergrenze definiert ist. Dimensionsspezifische Grenzwerte ermöglichen eine abweichende Behandlung einzelner Merkmalsarten. Für eine Merkmalsdimension $d$ ergibt sich die gespeicherte Wertemenge damit aus: \begin{equation} P_{m,d} = \operatorname{Top}_{k_d} \left( \left\{ \left(v,n_{m,d,v}\right) \right\} \right), \label{eq:kandidatenprofil_aggregation} \end{equation} wobei $P_{m,d}$ die gespeicherten Werte der Dimension $d$ für das Makro $m$, $n_{m,d,v}$ die Häufigkeit des Wertes $v$ und $k_d$ die für diese Dimension konfigurierte Obergrenze bezeichnet. Das Kandidatenprofil enthält neben den aggregierten Merkmalsrepräsentationen weitere technische Angaben: \begin{description} \item[\texttt{target}] eindeutige Bezeichnung des repräsentierten Zielmakros, \item[\texttt{family}] zugeordnete Makrofamilie, \item[\texttt{family\_prefix}] übergeordnetes Präfix zur Kennzeichnung verwandter Kandidaten, \item[\texttt{support}] Anzahl der für die Profilbildung verwendeten Trainingsartikel, \item[Singleton-Kennzeichen] Kennzeichnung eines Profils, das nur auf einem Trainingsartikel beziehungsweise einer unteilbaren Zusammenhangskomponente basiert, \item[logarithmierter Support] transformierte Supportinformation zur begrenzten Abbildung unterschiedlich großer Kandidatengruppen, \item[Fingerprint-Anzahl] Anzahl unterschiedlicher Artikel-Fingerprints innerhalb des Kandidatenprofils. \end{description} Die Supportangaben ergänzen die technische Profilrepräsentation. Sie ermöglichen die Unterscheidung zwischen Kandidatenprofilen mit umfangreicher Datenbasis und Profilen, die nur aus wenigen Trainingsartikeln gebildet wurden. Die Makroidentität und die Familieninformationen dienen der Profilverwaltung und der Labelzuweisung. Sie dürfen nicht als direkte Identitätsmerkmale in den numerischen Merkmalsvektor übernommen werden. Vor der Serialisierung werden unzulässige, redundante oder einer anderen Merkmalsrolle zugeordnete Dimensionen entfernt. Dadurch wird verhindert, dass dieselbe Information mehrfach im Kandidatenprofil vorliegt oder eine direkte Zielkennung in die späteren Paarvergleichsmerkmale eingeht. Die Funktion \texttt{candidate\_profiles\_for\_serialization()} überführt die internen Profilobjekte in eine speicherbare Struktur. Die Profile werden anschließend in der Datei \texttt{candidate\_profiles.json} abgelegt. Diese Datei ist Bestandteil des Runtime-Bundles und wird gemeinsam mit dem trainierten Modell geladen. Für die anschließende Merkmalsberechnung werden die serialisierten Profile durch \texttt{prepare\_runtime\_candidate\_caches()} in vorbereitete Laufzeitstrukturen überführt. Hierzu gehören insbesondere: \begin{itemize} \item flache Wertemengen für Mengenvergleiche, \item dimensionsbezogene Wertemengen, \item vorberechnete Merkmalsanzahlen, \item Struktur- und Hashinformationen, \item Support- und Familieninformationen. \end{itemize} Diese Laufzeit-Caches vermeiden eine wiederholte Aufbereitung unveränderlicher Kandidateninformationen bei jedem Anfrage-Kandidat-Vergleich. Sie sind technische Hilfsstrukturen und keine zusätzlichen fachlichen Profile. \subsection{Abgleich zwischen bisheriger Arbeit und Code} \label{subsec:abgleich_kandidatenprofile} \subsubsection*{Übereinstimmungen} Die bisherige Arbeit und der Code sehen übereinstimmend vor, die einem Makro zugeordneten Trainingsartikel zu einer gemeinsamen Kandidatenrepräsentation zusammenzuführen. Artikel aus dem Validierungs- und Testanteil werden nicht zur Profilbildung verwendet. Die Trennung zwischen Makrostruktur-, Artikel- und Lexikalmerkmalen bleibt bei der Aggregation erhalten. Ebenso werden Merkmale ausgeschlossen, die eine direkte Identifikation des Zielmakros ermöglichen. \subsubsection*{Technische Konkretisierungen} Das in der bisherigen Arbeit fachlich definierte Makroprofil wird im Code als \texttt{candidate\_profile} umgesetzt. Das Kandidatenprofil enthält neben den aggregierten Merkmalen zusätzliche Angaben zur technischen Verarbeitung: \begin{itemize} \item Makrofamilie und Familienpräfix, \item Anzahl der zugeordneten Trainingsartikel, \item Singleton-Kennzeichen, \item logarithmierter Support, \item Anzahl unterschiedlicher Fingerprints, \item vorbereitete Laufzeit-Caches. \end{itemize} Diese Angaben konkretisieren die Profilverwaltung, die Relevanzlabelbildung und die spätere Merkmalsberechnung. Sie erweitern das fachliche Makroprofil, verändern jedoch nicht dessen grundlegende Funktion als Repräsentation eines Makrokandidaten. \subsubsection*{Abweichungen} Die bisherige Arbeit vermittelt teilweise den Eindruck, dass alle beobachteten Merkmalswerte der zugeordneten Trainingsartikel vollständig in das Makroprofil übernommen werden. Der Code begrenzt die Anzahl der Profilwerte je Dimension. Bei Überschreitung der konfigurierten Obergrenze werden die am häufigsten beobachteten Werte übernommen. Seltene Ausprägungen können dadurch aus dem Profil entfallen. Das technische Kandidatenprofil ist folglich keine unbegrenzte Vereinigung sämtlicher beobachteter Artikelmerkmale. Eine weitere Abweichung besteht in der verwendeten Benennung. Die Arbeit verwendet den Begriff \textit{Makroprofil}, während der Code das Profil als \texttt{candidate\_profile} bezeichnet. Beide Begriffe beziehen sich auf dieselbe fachliche Entität, unterscheiden sich jedoch hinsichtlich ihrer Abstraktion: \begin{description} \item[Makroprofil] fachliche Beschreibung der für einen Makrokandidaten zusammengeführten Eigenschaften, \item[Kandidatenprofil] technische Repräsentation des Makroprofils einschließlich Support-, Familien-, Fingerprint- und Laufzeitinformationen. \end{description} Die bisherige Beschreibung der Makroprofile enthält zudem nicht alle im Code erzeugten Cache- und Supportinformationen. Diese Zusatzinformationen sind Implementierungsdetails und keine eigenständigen fachlichen Makromerkmale. \subsubsection*{Konsequenz für die Dokumentation} Der Begriff \textit{Makroprofil} wird für die fachliche Beschreibung beibehalten. Bei der technischen Umsetzung wird präzisiert, dass das Makroprofil im Code als Kandidatenprofil realisiert wird. Die Profilaggregation wird als häufigkeitsbasierte und begrenzte Aggregation dokumentiert. Die verwendeten Obergrenzen sind Teil der aktiven Laufkonfiguration und müssen im Run-Manifest beziehungsweise in der Trainingskonfiguration nachvollziehbar gespeichert werden. Support-, Familien- und Fingerprintinformationen sind getrennt von den aggregierten technischen Merkmalen darzustellen. Insbesondere die Zielmakrobezeichnung darf ausschließlich zur Verwaltung, Labelzuweisung und Evaluation verwendet werden. Sie darf nicht als trainierbares Vergleichsmerkmal in die LambdaMART-Matrix eingehen. \subsection{Zusammenhang mit der bisherigen Arbeit} \label{subsec:zusammenhang_kandidatenprofile} Die Kandidatenprofilbildung setzt die in Kapitel~6 definierte Makroprofilbildung technisch um. Artefakt A1 stellt die gruppierten Trainingsartikel und ihre Fingerprints bereit. Artefakt A3 definiert die getrennten Artikel-, Makrostruktur- und Lexikalrepräsentationen sowie die ausschließliche Verwendung des Trainingsanteils. Die in Abschnitt~\ref{sec:technische_umsetzung_der_datenrepraesentationen} beschriebenen Rollenrepräsentationen bilden die Eingabe der Profilaggregation. Die daraus erzeugte Datei \texttt{candidate\_profiles.json} ist Bestandteil des Artefakts A4. Die Kandidatenprofile bilden den festen Kandidatenraum für die nachfolgende Matrixerzeugung. Abschnitt~7.5 beschreibt, wie die Anfragerepräsentationen mit diesen Profilen kombiniert und in numerische Anfrage-Kandidat-Merkmale überführt werden. \subsection{Darstellung der Kandidatenprofilbildung} \label{subsec:darstellung_kandidatenprofilbildung} Abbildung~\ref{fig:technische_kandidatenprofilbildung} zeigt die Gruppierung der Trainingsartikel, die getrennte Aggregation der Rollenrepräsentationen und die Erweiterung um technische Profilinformationen. Die Konfigurationsgrenzen wirken ausschließlich auf die Auswahl der gespeicherten Profilwerte. 

% \begin{figure}[H] \centering \resizebox{\textwidth}{!}{\begin{tikzpicture}[ >=Stealth, font=\small, input/.style={ draw=blue!65!black, fill=blue!8, rounded corners=2mm, text width=4.4cm, minimum height=1.35cm, align=center, inner sep=6pt }, representation/.style={ draw=gray!70!black, fill=gray!8, rounded corners=2mm, text width=3.9cm, minimum height=1.35cm, align=center, inner sep=6pt }, process/.style={ draw=green!55!black, fill=green!9, rounded corners=2mm, text width=4.6cm, minimum height=1.4cm, align=center, inner sep=7pt }, config/.style={ draw=red!65!black, fill=red!4, dashed, rounded corners=2mm, text width=4.4cm, minimum height=1.25cm, align=center, inner sep=6pt }, profile/.style={ draw=orange!75!black, fill=orange!12, rounded corners=2mm, text width=5.1cm, minimum height=1.55cm, align=center, inner sep=7pt }, runtime/.style={ draw=violet!65!black, fill=violet!7, rounded corners=2mm, text width=4.8cm, minimum height=1.35cm, align=center, inner sep=6pt }, flow/.style={ ->, thick }, control/.style={ ->, thick, dashed, red!65!black } ]
%     % Trainingsdaten 
%     \node[input] (trainingArticles) at (0,6.6) { \textbf{Trainingsartikel eines Zielmakros}\\ gruppiert nach Zielmakro; keine Validierungs- oder Testartikel }; 
%     % Rollenrepräsentationen 
%     \node[representation] (macroShape) at (-5.1,4.3) { \textbf{Makrostruktur}\\ Funktions-, Anschluss- und Strukturmerkmale }; \node[representation] (articleFeatures) at (0,4.3) { \textbf{Artikelmerkmale}\\ technische und physikalische Eigenschaften }; \node[representation] (lexicalFeatures) at (5.1,4.3) { \textbf{Lexikalmerkmale}\\ Produkt-, Typ-, Kommunikations- und Technologiemerkmale }; \coordinate (roleSplit) at (0,5.5); \draw[thick] (trainingArticles.south) -- (roleSplit); \draw[flow] (roleSplit) -| (macroShape.north); \draw[flow] (roleSplit) -- (articleFeatures.north); \draw[flow] (roleSplit) -| (lexicalFeatures.north); 
%     % Aggregation 
%     \node[process] (aggregation) at (0,2.1) { \textbf{Dimensionsbezogene Profilaggregation}\\ Häufigkeitszählung, Rollenbereinigung und Auswahl der häufigsten Werte }; \coordinate (aggregationJoin) at (0,3.2); \draw[thick] (macroShape.south) |- (aggregationJoin); \draw[thick] (articleFeatures.south) -- (aggregationJoin); \draw[thick] (lexicalFeatures.south) |- (aggregationJoin); \draw[flow] (aggregationJoin) -- (aggregation.north); 
%     % Grenzen 
%     \node[config] (limits) at (6.2,2.1) { \textbf{Profilgrenzen}\\ Standardgrenze und dimensionsspezifische Obergrenzen }; \draw[control] (limits.west) -- (aggregation.east); 
%     % Technische Zusatzinformationen 
%     \node[process] (profileMetadata) at (-5.8,-0.2) { \textbf{Technische Profilinformationen}\\ Familie, Familienpräfix, Support, Singleton-Kennzeichen und Fingerprints }; 
%     % Kandidatenprofil 
%     \node[profile] (candidateProfile) at (0,-2.0) { \textbf{Kandidatenprofil}\\ aggregierte Rollenrepräsentationen und technische Profilinformationen }; \draw[flow] (aggregation.south) -- (candidateProfile.north); \draw[flow] (profileMetadata.south) |- (candidateProfile.west); 
%     % Serialisierung und Cache 
%     \node[profile] (serializedProfiles) at (-3.2,-4.5) { \textbf{\texttt{candidate\_profiles.json}}\\ serialisierte Profile als Bestandteil des Runtime-Bundles }; \node[runtime] (runtimeCaches) at (3.2,-4.5) { \textbf{Laufzeit-Caches}\\ flache Mengen, Dimensionswerte, Zähl- und Strukturinformationen }; \draw[flow] (candidateProfile.south) |- (serializedProfiles.north); \draw[flow] (candidateProfile.south) |- (runtimeCaches.north); \end{tikzpicture} } \caption{ Technische Bildung der Kandidatenprofile aus den Rollenrepräsentationen der Trainingsartikel. Die dimensionsbezogene Aggregation wird durch konfigurierte Profilgrenzen gesteuert und um Support-, Familien- und Fingerprintinformationen ergänzt. } \label{fig:technische_kandidatenprofilbildung} \end{figure}


% \section{Technische Erzeugung der Anfrage-Kandidat-Matrix} \label{sec:technische_erzeugung_der_anfrage_kandidat_matrix} Die Anfrage-Kandidat-Matrix bildet die unmittelbare Eingabe des LambdaMART-Modells. Für jeden Anfrageartikel wird eine Rankinggruppe erzeugt, die den Anfrageartikel mit den verfügbaren Kandidatenprofilen kombiniert. Aus jeder Kombination entsteht eine numerische Featurezeile mit fester Merkmalsreihenfolge. \subsection{Technische Schwerpunkte} \label{subsec:technische_schwerpunkte_anfrage_kandidat_matrix} Die Matrixerzeugung umfasst: \begin{itemize} \item Laden der splitbezogenen Anfrageartikel, \item Ableitung der Anfragerepräsentationen, \item Laden der Kandidatenprofile, \item Kombination jeder Anfrage mit den zugelassenen Kandidaten, \item Berechnung mengenbasierter, struktureller, lexikalischer und numerischer Paarmerkmale, \item Zuordnung eines Relevanzlabels zu jedem Paar, \item Bildung einer Rankinggruppe je Anfrage, \item Festlegung einer einheitlichen Feature-Reihenfolge, \item Erzeugung getrennter Trainings-, Validierungs- und Testmatrizen. \end{itemize} \subsection{Technische Umsetzung} \label{subsec:technische_umsetzung_anfrage_kandidat_matrix} Die Matrixerzeugung verarbeitet die Artikel aus den zuvor gebildeten Trainings-, Validierungs- und Testanteilen getrennt. Jeder Artikel bildet eine Anfrage $q$. Die in \texttt{candidate\_profiles.json} gespeicherten Kandidatenprofile bilden den Kandidatenraum $C$. Für jede Anfrage wird eine Rankinggruppe \begin{equation} G_q = \left\{ \left(q,c\right) \mid c \in C \right\} \label{eq:rankinggruppe_anfrage} \end{equation} erzeugt. Alle Zeilen einer Gruppe beziehen sich damit auf dieselbe Anfrage, unterscheiden sich jedoch hinsichtlich des bewerteten Kandidaten. Die Anfrageseite wird in der Implementierung durch getrennte Repräsentationen abgebildet: \begin{itemize} \item \texttt{query\_macro\_shape} für ableitbare Strukturmerkmale, \item \texttt{query\_article\_vector} für technische Artikeleigenschaften, \item \texttt{query\_lexical\_vector} für kontrollierte lexikalische Merkmale. \end{itemize} Auf der Kandidatenseite werden die entsprechenden Bestandteile des Kandidatenprofils verwendet: \begin{itemize} \item \texttt{candidate\_macro\_shape}, \item \texttt{candidate\_article\_vector}, \item \texttt{candidate\_lexical\_vector}. \end{itemize} Die Repräsentationen werden nicht unmittelbar zusammengefügt. Stattdessen berechnet die Pipeline für jedes Paar $(q,c)$ eine numerische Featurezeile \begin{equation} \boldsymbol{x}_{q,c} = f\left( q_{\mathrm{structure}}, q_{\mathrm{article}}, q_{\mathrm{lexical}}, c_{\mathrm{structure}}, c_{\mathrm{article}}, c_{\mathrm{lexical}} \right), \label{eq:anfrage_kandidat_featurezeile} \end{equation} wobei $f$ die implementierten Vergleichsfunktionen bezeichnet. Die Berechnung umfasst mehrere Merkmalsgruppen: \begin{description} \item[Mengenbasierte Merkmale] bestimmen Überschneidungen, Vereinigungen, Überdeckungsgrade und Jaccard-Ähnlichkeiten zwischen den Wertemengen einer Anfrage und eines Kandidaten. \item[Exakte Merkmale] prüfen die vollständige Übereinstimmung technischer Werte, Strukturkennungen, Schemata und Hashrepräsentationen. \item[Strukturmerkmale] vergleichen Funktionsschablonen, Anschlussrollen, Anschlussstrukturen, IO-Schemata sowie Kanal- und Signalzahlen. \item[Lexikalische Merkmale] vergleichen Produktcodes, Typserien, Kommunikationsstandards, Steckverbinder, Schutzklassen und technische Token. \item[Numerische Merkmale] bilden Differenzen, Abstände und Wertebereichsbeziehungen zwischen typisierten technischen Angaben ab. \item[Aggregierte Merkmale] fassen mehrere Einzelvergleiche zu übergeordneten Übereinstimmungs- und Abdeckungswerten zusammen. \item[Kandidatenbezogene Kontrollmerkmale] bilden Support-, Singleton- und Größeninformationen des Kandidatenprofils ab. \end{description} Mengenbasierte Vergleichsmerkmale werden unter anderem durch \texttt{add\_set\_pair\_features()} erzeugt. Exakte Struktur- und Hashvergleiche werden durch \texttt{add\_exact\_hash\_pair\_features()} ergänzt. Die Einzelfunktionen schreiben ihre Ergebnisse in eine gemeinsame Featurezeile. Für eine Anfragemenge $Q$ und den Kandidatenraum $C$ entsteht die Matrix \begin{equation} X = \left[ \boldsymbol{x}_{q,c} \right]_ {q \in Q,\;c \in C}. \label{eq:anfrage_kandidat_matrix} \end{equation} Jede Zeile wird um das zugehörige Relevanzlabel $y_{q,c}$ ergänzt. Zusätzlich wird die Anzahl der Kandidaten je Anfrage als Gruppeninformation gespeichert. LambdaMART benötigt diese Gruppengrenzen, damit ausschließlich die Kandidaten derselben Anfrage miteinander verglichen werden. Die Featurebezeichnungen werden in \texttt{feature\_names} gespeichert. Ihre Reihenfolge muss für Training, Validierung, Test und spätere Modellanwendung identisch bleiben. Eine veränderte Feature-Reihenfolge würde dazu führen, dass Modellwerte anderen Merkmalen zugeordnet werden. Die vollständige Modelleingabe eines Splits besteht damit aus: \begin{equation} D = \left( X, \boldsymbol{y}, \boldsymbol{g} \right), \label{eq:lambdamart_datensatz} \end{equation} wobei $X$ die Featurematrix, $\boldsymbol{y}$ die Relevanzlabels und $\boldsymbol{g}$ die Größen der Anfragegruppen bezeichnet. Für Training, Validierung und Test werden getrennte Instanzen dieser Datenstruktur erzeugt. Die Kandidatenprofile und die Featuredefinition bleiben dabei unverändert. Unterschiede bestehen ausschließlich in den jeweils als Anfragen verwendeten Artikeln und den daraus berechneten Paarmerkmalen. \subsection{Abgleich zwischen bisheriger Arbeit und Code} \label{subsec:abgleich_anfrage_kandidat_matrix} \subsubsection*{Übereinstimmungen} Die bisherige Arbeit und der Code sehen übereinstimmend vor, jede Anfrage mit mehreren Makrokandidaten zu vergleichen. Für jede Anfrage entsteht eine eigene Rankinggruppe. Der relevante Kandidat wird gemeinsam mit nicht relevanten oder abgestuft relevanten Kandidaten bewertet. Die in Kapitel~6 definierten Vergleichskategorien werden im Code durch numerische Paarmerkmale umgesetzt. Auch die getrennte Verarbeitung der Trainings-, Validierungs- und Testanfragen entspricht der bisherigen Konzeption. \subsubsection*{Technische Konkretisierungen} Die bisherige Arbeit beschreibt die Vergleichsmerkmale überwiegend fachlich. Der Code konkretisiert sie durch einzelne Merkmalsfunktionen und eine feste Feature-Reihenfolge. Zu den implementierten Vergleichen gehören neben mengenbasierten Ähnlichkeiten auch exakte Hashvergleiche, Strukturmerkmale, Zählmerkmale und technische Kontrollmerkmale. Die Anfragerepräsentation wird nicht als einzelner Vektor verarbeitet. Der Code verwendet getrennte Struktur-, Artikel- und Lexikalrepräsentationen. Erst die Vergleichsfunktionen überführen diese Repräsentationen in eine vollständig numerische Featurezeile. Die Implementierung ergänzt außerdem kandidatenbezogene Angaben wie Support, Profilgröße und Singleton-Kennzeichen. Diese Merkmale ermöglichen dem Modell, Unterschiede in der Datengrundlage der Kandidaten zu berücksichtigen. Sie dürfen jedoch keine direkte Makroidentität enthalten. \subsubsection*{Abweichungen} Die bisherige Arbeit beschreibt teilweise separat gespeicherte Anfrageprofile für Training, Validierung und Test. Im Code werden die Anfrageinformationen überwiegend beim Aufbau der jeweiligen Matrix aus den Artikeldaten abgeleitet. Ein eigenständiges Anfrageprofil-Repository unter den bisher verwendeten Bezeichnungen ist für die Matrixerzeugung nicht erforderlich. Zudem kann die bisherige Formulierung den Eindruck erwecken, dass Artikel-, Makro- und Lexikalvektoren direkt an LambdaMART übergeben werden. Tatsächlich verarbeitet das Modell ausschließlich die numerischen Paarmerkmale: \begin{equation} \text{Rollenrepräsentationen} \neq \text{LambdaMART-Eingabe}. \label{eq:rollenrepraesentation_ungleich_modelleingabe} \end{equation} Die direkte Modelleingabe ist die aus den Anfrage-Kandidat-Vergleichen erzeugte Featurematrix. Ein weiterer Prüfpunkt betrifft \texttt{query\_macro\_shape}. Strukturmerkmale auf der Anfrageseite sind nur zulässig, wenn sie ohne Kenntnis des Zielmakros aus dem Artikel beziehungsweise dem zugehörigen Datenblatt abgeleitet werden können. Eine Übernahme aus der bekannten Artikel-Makro-Zuordnung würde eine Zielfleckage verursachen. Die Kandidatenprofile werden ausschließlich aus dem Trainingsanteil gebildet. Für Validierungs- und Testanfragen werden keine splitbezogenen Kandidatenprofile aus den jeweiligen Evaluationsartikeln erzeugt. Damit bleibt der Kandidatenraum während Modellauswahl und Test unverändert. \subsubsection*{Konsequenz für die Dokumentation} Die bisher verwendeten Anfrageprofile werden als logische Anfragerepräsentationen beschrieben. Eine dauerhafte Speicherung ist nicht erforderlich, sofern die Repräsentationen reproduzierbar aus den splitbezogenen Artikeldaten erzeugt werden. Für alle Anfragemerkmale muss dokumentiert werden, aus welcher Eingabequelle sie stammen. Zulässig sind ausschließlich Informationen, die auch für eine neue Anfrage verfügbar oder daraus extrahierbar sind. Dazu zählen insbesondere: \begin{itemize} \item technische Angaben des Anfrageartikels, \item aus dem Datenblatt extrahierte Eigenschaften, \item aus diesen Angaben ableitbare Strukturmerkmale, \item kontrollierte lexikalische Merkmale. \end{itemize} Nicht zulässig sind: \begin{itemize} \item das bekannte Zielmakro, \item die Artikel-Makro-Zuordnung, \item aus dem Zielmakro übernommene Strukturmerkmale, \item Makronamen oder direkt daraus abgeleitete Token, \item Prozess- und Identitätsinformationen. \end{itemize} Die Featuredefinition und ihre Reihenfolge sind als Bestandteil des Artefakts A4 zu speichern. Modell und Featuredefinition dürfen nicht unabhängig voneinander versioniert oder ausgetauscht werden. \subsection{Zusammenhang mit der bisherigen Arbeit} \label{subsec:zusammenhang_anfrage_kandidat_matrix} Die Matrixerzeugung setzt die in Kapitel~6 definierten Anfrage-Kandidat-Vergleiche technisch um. Die in Abschnitt~\ref{sec:technische_umsetzung_der_datenrepraesentationen} beschriebenen Rollenrepräsentationen bilden die Anfrageseite. Die in Abschnitt~\ref{sec:technische_bildung_der_kandidatenprofile} erzeugten Kandidatenprofile bilden die Kandidatenseite. Artefakt A1 stellt die getrennten Trainings-, Validierungs- und Testanfragen bereit. Artefakt A3 definiert die fachlichen Vergleichsmerkmale und den Aufbau der Rankinggruppen. Die Implementierung überführt diese Vorgaben in die numerische Featurematrix $X$, die Relevanzlabels $\boldsymbol{y}$ und die Gruppenstruktur $\boldsymbol{g}$. Die erzeugten Matrizen bilden die technische Eingabe für das LambdaMART-Training. Die konkrete Zuordnung und Absicherung der Relevanzlabels wird im folgenden Abschnitt behandelt. \subsection{Darstellung der Anfrage-Kandidat-Matrix} \label{subsec:darstellung_anfrage_kandidat_matrix} Abbildung~\ref{fig:technische_anfrage_kandidat_matrix} zeigt die Verknüpfung einer Anfrage mit dem Kandidatenraum. Für jedes Anfrage-Kandidat-Paar wird eine numerische Featurezeile erzeugt. Featurezeilen derselben Anfrage werden zu einer Rankinggruppe zusammengefasst. 
% \begin{figure}[H] \centering \resizebox{\textwidth}{!}{ \begin{tikzpicture}[ >=Stealth, font=\small, query/.style={ draw=blue!65!black, fill=blue!8, rounded corners=2mm, text width=4.5cm, minimum height=1.5cm, align=center, inner sep=7pt }, candidate/.style={ draw=orange!75!black, fill=orange!10, rounded corners=2mm, text width=4.5cm, minimum height=1.5cm, align=center, inner sep=7pt }, process/.style={ draw=green!55!black, fill=green!9, rounded corners=2mm, text width=5.0cm, minimum height=1.5cm, align=center, inner sep=7pt }, feature/.style={ draw=violet!65!black, fill=violet!7, rounded corners=2mm, text width=5.1cm, minimum height=1.45cm, align=center, inner sep=7pt }, matrix/.style={ draw=gray!75!black, fill=gray!9, rounded corners=2mm, text width=5.4cm, minimum height=1.55cm, align=center, inner sep=7pt }, control/.style={ draw=red!65!black, fill=red!4, dashed, rounded corners=2mm, text width=4.7cm, minimum height=1.25cm, align=center, inner sep=6pt }, flow/.style={ ->, thick }, controlflow/.style={ ->, thick, dashed, red!65!black } ] 
%     % Anfrage und Kandidatenraum 
%     \node[query] (query) at (-4.2,6.0) { \textbf{Anfrageartikel}\\ \texttt{query\_macro\_shape}\\ \texttt{query\_article\_vector}\\ \texttt{query\_lexical\_vector} }; \node[candidate] (candidates) at (4.2,6.0) { \textbf{Kandidatenraum}\\ serialisierte Kandidatenprofile aus dem Trainingsanteil }; 
%     % Paarbildung 
%     \node[process] (pairing) at (0,3.8) { \textbf{Paarbildung}\\ Kombination der Anfrage mit jedem zugelassenen Kandidaten }; \draw[flow] (query.south) |- (pairing.west); \draw[flow] (candidates.south) |- (pairing.east); 
%     % Featureberechnung 
%     \node[process] (comparison) at (0,1.5) { \textbf{Paarvergleich}\\ Mengen-, Struktur-, Lexikal-, Zahlen-, Hash- und Kontrollmerkmale }; \draw[flow] (pairing.south) -- (comparison.north); 
%     % Zulässigkeitskontrolle 
%     \node[control] (availability) at (6.2,1.5) { \textbf{Anfrageprüfung}\\ nur ohne Kenntnis des Zielmakros verfügbare Merkmale }; \draw[controlflow] (availability.west) -- (comparison.east); 
%     % Featurezeile 
%     \node[feature] (row) at (0,-0.8) { \textbf{Numerische Featurezeile}\\ $\boldsymbol{x}_{q,c}$ mit fester Feature-Reihenfolge }; \draw[flow] (comparison.south) -- (row.north); 
%     % Label und Gruppe 
%     \node[feature] (label) at (-4.0,-3.0) { \textbf{Relevanzlabel}\\ $y_{q,c}$ }; \node[feature] (group) at (4.0,-3.0) { \textbf{Anfragegruppe}\\ Gruppengröße $g_q$ }; \draw[flow] (row.south) |- (label.north); \draw[flow] (row.south) |- (group.north); 
%     % Matrix 
%     \node[matrix] (matrix) at (0,-5.3) { \textbf{LambdaMART-Datensatz}\\ Featurematrix $X$, Labels $\boldsymbol{y}$ und Gruppenstruktur $\boldsymbol{g}$ }; \draw[flow] (label.south) |- (matrix.west); \draw[flow] (group.south) |- (matrix.east); \end{tikzpicture} } \caption{ Technische Erzeugung der Anfrage-Kandidat-Matrix. Aus jeder Kombination einer Anfrage mit einem Kandidatenprofil entsteht eine numerische Featurezeile. Featurezeilen derselben Anfrage bilden gemeinsam eine Rankinggruppe. } \label{fig:technische_anfrage_kandidat_matrix} \end{figure}


% \section{Implementierung der Relevanzlabels und Leakage-Sicherung} \label{sec:implementierung_der_relevanzlabels_und_leakage_sicherung} Für jedes Anfrage-Kandidat-Paar wird ein Relevanzlabel erzeugt. Das Label beschreibt die fachliche Beziehung zwischen der Anfrage und dem Kandidaten und bildet gemeinsam mit dem Paarmerkmalsvektor die Trainingsgrundlage des LambdaMART-Modells. Parallel dazu entfernt und kontrolliert die Pipeline Merkmale, aus denen das Zielmakro unmittelbar abgeleitet werden könnte. Die Labelbildung und die Leakage-Sicherung werden getrennt ausgeführt: Die Zielzuordnung bestimmt ausschließlich das Trainingslabel und darf nicht in den Paarmerkmalsvektor eingehen. \subsection{Technische Schwerpunkte} \label{subsec:technische_schwerpunkte_labels_leakage} Die Implementierung umfasst: \begin{itemize} \item Labelzuweisung für jedes Anfrage-Kandidat-Paar, \item Unterscheidung zwischen exaktem Zielmakro, gleicher Makrofamilie, gleichem Familienpräfix und nicht relevantem Kandidaten, \item konfigurierbare Umschaltung zwischen binärer und abgestufter Relevanz, \item Kopplung der Relevanzstufen an \texttt{LIGHTGBM\_LABEL\_GAIN}, \item Ausschluss definierter Leakage-Dimensionen, \item Entfernung redundanter und rollenfremder Merkmale, \item Prüfung der Featurebezeichnungen auf direkte Zielhinweise, \item Leave-one-out-Prüfung für Trainingsanfragen, \item Speicherung technischer Leakage- und Sicherheitsberichte. \end{itemize} \subsection{Technische Umsetzung} \label{subsec:technische_umsetzung_labels_leakage} Die Relevanz eines Kandidaten wird durch die Funktion \texttt{relevance\_label()} bestimmt. Hierzu werden das Zielmakro der Anfrage, der bewertete Kandidat sowie deren Familieninformationen miteinander verglichen. Die Implementierung unterscheidet folgende Relevanzstufen: \begin{description} \item[Exaktes Zielmakro] Der Kandidat entspricht dem der Anfrage zugeordneten Zielmakro. \item[Identische Makrofamilie] Kandidat und Zielmakro gehören derselben Makrofamilie an, sind jedoch nicht identisch. \item[Übereinstimmendes Familienpräfix] Kandidat und Zielmakro gehören zu einem übergeordnet verwandten Makrobereich. \item[Keine Relevanz] Zwischen Kandidat und Zielmakro besteht keine der definierten Beziehungen. \end{description} Die Relevanzstufen werden durch die Konstanten \texttt{RELEVANCE\_EXACT\_TARGET}, \texttt{RELEVANCE\_SAME\_FAMILY}, \texttt{RELEVANCE\_SAME\_PREFIX} und \texttt{RELEVANCE\_NONE} abgebildet. Für ein Anfrage-Kandidat-Paar $(q,c)$ ergibt sich das Label damit aus: \begin{equation} y_{q,c} = \begin{cases} r_{\mathrm{exact}}, & c = t_q, \\[2mm] r_{\mathrm{family}}, & \operatorname{family}(c) = \operatorname{family}(t_q), \\[2mm] r_{\mathrm{prefix}}, & \operatorname{prefix}(c) = \operatorname{prefix}(t_q), \\[2mm] r_{\mathrm{none}}, & \text{sonst}. \end{cases} \label{eq:relevanzlabel_anfrage_kandidat} \end{equation} Dabei bezeichnet $t_q$ das Zielmakro der Anfrage und $r_{\mathrm{exact}}$, $r_{\mathrm{family}}$, $r_{\mathrm{prefix}}$ sowie $r_{\mathrm{none}}$ die konfigurierten Relevanzwerte. Über \texttt{EXACT\_TARGET\_ONLY\_TRAINING} kann die abgestufte Labelbildung auf ein binäres Schema reduziert werden. In diesem Fall erhält ausschließlich das exakte Zielmakro ein positives Label. Familien- und Präfixbeziehungen werden wie nicht relevante Kandidaten behandelt. Die Labelwerte werden mit \texttt{LIGHTGBM\_LABEL\_GAIN} abgestimmt. Diese Konfiguration legt fest, welchen Gewinn LightGBM den einzelnen Relevanzstufen zuordnet. Die Position einer Relevanzstufe im Labelschema und der zugehörige Gain-Wert müssen daher konsistent definiert sein. Zusätzlich können schwierige Negativbeispiele stärker berücksichtigt werden. Dazu zählen insbesondere Kandidaten derselben Makrofamilie, die dem Zielmakro fachlich ähnlich sind, aber nicht dem exakten Ziel entsprechen. Diese Kandidaten unterstützen das Modell dabei, nicht nur offensichtlich unpassende, sondern auch eng verwandte Makros voneinander zu unterscheiden. \subsubsection{Ausschluss unzulässiger Merkmale} \label{subsubsec:ausschluss_unzulaessiger_merkmale} Die Zielzuordnung wird ausschließlich zur Bildung des Labels verwendet. Sie darf nicht Bestandteil des numerischen Paarmerkmalsvektors sein. Die in \texttt{LEAKAGE\_DIMENSIONS} festgelegten Dimensionen werden deshalb vor der Matrixerzeugung aus den trainierbaren Repräsentationen entfernt. Zu den auszuschließenden Angaben gehören insbesondere: \begin{itemize} \item Zielmakro und Makroname, \item direkte Artikel-Makro-Zuordnung, \item aus dem Makronamen abgeleitete Zieltoken, \item Artikel- und Dokumentidentitäten, \item Datei- und Dokumentpfade, \item Bearbeiter- und Prozessinformationen, \item Split- und Laufkennungen, \item Evaluations- und Ergebnisinformationen. \end{itemize} Familien- und Zielinformationen dürfen weiterhin für die Labelzuweisung, die Splitprüfung und die Evaluation verwendet werden. Sie sind jedoch von den trainierbaren Modellmerkmalen zu trennen. \subsubsection{Rollen- und Redundanzbereinigung} \label{subsubsec:rollen_und_redundanzbereinigung} Die Rollenbereinigung stellt sicher, dass eine Merkmalsdimension nur in der dafür vorgesehenen Repräsentation verarbeitet wird. Spezialisierte Artikel-, Makrostruktur- und Lexikaldimensionen werden aus allgemeinen Merkmalsstrukturen entfernt, sofern sie dort redundant enthalten sind. Damit wird verhindert, dass ein Merkmal mehrfach in die Paarmerkmalsberechnung eingeht. Die Bereinigung reduziert zugleich das Risiko, dass eine unzulässige Information über eine zweite Repräsentation erhalten bleibt. Nach der Bereinigung gilt für die verwendeten Merkmalsmengen: \begin{equation} F_{\mathrm{model}} = F_{\mathrm{available}} \setminus \left( F_{\mathrm{identity}} \cup F_{\mathrm{process}} \cup F_{\mathrm{target}} \cup F_{\mathrm{redundant}} \right). \label{eq:bereinigte_modellmerkmale} \end{equation} Dabei bezeichnet $F_{\mathrm{model}}$ die freigegebene Modellmerkmalsmenge. \subsubsection{Leave-one-out-Prüfung} \label{subsubsec:leave_one_out_pruefung} Bei Trainingsanfragen kann der Anfrageartikel Bestandteil des Kandidatenprofils seines eigenen Zielmakros sein. Dadurch könnten dessen Merkmale sowohl auf der Anfrage- als auch auf der Kandidatenseite auftreten. Die Funktion \texttt{build\_leave\_one\_out\_profile()} ermöglicht eine diagnostische Prüfung dieses Effekts. Hierbei wird für den jeweiligen Trainingsartikel ein temporäres Kandidatenprofil erzeugt, in dem dessen eigener Beitrag entfernt ist: \begin{equation} P_{t_q}^{(-q)} = \operatorname{aggregate} \left( A_{t_q}^{\mathrm{train}} \setminus \{q\} \right). \label{eq:leave_one_out_profil} \end{equation} Das Leave-one-out-Profil dient der Prüfung, ob die Modellbewertung wesentlich von der Eigenevidenz eines Trainingsartikels abhängt. Bei Kandidaten mit nur einem Trainingsartikel kann nach dessen Entfernung kein vollständiges Profil gebildet werden. Diese Singleton-Fälle werden deshalb gesondert gekennzeichnet. \subsubsection{Technische Leakage-Prüfung} \label{subsubsec:technische_leakage_pruefung} Nach der Merkmalsbereinigung prüft die Pipeline die erzeugten Featurebezeichnungen und Datenstrukturen auf weiterhin vorhandene Zielinformationen. Kontrolliert werden insbesondere: \begin{itemize} \item verbotene Dimensionen nach der Bereinigung, \item direkte Makro- und Labeltoken in Featurebezeichnungen, \item redundante Merkmalsrollen, \item unzulässige Identitäts- und Prozessfelder, \item Herkunft der anfrageseitigen Strukturmerkmale, \item Eigenevidenz bei Trainingsanfragen. \end{itemize} Die Ergebnisse werden in technischen Prüfberichten, unter anderem im \texttt{retrieval\_leakage\_safety\_report.json}, gespeichert. Diese Berichte dokumentieren die ausgeführten Kontrollen. Sie ersetzen nicht die Prüfung, ob jedes verwendete Anfragemerkmal auch im späteren Anwendungsfall verfügbar ist. \subsection{Abgleich zwischen bisheriger Arbeit und Code} \label{subsec:abgleich_labels_leakage} \subsubsection*{Übereinstimmungen} Die bisherige Arbeit und der Code verwenden die bekannte Artikel-Makro-Zuordnung zur Bildung der Relevanzlabels. Gleichzeitig wird diese Zuordnung von den trainierbaren Anfrage- und Kandidatenmerkmalen ausgeschlossen. Die Unterscheidung zwischen exaktem Zielmakro, verwandten Kandidaten und nicht relevanten Kandidaten entspricht der konzeptionellen Relevanzdefinition. Ebenso wird die Trennung von Trainings-, Validierungs- und Testdaten bei der Label- und Matrixerzeugung beibehalten. \subsubsection*{Technische Konkretisierungen} Der Code konkretisiert die fachliche Relevanzdefinition durch \texttt{relevance\_label()} und explizite Relevanzkonstanten. Die Labelstufen werden zusätzlich mit den Gain-Werten von LightGBM verknüpft. Die Leakage-Sicherung ist umfangreicher als der ausschließliche Entzug des Makronamens. Die Implementierung umfasst Dimensionsfilter, Rollenbereinigung, Featureprüfungen, Sicherheitsberichte und eine Leave-one-out-Diagnostik. Der Code unterscheidet außerdem zwischen der technischen Makroidentität eines Kandidaten und dessen trainierbaren Eigenschaften. Die Identität wird zur Verwaltung, Labelbildung und Evaluation benötigt, bleibt jedoch von der Modelleingabe getrennt. \subsubsection*{Abweichungen} Die bisherige Arbeit beschreibt überwiegend ein abgestuftes Relevanzschema. Der Code unterstützt zusätzlich ein binäres Schema über \texttt{EXACT\_TARGET\_ONLY\_TRAINING}. Das Vorhandensein beider Varianten legt nicht fest, welche davon im finalen Trainingslauf verwendet wurde. Die aktive Labelpolitik muss deshalb aus der wirksamen Laufkonfiguration und der Trainingszusammenfassung hervorgehen. Eine nachträgliche Auswahl anhand der Testergebnisse ist unzulässig. Eine weitere Abweichung betrifft die Eigenevidenz in den Trainingsprofilen. Die bisherige Formulierung legt zwar fest, dass Kandidatenprofile ausschließlich aus Trainingsartikeln entstehen, behandelt aber nicht durchgehend den Fall, dass der Anfrageartikel selbst Bestandteil des Zielkandidatenprofils ist. Der Code sieht hierfür eine Leave-one-out-Diagnostik vor. Zudem verwendet die Implementierung Familienkennungen und Familienpräfixe für die Labelbildung. Diese Angaben dürfen nicht gleichzeitig als direkte Paarmerkmale verwendet werden, da das Modell sonst Teile der Labelregel als Eingabe erhalten würde. \subsubsection*{Konsequenz für die Dokumentation} Für das finale Modell sind folgende Angaben eindeutig zu dokumentieren: \begin{itemize} \item aktives binäres oder abgestuftes Labelschema, \item numerischer Wert jeder Relevanzstufe, \item verwendete \texttt{label\_gain}-Konfiguration, \item Behandlung familiennaher negativer Kandidaten, \item Umfang der Leakage-Dimensionen, \item Verwendung der Leave-one-out-Prüfung, \item Behandlung von Singleton-Kandidaten. \end{itemize} Die Labelbildung ist strikt von der Merkmalsbildung zu trennen. Zielmakro, Familie und Familienpräfix dürfen zur Bestimmung des Labels verwendet werden. Als Modelleingabe sind sie nur dann zulässig, wenn daraus keine Ziel- oder Labelinformation hervorgeht. Direkte Identitäts- und Beziehungskennungen bleiben ausgeschlossen. Die Zuordnung zwischen Datenblatt und Artikel darf ausschließlich für Tests, Evaluation und Optimierung verwendet werden. Sie darf weder als Trainingssignal noch zur Konstruktion der Trainingsmerkmale eingesetzt werden. \subsection{Zusammenhang mit der bisherigen Arbeit} \label{subsec:zusammenhang_labels_leakage} Die Relevanzlabels setzen die in Kapitel~6 definierte fachliche Beziehung zwischen Anfrage und Makrokandidat technisch um. Die Artikel-Makro-Zuordnung aus Artefakt A1 stellt hierfür die Referenzinformation bereit. Der in Artefakt A2 ausgewählte LambdaMART-Ansatz benötigt diese Labels gemeinsam mit den Anfragegruppen und Paarmerkmalsvektoren. Die in Abschnitt~\ref{sec:technische_erzeugung_der_anfrage_kandidat_matrix} erzeugten Anfrage-Kandidat-Zeilen werden in diesem Abschnitt um die Relevanzlabels ergänzt und auf unzulässige Merkmale geprüft. Damit entsteht die freigegebene Trainingsgrundlage für das in Abschnitt~\ref{sec:implementierung_des_lambdamart_trainings} beschriebene Modelltraining. Die Leakage-Sicherung konkretisiert die in den vorherigen Kapiteln formulierte Anforderung, dass das Modell ausschließlich aus technischen Vergleichsmerkmalen lernen darf. Identitätsinformationen und bekannte Zielbeziehungen bleiben auf Labelbildung, Verwaltung und Evaluation beschränkt. \subsection{Darstellung der Labelbildung und Leakage-Sicherung} \label{subsec:darstellung_labels_leakage} Abbildung~\ref{fig:relevanzlabels_leakage_sicherung} zeigt die getrennten Verarbeitungspfade für Label- und Merkmalsbildung. Die bekannte Zielzuordnung fließt ausschließlich in die Labelbildung ein. Die technischen Repräsentationen werden unabhängig davon bereinigt und geprüft. Erst danach werden Featurezeile, Label und Anfragegruppe zum LambdaMART-Datensatz zusammengeführt.
% \begin{figure}[H] \centering \resizebox{\textwidth}{!}{ \begin{tikzpicture}[ >=Stealth, font=\small, input/.style={ draw=blue!65!black, fill=blue!8, rounded corners=2mm, text width=4.5cm, minimum height=1.35cm, align=center, inner sep=6pt }, label/.style={ draw=orange!75!black, fill=orange!11, rounded corners=2mm, text width=4.8cm, minimum height=1.4cm, align=center, inner sep=7pt }, process/.style={ draw=green!55!black, fill=green!9, rounded corners=2mm, text width=4.8cm, minimum height=1.4cm, align=center, inner sep=7pt }, control/.style={ draw=red!65!black, fill=red!4, dashed, rounded corners=2mm, text width=4.8cm, minimum height=1.35cm, align=center, inner sep=6pt }, output/.style={ draw=violet!65!black, fill=violet!7, rounded corners=2mm, text width=5.1cm, minimum height=1.45cm, align=center, inner sep=7pt }, dataset/.style={ draw=gray!75!black, fill=gray!9, rounded corners=2mm, text width=5.5cm, minimum height=1.5cm, align=center, inner sep=7pt }, flow/.style={ ->, thick }, controlflow/.style={ ->, thick, dashed, red!65!black } ] 
%     % Getrennte Eingaben 
%     \node[input] (targetRelation) at (-4.2,6.6) { \textbf{Referenzinformation}\\ Zielmakro, Familie und Familienpräfix }; \node[input] (representations) at (4.2,6.6) { \textbf{Technische Repräsentationen}\\ Anfrage- und Kandidatenmerkmale }; 
%     % Labelpfad 
%     \node[label] (labelPolicy) at (-4.2,4.2) { \textbf{Labelbildung}\\ exakt, gleiche Familie, gleiches Präfix oder nicht relevant }; \draw[flow] (targetRelation.south) -- (labelPolicy.north); \node[label] (labelGain) at (-4.2,2.0) { \textbf{Relevanzlabel}\\ $y_{q,c}$ und zugehöriger LightGBM-\texttt{label\_gain} }; \draw[flow] (labelPolicy.south) -- (labelGain.north); 
%     % Merkmalspfad 
%     \node[process] (featureCleaning) at (4.2,4.2) { \textbf{Merkmalsbereinigung}\\ Leakage-Dimensionen, Rollen und Redundanzen entfernen }; \draw[flow] (representations.south) -- (featureCleaning.north); \node[control] (leakageChecks) at (4.2,2.0) { \textbf{Leakage-Prüfung}\\ Zieltoken, Identitäten, Featureherkunft und Eigenevidenz }; \draw[controlflow] (featureCleaning.south) -- (leakageChecks.north); 
%     % Trennung der beiden Pfade 
%     \node[control] (separation) at (0,4.2) { \textbf{Strikte Trennung}\\ Referenzinformationen werden nicht zu Modellmerkmalen. }; \draw[controlflow] (targetRelation.east) |- (separation.north); \draw[controlflow] (separation.south) |- (featureCleaning.west); 
%     % Freigegebene Featurezeile 
%     \node[output] (approvedFeatures) at (4.2,-0.4) { \textbf{Freigegebene Featurezeile}\\ $\boldsymbol{x}_{q,c}$ ohne Ziel- und Identitätsinformationen }; \draw[flow] (leakageChecks.south) -- (approvedFeatures.north); 
%     % Zusammenführung 
%     \node[dataset] (trainingRow) at (0,-3.0) { \textbf{Freigegebene Trainingszeile}\\ Featurezeile $\boldsymbol{x}_{q,c}$, Relevanzlabel $y_{q,c}$ und Anfragegruppe }; \draw[flow] (labelGain.south) |- (trainingRow.west); \draw[flow] (approvedFeatures.south) |- (trainingRow.east); 
%     % Bericht 
%     \node[control] (report) at (6.6,-3.0) { \textbf{Prüfbericht}\\ Leakage-Sicherheitsbericht und Leave-one-out-Diagnostik }; \draw[controlflow] (leakageChecks.east) -| (report.north); \end{tikzpicture} } \caption{ Getrennte Bildung der Relevanzlabels und Modellmerkmale. Ziel- und Familieninformationen bestimmen ausschließlich das Label. Die Paarmerkmale werden unabhängig davon bereinigt und auf unzulässige Zielinformationen geprüft. } \label{fig:relevanzlabels_leakage_sicherung} \end{figure}


% \section{Implementierung des LambdaMART-Trainings} \label{sec:implementierung_des_lambdamart_trainings} Das Modelltraining verarbeitet die zuvor erzeugten Anfrage-Kandidat-Matrizen. Für jede Anfrage wird eine Rankinggruppe übergeben, innerhalb derer LambdaMART die Kandidaten anhand ihrer Relevanzlabels und Paarmerkmale ordnet. Die Implementierung nutzt den LightGBM-Ranker mit dem Lernziel \texttt{lambdarank}. Mehrere Hyperparameterkonfigurationen werden auf dem Trainingsanteil gelernt und ausschließlich anhand des Validierungsanteils miteinander verglichen. Der Testanteil wird erst nach Festlegung der Modellkonfiguration ausgewertet. \subsection{Technische Schwerpunkte} \label{subsec:technische_schwerpunkte_lambdamart_training} Die Trainingsimplementierung umfasst: \begin{itemize} \item Laden der Trainings- und Validierungsmatrizen, \item Übergabe der Featurematrix, Relevanzlabels und Anfragegruppen, \item Initialisierung des LightGBM-Rankers, \item Verwendung von \texttt{lambdarank} als Lernziel, \item Übergabe der konfigurierten Relevanzgewinne, \item Training mehrerer Hyperparameterkonfigurationen, \item Bewertung der Konfigurationen auf dem Validierungsanteil, \item Auswahl der besten Modellkonfiguration, \item einmalige Bewertung auf dem Testanteil, \item Serialisierung des Modells und der Trainingsmetadaten. \end{itemize} \subsection{Technische Umsetzung} \label{subsec:technische_umsetzung_lambdamart_training} Für jeden Datensplit liegen die Featurematrix $X$, der Labelvektor $\boldsymbol{y}$ und die Gruppenstruktur $\boldsymbol{g}$ vor. Die Gruppenstruktur enthält die Anzahl der Kandidaten je Anfrage und stellt sicher, dass LambdaMART ausschließlich Kandidaten derselben Anfrage miteinander vergleicht. Der Trainingsdatensatz wird formal durch \begin{equation} D_{\mathrm{train}} = \left( X_{\mathrm{train}}, \boldsymbol{y}_{\mathrm{train}}, \boldsymbol{g}_{\mathrm{train}} \right) \label{eq:lambdamart_trainingsdaten} \end{equation} beschrieben. Der Validierungs- und Testdatensatz besitzen dieselbe Struktur, enthalten jedoch ausschließlich die Anfragen des jeweiligen Splits. \subsubsection{Initialisierung des Rankingmodells} \label{subsubsec:initialisierung_rankingmodell} Das Rankingmodell wird als LightGBM-Ranker initialisiert. Als Lernziel wird \texttt{lambdarank} verwendet. Das Modell erzeugt einen Score \begin{equation} s_{q,c} = f_{\theta} \left( \boldsymbol{x}_{q,c} \right) \label{eq:lambdamart_score} \end{equation} für jedes Anfrage-Kandidat-Paar. Dabei bezeichnet $\boldsymbol{x}_{q,c}$ den Paarmerkmalsvektor und $\theta$ die gelernten Modellparameter. Die vorhergesagten Scores werden innerhalb jeder Anfragegruppe absteigend sortiert. Der Score besitzt keine eigenständige fachliche Einheit. Er dient ausschließlich zur Bestimmung der Rangfolge der Kandidaten. Die Relevanzstufen werden über \texttt{LIGHTGBM\_LABEL\_GAIN} mit den für das Training verwendeten Gewinnen verknüpft. Die Gain-Konfiguration muss mit der in Abschnitt~\ref{sec:implementierung_der_relevanzlabels_und_leakage_sicherung} festgelegten Labelpolitik übereinstimmen. \subsubsection{Hyperparameterkonfiguration} \label{subsubsec:hyperparameterkonfiguration_lambdamart} Die Implementierung trainiert mehrere vorher festgelegte Hyperparameterkombinationen. Variiert werden insbesondere: \begin{description} \item[\texttt{n\_estimators}] Anzahl der nacheinander erzeugten Entscheidungsbäume, \item[\texttt{learning\_rate}] Gewicht des Beitrags eines neu hinzugefügten Baums, \item[\texttt{num\_leaves}] maximale Anzahl der Blätter eines Entscheidungsbaums, \item[\texttt{min\_child\_samples}] Mindestanzahl von Trainingsbeispielen in einem Blatt, \item[\texttt{reg\_alpha}] Stärke der L1-Regularisierung, \item[\texttt{reg\_lambda}] Stärke der L2-Regularisierung, \item[\texttt{min\_split\_gain}] erforderlicher Mindestgewinn für eine weitere Baumaufteilung, \item[\texttt{lambdarank\_truncation\_level}] Rangposition, bis zu der die LambdaMART-Optimierung besonders berücksichtigt wird. \end{description} Die Regularisierung begrenzt die Modellkomplexität und reduziert das Risiko einer zu starken Anpassung an die Trainingsanfragen. Das Truncation-Level richtet die Optimierung auf den für die Makroempfehlung relevanten oberen Bereich der Rangliste aus. Die getesteten Parameterkombinationen werden vor der abschließenden Testauswertung festgelegt. Eine Erweiterung oder Anpassung des Suchraums anhand der Testmetriken ist nicht zulässig. \subsubsection{Training und Validierung} \label{subsubsec:training_und_validierung_lambdamart} Jede Hyperparameterkonfiguration wird ausschließlich auf $D_{\mathrm{train}}$ trainiert. Anschließend erzeugt das Modell für die Anfragen des Validierungsanteils Kandidatenscores. Die Kandidaten werden innerhalb ihrer Anfragegruppe nach diesen Scores sortiert. Die Bewertung einer Konfiguration erfolgt anhand der in der Trainingspipeline festgelegten Validierungsmetriken. Dazu gehören insbesondere: \begin{itemize} \item Recall@1, \item Recall@10, \item NDCG@10. \end{itemize} Die Modellkonfiguration wird anhand einer festgelegten Auswahlregel bestimmt. Abhängig von der aktiven Implementierung kann hierfür eine einzelne Zielmetrik oder ein vorab definierter Gesamtscore verwendet werden. Die Auswahlregel muss vor Einsicht in die Testmetriken feststehen und gemeinsam mit der Trainingskonfiguration gespeichert werden. Für eine Menge getesteter Konfigurationen $\Theta=\{\theta_1,\ldots,\theta_n\}$ ergibt sich die Auswahl formal durch \begin{equation} \theta^{*} = \underset{\theta \in \Theta}{\operatorname{arg\,max}} \; M_{\mathrm{valid}} \left( f_{\theta} \right), \label{eq:auswahl_lambdamart_modell} \end{equation} wobei $M_{\mathrm{valid}}$ die festgelegte Validierungs-Auswahlfunktion bezeichnet. \subsubsection{Abschließende Testbewertung} \label{subsubsec:abschliessende_testbewertung_lambdamart} Nach der Auswahl von $\theta^{*}$ wird die festgelegte Modellkonfiguration auf dem unabhängigen Testanteil bewertet. Der Testanteil wird nicht zur Auswahl von Features, Relevanzlabels, Profilgrenzen, Hyperparametern oder Auswahlmetriken verwendet. Die Testvorhersagen werden getrennt von den Validierungsvorhersagen gespeichert. Dadurch bleibt nachvollziehbar, auf welcher Datengrundlage die Modellauswahl und die abschließende Bewertung durchgeführt wurden. Die technische Berechnung der Testmetriken erfolgt innerhalb der Pipeline. Ihre wissenschaftliche Interpretation und der Vergleich mit den definierten Zielwerten sind Gegenstand von Kapitel~8. \subsubsection{Speicherung des Trainingsergebnisses} \label{subsubsec:speicherung_trainingsergebnis} Nach Abschluss des Trainings werden das ausgewählte Modell und die zugehörigen Laufdaten gespeichert. Dazu gehören insbesondere: \begin{itemize} \item die serialisierte LambdaMART-Modelldatei, \item die ausgewählten Hyperparameter, \item die verwendete Label- und Gain-Konfiguration, \item die verbindliche Feature-Reihenfolge, \item die Validierungsmetriken der getesteten Konfigurationen, \item die Testvorhersagen des ausgewählten Modells, \item die Kennung und Konfiguration des Trainingslaufs. \end{itemize} Die ausgewählten Parameter werden unter anderem in \texttt{best\_lightgbm\_params} und der \texttt{training\_summary.json} dokumentiert. Die Modelldatei wird gemeinsam mit den Kandidatenprofilen und Featuredefinitionen dem Artefakt A4 zugeordnet. \subsection{Abgleich zwischen bisheriger Arbeit und Code} \label{subsec:abgleich_lambdamart_training} \subsubsection*{Übereinstimmungen} Die bisherige Arbeit und der Code verwenden LambdaMART zur Bewertung und Sortierung der Makrokandidaten. Das Modell wird mit numerischen Merkmalen, Relevanzlabels und Anfragegruppen trainiert. Die Datensplits erfüllen unterschiedliche Aufgaben: \begin{itemize} \item Der Trainingsanteil dient zum Lernen der Modellparameter. \item Der Validierungsanteil dient zur Auswahl der Hyperparameterkonfiguration. \item Der Testanteil dient zur abschließenden, unabhängigen Bewertung. \end{itemize} Diese Trennung entspricht der in den vorherigen Kapiteln festgelegten Trainings- und Evaluationslogik. \subsubsection*{Technische Konkretisierungen} Der Code konkretisiert den Modellansatz durch den LightGBM-Ranker, das Lernziel \texttt{lambdarank} und einen definierten Hyperparametersuchraum. Die Gruppengrößen werden explizit an LightGBM übergeben. Damit wird die im bisherigen Text fachlich beschriebene Anfragegruppe technisch als zusammenhängender Bereich der Featurematrix abgebildet. Der Parameter \texttt{lambdarank\_truncation\_level} richtet das Training auf die vorderen Rangpositionen aus. Dies konkretisiert die Anforderung, insbesondere die Qualität der Top-Empfehlungen zu verbessern. Die Implementierung speichert neben dem trainierten Modell auch die ausgewählten Parameter, Featurebezeichnungen, Vorhersagen und Metriken. Dadurch bleibt die Modellauswahl nachvollziehbar. \subsubsection*{Abweichungen} Die bisherige Arbeit beschreibt Artikel-, Makro- und Lexikalvektoren teilweise als unmittelbare Eingaben des LambdaMART-Modells. Der Code übergibt dagegen ausschließlich die in Abschnitt~\ref{sec:technische_erzeugung_der_anfrage_kandidat_matrix} erzeugten Paarmerkmalsvektoren. Damit gilt: \begin{equation} f_{\theta} \left( \boldsymbol{x}_{q,c} \right) \neq f_{\theta} \left( q_{\mathrm{article}}, c_{\mathrm{profile}} \right). \label{eq:lambdamart_verarbeitet_paarmerkmale} \end{equation} LambdaMART erhält keine fachlichen Objekte, sondern ausschließlich eine numerische Featurematrix mit fester Spaltenreihenfolge. Eine weitere Abweichung besteht in den vorhandenen Modellkonfigurationen. Einzelne Konfigurationsdateien enthalten Parameter früherer oder alternativer neuronaler Ansätze, etwa Embedding-Dimensionen, versteckte Schichten, Dropout oder Temperaturparameter. Diese Einstellungen sind keine LambdaMART-Parameter und dürfen nicht als Bestandteil des finalen Rankingmodells beschrieben werden, sofern sie im ausgeführten Trainingslauf nicht verwendet werden. Der Code kann Validierungs- und Testmetriken innerhalb desselben Pipeline-Laufs berechnen. Methodisch entscheidend ist nicht der Zeitpunkt der technischen Berechnung, sondern ihre Verwendung. Die Testmetriken dürfen keine Rückwirkung auf Hyperparametersuche, Featureauswahl, Labelpolitik oder Profilbildung haben. Die Implementierung kann außerdem mehrere Modellkennzahlen zu einem Auswahlwert zusammenführen. Falls im bisherigen Text eine einzelne Metrik als alleinige Auswahlgrundlage genannt wird, muss die Darstellung an die tatsächlich aktive Auswahlfunktion angepasst werden. \subsubsection*{Konsequenz für die Dokumentation} Für den finalen Trainingslauf sind mindestens folgende Angaben eindeutig zu dokumentieren: \begin{itemize} \item verwendete LightGBM-Version, \item aktives Lernziel, \item vollständiger Hyperparametersuchraum, \item ausgewählte Hyperparameter, \item verwendete Label- und Gain-Konfiguration, \item verwendetes Truncation-Level, \item Seed und weitere Zufallseinstellungen, \item Validierungs-Auswahlregel, \item Zeitpunkt und Umfang der Testauswertung. \end{itemize} Die Darstellung des Trainings muss klar zwischen drei Ebenen unterscheiden: \begin{description} \item[Modellparameter] werden während des Trainings aus dem Trainingsdatensatz gelernt. \item[Hyperparameter] werden anhand des Validierungsdatensatzes ausgewählt. \item[Testmetriken] werden erst nach Festlegung der Modell- und Hyperparameterkonfiguration zur abschließenden Bewertung verwendet. \end{description} Parameter neuronaler oder nicht verwendeter Legacy-Modelle werden nicht dem Artefakt A4 zugerechnet. Maßgeblich sind ausschließlich die im finalen LambdaMART-Lauf wirksamen Einstellungen. \subsection{Zusammenhang mit der bisherigen Arbeit} \label{subsec:zusammenhang_lambdamart_training} Artefakt A1 stellt die getrennten Trainings-, Validierungs- und Testdaten bereit. Artefakt A2 legt LambdaMART als Rankingverfahren fest. Artefakt A3 definiert die Rankinggruppen, Vergleichsmerkmale und den Ablauf der Kandidatenbewertung. Die in Abschnitt~\ref{sec:technische_erzeugung_der_anfrage_kandidat_matrix} erzeugten Featurematrizen bilden die numerische Modelleingabe. Die in Abschnitt~\ref{sec:implementierung_der_relevanzlabels_und_leakage_sicherung} beschriebenen Labels und Leakage-Prüfungen stellen die zulässige Trainingsgrundlage sicher. Das ausgewählte LambdaMART-Modell bildet zusammen mit den Kandidatenprofilen, Featuredefinitionen und Laufmetadaten den Kern des Artefakts A4. Die vollständige Zusammensetzung dieses Artefakts wird im folgenden Abschnitt beschrieben. \subsection{Darstellung des LambdaMART-Trainings} \label{subsec:darstellung_lambdamart_training} Abbildung~\ref{fig:technisches_lambdamart_training} zeigt die Trennung zwischen Modelltraining, Hyperparameterauswahl und abschließender Testbewertung. Der Testanteil besitzt keinen Rückführungspfad zur Modellauswahl. 
% \begin{figure}[H] \centering \resizebox{\textwidth}{!}{ \begin{tikzpicture}[ >=Stealth, font=\small, data/.style={ draw=blue!65!black, fill=blue!8, rounded corners=2mm, text width=4.2cm, minimum height=1.45cm, align=center, inner sep=7pt }, config/.style={ draw=gray!70!black, fill=gray!8, rounded corners=2mm, text width=4.4cm, minimum height=1.4cm, align=center, inner sep=7pt }, training/.style={ draw=green!55!black, fill=green!9, rounded corners=2mm, text width=4.7cm, minimum height=1.5cm, align=center, inner sep=7pt }, validation/.style={ draw=orange!75!black, fill=orange!11, rounded corners=2mm, text width=4.8cm, minimum height=1.5cm, align=center, inner sep=7pt }, test/.style={ draw=violet!65!black, fill=violet!7, rounded corners=2mm, text width=4.8cm, minimum height=1.5cm, align=center, inner sep=7pt }, output/.style={ draw=red!65!black, fill=red!5, rounded corners=2mm, text width=5.0cm, minimum height=1.45cm, align=center, inner sep=7pt }, flow/.style={ ->, thick }, evalflow/.style={ ->, thick, dashed, violet!65!black } ]
%     % Eingabedaten 
%     \node[data] (trainData) at (-5.2,6.3) { \textbf{Trainingsdaten}\\ $X_{\mathrm{train}}$, $\boldsymbol{y}_{\mathrm{train}}$ und $\boldsymbol{g}_{\mathrm{train}}$ }; \node[data] (validData) at (0,6.3) { \textbf{Validierungsdaten}\\ $X_{\mathrm{valid}}$, $\boldsymbol{y}_{\mathrm{valid}}$ und $\boldsymbol{g}_{\mathrm{valid}}$ }; \node[data] (testData) at (5.2,6.3) { \textbf{Testdaten}\\ $X_{\mathrm{test}}$, $\boldsymbol{y}_{\mathrm{test}}$ und $\boldsymbol{g}_{\mathrm{test}}$ }; 
%     % Konfiguration 
%     \node[config] (searchSpace) at (-5.2,4.0) { \textbf{Hyperparametersuchraum}\\ Bäume, Lernrate, Blätter, Mindestblattgröße und Regularisierung }; \draw[flow] (trainData.south) -- (searchSpace.north); 
%     % Training 
%     \node[training] (models) at (-2.8,1.7) { \textbf{LambdaMART-Training}\\ ein LightGBM-Modell je Hyperparameterkonfiguration }; \draw[flow] (searchSpace.south) |- (models.west); 
%     % Validierung 
%     \node[validation] (validation) at (2.8,1.7) { \textbf{Validierungsvergleich}\\ Bewertung aller Konfigurationen anhand der festgelegten Auswahlregel }; \draw[flow] (models.east) -- (validation.west); \draw[flow] (validData.south) -- (validation.north); 
%     % Auswahl 
%     \node[output] (selectedModel) at (0,-0.8) { \textbf{Ausgewählte Modellkonfiguration}\\ $\theta^{*}$ mit dokumentierten Hyperparametern und Validierungsmetriken }; \draw[flow] (validation.south) |- (selectedModel.east); 
%     % Test 
%     \node[test] (testEvaluation) at (5.2,-0.8) { \textbf{Einmalige Testbewertung}\\ Vorhersagen und Metriken ohne Rückwirkung auf die Modellauswahl }; \draw[evalflow] (testData.south) -- (testEvaluation.north); \draw[evalflow] (selectedModel.east) -- (testEvaluation.west); 
%     % Persistierung 
%     \node[output] (artifacts) at (0,-3.3) { \textbf{Trainingsergebnis}\\ Modelldatei, \texttt{best\_lightgbm\_params}, Feature-Reihenfolge und \texttt{training\_summary.json} }; \draw[flow] (selectedModel.south) -- (artifacts.north); \draw[evalflow] (testEvaluation.south) |- (artifacts.east); \end{tikzpicture} } \caption{ Technische Durchführung des LambdaMART-Trainings. Modellparameter werden aus den Trainingsdaten gelernt, Hyperparameter anhand der Validierungsdaten ausgewählt und erst danach einmalig auf dem Testanteil bewertet. } \label{fig:technisches_lambdamart_training} \end{figure}

% \section{Erzeugung des Artefakts A4} \label{sec:erzeugung_des_artefakts_a4} Das Ergebnis der Modellimplementierung ist das Artefakt A4. Es umfasst neben dem trainierten LambdaMART-Modell alle Dateien, die für eine konsistente Modellanwendung und die Reproduktion des Trainingslaufs erforderlich sind. A4 wird in ein Runtime-Bundle und ein Nachweispaket unterteilt. Das Runtime-Bundle enthält die für die Kandidatenbewertung benötigten Bestandteile. Das Nachweispaket dokumentiert die Erzeugung, Konfiguration und Bewertung des Modells. \subsection{Technische Schwerpunkte} \label{subsec:technische_schwerpunkte_artefakt_a4} Die Erzeugung von A4 umfasst: \begin{itemize} \item Serialisierung des trainierten LambdaMART-Modells, \item Speicherung der Kandidatenprofile, \item Speicherung der Featurebezeichnungen und Feature-Reihenfolge, \item Dokumentation der ausgewählten Modellparameter, \item Speicherung der aktiven Laufkonfiguration, \item Erstellung eines Run-Manifests, \item Speicherung der Trainingszusammenfassung, \item Ablage der Validierungs- und Testvorhersagen, \item Speicherung der Metriken und technischen Prüfberichte, \item Zuordnung aller Dateien zu einem eindeutigen Trainingslauf. \end{itemize} \subsection{Technische Umsetzung} \label{subsec:technische_umsetzung_artefakt_a4} Die im Training erzeugten Dateien werden in einem laufbezogenen Ergebnisverzeichnis gespeichert. Dieses Verzeichnis fasst die Modelldatei, die Kandidatenprofile, die Featuredefinitionen sowie die zugehörigen Konfigurations- und Nachweisdateien zusammen. Die eindeutige Zuordnung zu einem Trainingslauf verhindert, dass Dateien unterschiedlicher Modellversionen miteinander kombiniert werden. Insbesondere das Modell, die Kandidatenprofile und die Feature-Reihenfolge müssen aus demselben Lauf stammen. \subsubsection{Runtime-Bundle} \label{subsubsec:runtime_bundle_a4} Das Runtime-Bundle enthält die Bestandteile, die zur Bewertung einer neuen Anfrage erforderlich sind: \begin{description} \item[LambdaMART-Modell] enthält die während des Trainings gelernten Baumstrukturen und Modellparameter. \item[\texttt{candidate\_profiles.json}] enthält die aus dem Trainingsanteil erzeugten Kandidatenprofile. \item[\texttt{feature\_names.json}] enthält die verbindlichen Bezeichnungen und die Reihenfolge der numerischen Eingabemerkmale. \item[Modellkonfiguration] enthält die für die Modellanwendung relevanten LambdaMART-Parameter. \item[Featurekonfiguration] legt fest, welche Vergleichsmerkmale erzeugt und in welcher Form diese verarbeitet werden. \item[Laufzeitkonfiguration] enthält die für das Laden und Ausführen des Modells erforderlichen Pfad- und Systemeinstellungen. \end{description} Die Feature-Reihenfolge wird gemeinsam mit dem Modell gespeichert, da LightGBM die Eingabespalten anhand ihrer Position verarbeitet. Für die Laufzeitanwendung muss deshalb gelten: \begin{equation} F_{\mathrm{runtime}} = F_{\mathrm{training}}, \label{eq:featuregleichheit_runtime_training} \end{equation} wobei $F_{\mathrm{training}}$ die beim Training und $F_{\mathrm{runtime}}$ die bei der Anwendung verwendete geordnete Featureliste bezeichnet. Neben der Anzahl müssen auch Bedeutung, Datentyp und Position der Merkmale übereinstimmen. Eine abweichende Feature-Reihenfolge kann zu formal gültigen, jedoch inhaltlich fehlerhaften Vorhersagen führen. Die Kandidatenprofile bilden den zur Laufzeit bewerteten Kandidatenraum. Sie werden gemeinsam mit dem Modell geladen und in die in Abschnitt~\ref{sec:technische_bildung_der_kandidatenprofile} beschriebenen Laufzeit-Caches überführt. Ein Runtime-Bundle ist daher nur gültig, wenn mindestens folgende Konsistenzbedingungen erfüllt sind: \begin{equation} B_{\mathrm{runtime}} = \left( M, P, F, C \right), \label{eq:runtime_bundle} \end{equation} wobei $M$ das Modell, $P$ die Kandidatenprofile, $F$ die geordnete Featuredefinition und $C$ die aktive Laufzeitkonfiguration bezeichnet. \subsubsection{Nachweispaket} \label{subsubsec:nachweispaket_a4} Das Nachweispaket enthält die zur Reproduktion und Bewertung des Trainingslaufs erforderlichen Dateien. Dazu gehören insbesondere: \begin{description} \item[\texttt{run\_manifest.json}] dokumentiert die Kennung, Konfiguration und verwendeten Datenstände des Trainingslaufs. \item[\texttt{training\_summary.json}] enthält die ausgewählte Modellkonfiguration, die Trainingsstatistiken und die zusammengefassten Validierungs- und Testergebnisse. \item[\texttt{performance\_metrics.json}] enthält die technisch berechneten Rankingmetriken. \item[Validierungsvorhersagen] enthalten die Scores und Rangpositionen des Validierungsanteils. \item[Testvorhersagen] enthalten die Scores und Rangpositionen des unabhängigen Testanteils. \item[Prüfberichte] dokumentieren Datenqualität, Split-Konsistenz, Featureverwendung und Leakage-Kontrollen. \end{description} Die Vorhersage- und Metrikdateien sind nicht für die Bewertung neuer Anfragen erforderlich. Sie dokumentieren den Trainingslauf und bilden die Datengrundlage für die Evaluation in Kapitel~8. \subsubsection{Konsistenzprüfung der Artefakte} \label{subsubsec:konsistenzpruefung_artefakt_a4} Vor der Verwendung von A4 müssen die Bestandteile des Runtime-Bundles auf Konsistenz geprüft werden. Die Prüfung umfasst: \begin{itemize} \item Übereinstimmung der Run-Kennung, \item Vorhandensein der Modelldatei, \item Vorhandensein der Kandidatenprofile, \item Übereinstimmung der Featureanzahl, \item Übereinstimmung der Feature-Reihenfolge, \item Kompatibilität der Modell- und Featurekonfiguration, \item Übereinstimmung des Daten- und Modellstands, \item Ausschluss unvollständiger oder temporärer Dateien. \end{itemize} Ein Modell darf nicht mit Kandidatenprofilen oder Featuredefinitionen eines anderen Trainingslaufs kombiniert werden. Andernfalls ist nicht sichergestellt, dass die Kandidaten mit denselben Merkmalsdimensionen und derselben Datenbasis verarbeitet werden wie während des Trainings. Die Konsistenzprüfung kann formal durch \begin{equation} \operatorname{valid}(A4) = \operatorname{sameRun}(M,P,F,C) \land \operatorname{sameFeatures}(M,F) \land \operatorname{complete}(M,P,F,C) \label{eq:konsistenzpruefung_a4} \end{equation} beschrieben werden. Nur ein vollständig erzeugtes und geprüftes Bundle wird für die Modellanwendung freigegeben. \subsubsection{Laufzeitverwendung} \label{subsubsec:laufzeitverwendung_artefakt_a4} Zur Laufzeit werden das Modell, die Kandidatenprofile und die Featuredefinition geladen. Für eine neue Anfrage werden dieselben Verarbeitungsschritte wie bei der Trainingsmatrix ausgeführt: \begin{enumerate} \item Erzeugung der zulässigen Anfragerepräsentationen, \item Vergleich mit jedem Kandidatenprofil, \item Bildung der Paarmerkmalsvektoren, \item Anordnung der Merkmale entsprechend \texttt{feature\_names.json}, \item Berechnung der LambdaMART-Scores, \item Sortierung der Kandidaten nach dem Modellscore. \end{enumerate} Die Laufzeitverarbeitung verwendet keine bekannten Relevanzlabels. Diese stehen nur während Training und Evaluation zur Verfügung. Ebenso werden Validierungs- und Testvorhersagen nicht als Eingabe der Modellanwendung verwendet. \subsection{Abgleich zwischen bisheriger Arbeit und Code} \label{subsec:abgleich_artefakt_a4} \subsubsection*{Übereinstimmungen} Die bisherige Arbeit und der Code betrachten das trainierte Empfehlungssystem als Ergebnis der Implementierung. Das Modell bewertet die Makrokandidaten anhand der in Kapitel~6 festgelegten Anfrage-Kandidat-Merkmale. Die Kandidatenprofile werden aus dem Trainingsanteil erzeugt und gemeinsam mit dem Modell für die spätere Kandidatenbewertung bereitgestellt. Die getrennte Speicherung der Trainings-, Validierungs- und Testergebnisse unterstützt die vorgesehene Evaluation. \subsubsection*{Technische Konkretisierungen} Der Code konkretisiert A4 durch mehrere voneinander abhängige Dateien. Die Modelldatei bildet nur den gelernten Rankingalgorithmus ab. Die vollständige Modellanwendung benötigt zusätzlich die Kandidatenprofile, die Featuredefinition und die aktive Laufzeitkonfiguration. Das Run-Manifest und die Trainingszusammenfassung stellen die Verbindung zwischen diesen Dateien her. Dadurch kann nachvollzogen werden, aus welchem Datenstand und welcher Konfiguration eine Modellversion entstanden ist. Die gespeicherten Validierungs- und Testvorhersagen ermöglichen eine nachträgliche Evaluation, ohne das Modelltraining erneut auszuführen. Sie gehören zum Nachweispaket, nicht zur produktiven Modelleingabe. \subsubsection*{Abweichungen} Die bisherige Formulierung kann den Eindruck vermitteln, A4 bestehe ausschließlich aus dem trainierten LambdaMART-Modell. Dies entspricht nicht der technischen Umsetzung. Ohne Kandidatenprofile und Featuredefinition kann das Modell keine neue Anfrage konsistent bewerten. Weiterhin muss zwischen Trainingsartefakten und Laufzeitartefakten unterschieden werden. Validierungsmetriken, Testvorhersagen und Leakage-Berichte dokumentieren den Trainingslauf, werden jedoch für die produktive Scoreberechnung nicht benötigt. Der Code erzeugt laufbezogene Ergebnisdateien. Ein vollständig atomarer Austausch eines produktiven Runtime-Bundles ist davon zu unterscheiden. Das erfolgreiche Schreiben einzelner Dateien garantiert noch nicht, dass zu jedem Zeitpunkt ausschließlich ein vollständiges und konsistentes Bundle geladen wird. Die bisherige Arbeit sollte A4 daher nicht nur als \textit{trainiertes Modell}, sondern als versionierte Einheit aus Modell, Kandidatenraum, Featuredefinition und Konfiguration beschreiben. \subsubsection*{Konsequenz für die Dokumentation} Artefakt A4 wird in zwei Bestandteile gegliedert: \begin{enumerate} \item das Runtime-Bundle für die Modellanwendung, \item das Nachweispaket für Reproduzierbarkeit und Evaluation. \end{enumerate} Das minimale Runtime-Bundle umfasst: \begin{equation} A4_{\mathrm{runtime}} = \left\{ \text{Modell}, \text{Kandidatenprofile}, \text{Featuredefinition}, \text{Laufzeitkonfiguration} \right\}. \label{eq:a4_runtime} \end{equation} Das vollständige Nachweispaket umfasst zusätzlich: \begin{equation} A4_{\mathrm{evidence}} = \left\{ \text{Manifest}, \text{Trainingszusammenfassung}, \text{Metriken}, \text{Vorhersagen}, \text{Prüfberichte} \right\}. \label{eq:a4_evidence} \end{equation} Nicht aktive Legacy-Konfigurationen und alternative Modellkomponenten sind nicht Bestandteil des finalen A4. Maßgeblich sind ausschließlich die Dateien und Parameter des ausgewählten LambdaMART-Laufs. Für einen produktiven Modellwechsel ist ein versionierter und atomarer Bundle-Austausch vorzusehen. Ein neues Bundle darf erst aktiviert werden, wenn sämtliche Bestandteile vollständig geschrieben und erfolgreich geprüft wurden. \subsection{Zusammenhang mit der bisherigen Arbeit} \label{subsec:zusammenhang_artefakt_a4} Artefakt A4 führt die Ergebnisse der vorherigen Abschnitte zusammen. Die in Abschnitt~\ref{sec:technische_umsetzung_der_datenrepraesentationen} beschriebenen Datenrepräsentationen bilden die Grundlage der Kandidaten- und Anfrageverarbeitung. Die in Abschnitt~\ref{sec:technische_bildung_der_kandidatenprofile} erzeugten Kandidatenprofile bilden den gespeicherten Kandidatenraum. Die in Abschnitt~\ref{sec:technische_erzeugung_der_anfrage_kandidat_matrix} festgelegte Featuredefinition bestimmt die Modelleingabe. Die Label- und Leakage-Regeln aus Abschnitt~\ref{sec:implementierung_der_relevanzlabels_und_leakage_sicherung} sichern die Trainingsgrundlage ab. Das in Abschnitt~\ref{sec:implementierung_des_lambdamart_trainings} ausgewählte LambdaMART-Modell bildet gemeinsam mit diesen Bestandteilen das Runtime-Bundle. Manifest, Metriken, Vorhersagen und Prüfberichte ergänzen dieses um das reproduzierbare Nachweispaket. A4 stellt damit das technische Ergebnis von Kapitel~7 und die Eingangsgrundlage für die Evaluation in Kapitel~8 dar. \subsection{Darstellung des Artefakts A4} \label{subsec:darstellung_artefakt_a4} Abbildung~\ref{fig:technische_erzeugung_artefakt_a4} zeigt die Zusammenführung der Modellergebnisse zum Artefakt A4. Runtime-Bundle und Nachweispaket bleiben getrennt, werden jedoch über dieselbe Run-Kennung und Konfiguration verbunden. 
% \begin{figure}[H] \centering \resizebox{\textwidth}{!}{ \begin{tikzpicture}[ >=Stealth, font=\small, input/.style={ draw=blue!65!black, fill=blue!8, rounded corners=2mm, text width=4.2cm, minimum height=1.4cm, align=center, inner sep=7pt }, runtime/.style={ draw=green!55!black, fill=green!9, rounded corners=2mm, text width=5.1cm, minimum height=1.55cm, align=center, inner sep=7pt }, evidence/.style={ draw=gray!70!black, fill=gray!8, rounded corners=2mm, text width=5.1cm, minimum height=1.55cm, align=center, inner sep=7pt }, control/.style={ draw=red!65!black, fill=red!4, dashed, rounded corners=2mm, text width=5.0cm, minimum height=1.35cm, align=center, inner sep=7pt }, artifact/.style={ draw=orange!75!black, fill=orange!12, rounded corners=2mm, text width=5.4cm, minimum height=1.55cm, align=center, inner sep=7pt }, next/.style={ draw=violet!65!black, fill=violet!7, rounded corners=2mm, text width=4.9cm, minimum height=1.35cm, align=center, inner sep=7pt }, flow/.style={ ->, thick }, controlflow/.style={ ->, thick, dashed, red!65!black }, nextflow/.style={ ->, thick, dashed, violet!65!black } ] 
%     % Eingangsartefakte 
%     \node[input] (model) at (-5.2,6.3) { \textbf{Trainiertes Modell}\\ LambdaMART-Modell und ausgewählte Hyperparameter }; \node[input] (profiles) at (0,6.3) { \textbf{Kandidatenraum}\\ \texttt{candidate\_profiles.json} }; \node[input] (features) at (5.2,6.3) { \textbf{Featuredefinition}\\ \texttt{feature\_names.json} und aktive Featurekonfiguration }; 
%     % Runtime-Bundle 
%     \node[runtime] (runtimeBundle) at (-3.1,3.8) { \textbf{Runtime-Bundle}\\ Modell, Kandidatenprofile, Feature-Reihenfolge und Laufzeitkonfiguration }; \draw[flow] (model.south) |- (runtimeBundle.west); \draw[flow] (profiles.south) -- (runtimeBundle.north); \draw[flow] (features.south) |- (runtimeBundle.east); 
%     % Nachweispaket 
%     \node[evidence] (evidenceBundle) at (3.1,3.8) { \textbf{Nachweispaket}\\ Run-Manifest, Trainingszusammenfassung, Metriken, Vorhersagen und Prüfberichte }; \draw[flow] (model.south) |- (evidenceBundle.west); \draw[flow] (features.south) |- (evidenceBundle.east); 
%     % Konsistenzprüfung 
%     \node[control] (validation) at (0,1.3) { \textbf{Bundle-Konsistenzprüfung}\\ gleiche Run-Kennung, vollständige Dateien, identische Feature-Reihenfolge und kompatible Konfiguration }; \draw[controlflow] (runtimeBundle.south) |- (validation.west); \draw[controlflow] (evidenceBundle.south) |- (validation.east); 
%     % A4 
%     \node[artifact] (a4) at (0,-1.2) { \textbf{Artefakt A4}\\ versioniertes Runtime-Bundle und reproduzierbares Nachweispaket }; \draw[flow] (validation.south) -- (a4.north); 
%     % Nachfolgende Verwendungen 
%     \node[next] (runtimeUse) at (-3.4,-3.8) { \textbf{Modellanwendung}\\ Bewertung und Sortierung neuer Anfragen }; \node[next] (chapter8) at (3.4,-3.8) { \textbf{Kapitel 8}\\ Evaluation und Interpretation der gespeicherten Ergebnisse }; \draw[nextflow] (a4.south) |- (runtimeUse.north); \draw[nextflow] (a4.south) |- (chapter8.north); \end{tikzpicture} } \caption{ Technische Erzeugung des Artefakts A4 aus dem trainierten LambdaMART-Modell, den Kandidatenprofilen und der Featuredefinition. Das Runtime-Bundle dient der Modellanwendung, während das Nachweispaket die Reproduzierbarkeit und Evaluation unterstützt. } \label{fig:technische_erzeugung_artefakt_a4} \end{figure}

% \section{Technische Sicherstellung der Reproduzierbarkeit} \label{sec:technische_sicherstellung_der_reproduzierbarkeit} Die Reproduzierbarkeit der Implementierung wird durch laufbezogene Ergebnisverzeichnisse, gespeicherte Konfigurationen, Protokolldateien und technische Prüfberichte sichergestellt. Jeder Trainingslauf erhält eine eindeutige Kennung, über die Eingaben, Konfigurationen, Zwischenstände und Ergebnisse miteinander verknüpft werden. Ergänzend überwacht die Implementierung den Pipelinefortschritt und die Ressourcenauslastung. Checkpoints unterstützen die kontrollierte Fortsetzung unterbrochener Verarbeitungsschritte. Neue Artikel und Makros können importiert und durch einen erneuten Pipeline-Lauf in die Datenbasis und das Modell übernommen werden. \subsection{Technische Schwerpunkte} \label{subsec:technische_schwerpunkte_reproduzierbarkeit} Die technische Reproduzierbarkeit umfasst: \begin{itemize} \item eindeutige Kennzeichnung jedes Pipeline-Laufs, \item laufbezogene Speicherung der Konfigurationen und Ergebnisartefakte, \item Dokumentation der verwendeten Datenstände, \item Speicherung der aktiven Feature-, Label- und Modellkonfiguration, \item Protokollierung des Pipelinefortschritts, \item Speicherung von Fehler- und Statusmeldungen, \item Erfassung der CPU-, RAM- und GPU-Auslastung, \item Verwendung von Checkpoints zur Wiederaufnahme unterbrochener Verarbeitungsschritte, \item Import neuer Makros und erneute Ausführung der Pipeline, \item Prüfung eines neuen Runtime-Bundles vor dessen produktiver Aktivierung. \end{itemize} \subsection{Technische Umsetzung} \label{subsec:technische_umsetzung_reproduzierbarkeit} Jeder vollständige Pipeline-Lauf wird einem eigenen Ergebnisverzeichnis zugeordnet. In diesem Verzeichnis werden das trainierte Modell, die Kandidatenprofile, Konfigurationen, Vorhersagen, Metriken und Prüfberichte gemeinsam gespeichert. Eine eindeutige Run-Kennung verbindet die Artefakte eines Laufs. Dadurch wird verhindert, dass Kandidatenprofile, Featuredefinitionen oder Evaluationsdateien unterschiedlicher Modellstände ohne Kennzeichnung zusammengeführt werden. Die Reproduzierbarkeit eines Laufs setzt mindestens folgende Informationen voraus: \begin{itemize} \item Kennung und Zeitpunkt des Pipeline-Laufs, \item verwendeter Datenmodus, \item Eingabe- und Ausgabepfade, \item Version oder Fingerprint der Eingabedaten, \item Splitkonfiguration und Zufallsseed, \item aktive Featurekonfiguration, \item verwendete Profilgrenzen, \item aktive Relevanzlabel- und Gain-Konfiguration, \item getestete und ausgewählte LambdaMART-Hyperparameter, \item verwendete Software- und Bibliotheksversionen. \end{itemize} Die zentrale Zusammenfassung dieser Angaben erfolgt über das \texttt{run\_manifest.json} und die \texttt{training\_summary.json}. Das Run-Manifest dokumentiert den technischen Kontext des Laufs. Die Trainingszusammenfassung enthält die ausgewählte Modellkonfiguration und die während des Trainings erzeugten Kennzahlen. Ein Lauf ist reproduzierbar, wenn dieselben Eingabedaten, Konfigurationen und Softwarestände erneut verarbeitet werden können: \begin{equation} R = \left( D, S, F, L, H, E \right), \label{eq:reproduzierbarer_lauf} \end{equation} wobei $D$ den Datenstand, $S$ die Splitkonfiguration, $F$ die Featurekonfiguration, $L$ die Labelkonfiguration, $H$ die Hyperparameter und $E$ die Ausführungsumgebung bezeichnet. \subsubsection{Status- und Fehlerprotokollierung} \label{subsubsec:status_und_fehlerprotokollierung} Die Pipeline schreibt ihren aktuellen Ausführungszustand in die Datei \texttt{live\_status.jsonl}. Die einzelnen Einträge enthalten Statusinformationen zu den ausgeführten Verarbeitungsschritten. Dadurch können Fortschritt, Laufzeit und Fehlerzustände während der Ausführung nachvollzogen werden. Die Anwendungs- und Dashboardmodule lesen diese Statusdatei ein und stellen den aktuellen Pipelinezustand dar. Zu den beteiligten Modulen gehören insbesondere: \begin{itemize} \item \texttt{app\_1020\_evaluation\_dashboard.py}, \item \texttt{app\_1021\_dashboard\_evaluation.py}, \item \texttt{app\_1022\_dashboard\_overview.py}, \item \texttt{app\_1024\_dashboard\_log\_viewer.py}. \end{itemize} Die grafische Oberfläche übernimmt keine Modelllogik. Sie startet die Pipeline als separaten Prozess, liest die erzeugten Status- und Ergebnisdateien und stellt diese für die technische Kontrolle zusammen. Fehler werden gemeinsam mit dem betroffenen Verarbeitungsschritt protokolliert. Ein fehlgeschlagener Pipeline-Lauf darf nicht als gültiges Artefakt A4 freigegeben werden. Bereits erzeugte Teildateien bleiben dem fehlgeschlagenen Lauf zugeordnet und dürfen nicht mit dem zuletzt freigegebenen Runtime-Bundle vermischt werden. \subsubsection{Ressourcenmonitoring} \label{subsubsec:ressourcenmonitoring} Während der Pipelineausführung werden die verwendeten Systemressourcen erfasst. Die Implementierung berücksichtigt insbesondere: \begin{itemize} \item CPU-Auslastung, \item Arbeitsspeicherbelegung, \item GPU-Auslastung, sofern eine GPU verfügbar ist, \item GPU-Speicherbelegung, \item Laufzeit der Pipeline und ihrer Teilschritte. \end{itemize} Die Messwerte werden unter anderem in \texttt{resource\_monitoring.json} und \texttt{resource\_monitoring.csv} gespeichert. Sie dienen der technischen Analyse des Trainingslaufs und ermöglichen den Vergleich des Ressourcenbedarfs mehrerer Runs. Das Ressourcenmonitoring beeinflusst weder die Relevanzlabels noch die Modellmerkmale. Die Messwerte gehören deshalb zum Nachweispaket und nicht zum Runtime-Bundle. \subsubsection{Checkpoints und Wiederanlauf} \label{subsubsec:checkpoints_und_wiederanlauf} Checkpoints dokumentieren erfolgreich abgeschlossene Pipelineabschnitte. Sie ermöglichen die kontrollierte Fortsetzung nach einem Fehler, ohne unveränderte und bereits abgeschlossene Schritte erneut auszuführen. Vor der Wiederverwendung eines Checkpoints muss geprüft werden, ob die zugehörigen Eingaben und Konfigurationen unverändert geblieben sind. Ein Checkpoint ist nur gültig, wenn \begin{equation} \operatorname{validCheckpoint} = \operatorname{sameData} \land \operatorname{sameConfig} \land \operatorname{completeOutput} \label{eq:gueltiger_checkpoint} \end{equation} erfüllt ist. Bei Änderungen an Rohdaten, Splits, Featuredefinitionen, Profilgrenzen, Labelregeln oder Modellparametern müssen alle davon abhängigen Pipeline-Stufen erneut ausgeführt werden. Ein technisch vorhandener Checkpoint darf nicht verwendet werden, wenn seine Ergebnisse nicht mehr zur aktiven Konfiguration passen. Die Checkpoints dienen damit der Fehlerbehandlung innerhalb eines definierten Laufs. Sie ersetzen keine vollständige Versionierung der erzeugten Modell- und Datenartefakte. \subsubsection{Integration neuer Artikel und Makros} \label{subsubsec:integration_neuer_artikel_und_makros} Der Import neuer Makros wird durch \texttt{tool\_1110\_import\_new\_macro.py} unterstützt. Das Werkzeug übernimmt neue EMA-Dateien in den vorgesehenen Datenbestand und stößt anschließend die abhängigen Pipeline-Schritte an. Der technische Ablauf umfasst: \begin{enumerate} \item Import und Prüfung der neuen Quelldateien, \item Erkennung neuer oder veränderter Artikel und Makros, \item erneute Datenaufbereitung, \item erneute gruppierte Datenaufteilung, \item Neubildung der Kandidatenprofile, \item erneute Erzeugung der Rankingmatrizen, \item erneutes LambdaMART-Training, \item erneute technische Evaluation, \item Erzeugung eines neuen versionierten Artefakts A4. \end{enumerate} Die Implementierung ruft hierfür die vorgesehenen Aufbereitungs-, Split-, Trainings- und Evaluationsschritte auf. Ein neuer Datenstand verändert damit nicht das bestehende Runtime-Bundle, sondern erzeugt einen neuen Lauf mit eigenständiger Run-Kennung. Neue Artikel und Makros dürfen nicht direkt in bereits serialisierte Kandidatenprofile eingefügt werden. Die Profile müssen aus dem aktuellen Trainingsanteil reproduzierbar neu gebildet werden. Andernfalls wären die gespeicherten Profile nicht mehr eindeutig auf einen Daten- und Splitstand zurückzuführen. \subsubsection{Freigabe eines neuen Runtime-Bundles} \label{subsubsec:freigabe_neues_runtime_bundle} Ein erfolgreich beendeter Trainingslauf wird nicht automatisch zum produktiven Modellstand. Vor der Aktivierung sind mindestens folgende Prüfungen erforderlich: \begin{itemize} \item vollständige Erzeugung aller Pflichtdateien, \item erfolgreiche Bundle-Konsistenzprüfung, \item bestandene Leakage- und Splitprüfungen, \item Übereinstimmung von Modell und Featuredefinition, \item erfolgreiche technische Evaluation, \item Erreichen der festgelegten Qualitätsanforderungen, \item eindeutige Versionierung des neuen Runtime-Bundles. \end{itemize} Die produktive Aktivierung muss atomar erfolgen. Modell, Kandidatenprofile, Featuredefinition und Laufzeitkonfiguration dürfen nicht einzeln ausgetauscht werden. Stattdessen wird die Referenz auf das vollständige Bundle erst nach erfolgreicher Prüfung auf die neue Version umgestellt. Der Freigabeprozess lässt sich formal durch \begin{equation} B_{\mathrm{active}}^{(t+1)} = \begin{cases} B_{\mathrm{new}}, & \operatorname{valid}(B_{\mathrm{new}}) \land \operatorname{approved}(B_{\mathrm{new}}), \\[1mm] B_{\mathrm{active}}^{(t)}, & \text{sonst} \end{cases} \label{eq:atomare_bundle_freigabe} \end{equation} beschreiben. Bei einer fehlgeschlagenen Prüfung bleibt das bisherige Runtime-Bundle aktiv. Dadurch wird verhindert, dass ein unvollständiger oder inkonsistenter Modellstand produktiv verwendet wird. \subsection{Abgleich zwischen bisheriger Arbeit und Code} \label{subsec:abgleich_reproduzierbarkeit} \subsubsection*{Übereinstimmungen} Die bisherige Arbeit und der Code sehen eine wiederholbare Ausführung der Datenaufbereitung, des Splits, des Trainings und der Evaluation vor. Die Speicherung laufbezogener Konfigurationen und Ergebnisse unterstützt die geforderte Nachvollziehbarkeit der Modellbildung. Der Import neuer Makros und die erneute Ausführung abhängiger Pipeline-Stufen entsprechen dem vorgesehenen Aktualisierungsprozess. Auch die getrennte Speicherung der Laufartefakte ist mit der Versionierung des Empfehlungssystems vereinbar. \subsubsection*{Technische Konkretisierungen} Der Code konkretisiert die Reproduzierbarkeit durch Run-Manifeste, Trainingszusammenfassungen, Statusdateien, Logdateien, Ressourcenmessungen und Checkpoints. Die Dashboardmodule stellen diese Informationen dar und ermöglichen die technische Überwachung der Pipeline. Das Importwerkzeug automatisiert die Aufnahme neuer EMA-Dateien und den Aufruf der nachfolgenden Verarbeitungsschritte. Die Implementierung erweitert die bisherige Architektur damit um operative Funktionen zur Ausführung, Fehleranalyse und Ressourcenüberwachung. Diese Funktionen verändern das Rankingmodell nicht, sichern jedoch dessen reproduzierbare Erzeugung ab. \subsubsection*{Abweichungen} Die bisherige Arbeit fordert einen reproduzierbaren nächtlichen Pipeline-Neulauf mit versionierten Artefakten und einem atomaren Runtime-Bundle-Release. Der Code unterstützt den erneuten Pipeline-Lauf, die laufbezogene Speicherung und die technische Überwachung. Aus dem vorhandenen Import-, Trainings- und Dashboardablauf ergibt sich jedoch nicht automatisch ein vollständiger atomarer Produktiv-Release. Das Erzeugen einzelner Dateien und das erfolgreiche Beenden der Pipeline stellen noch nicht sicher, dass Modell, Kandidatenprofile, Featuredefinition und Konfiguration gemeinsam aktiviert werden. Eine weitere Abweichung besteht zwischen Checkpoint-Wiederaufnahme und vollständiger Reproduzierbarkeit. Checkpoints ermöglichen die Fortsetzung eines unterbrochenen Laufs. Sie belegen jedoch nicht, dass ein Lauf ausgehend von unveränderten Rohdaten vollständig reproduziert werden kann. Dafür müssen zusätzlich Datenstände, Konfigurationen, Seeds und Softwareversionen gespeichert werden. Das Ressourcenmonitoring ist im Code umfangreicher ausgeprägt als in der bisherigen fachlichen Beschreibung. CPU-, RAM- und GPU-Werte sind technische Nachweisdaten, jedoch keine fachlichen Qualitätskennzahlen des Empfehlungssystems. \subsubsection*{Konsequenz für die Dokumentation} Die technische Dokumentation unterscheidet künftig zwischen: \begin{description} \item[Wiederanlauf] Fortsetzung eines unterbrochenen Pipeline-Laufs mithilfe gültiger Checkpoints. \item[Reproduktion] erneute vollständige Erzeugung eines Modellstands aus versionierten Daten, Konfigurationen und Softwareständen. \item[Aktualisierung] Erzeugung eines neuen Modells nach dem Import neuer oder veränderter Artikel und Makros. \item[Freigabe] kontrollierte und atomare Aktivierung eines vollständig geprüften Runtime-Bundles. \end{description} Der vorhandene Code deckt Pipeline-Neulauf, Statusprotokollierung, Monitoring und Checkpoints ab. Der atomare produktive Bundle-Austausch ist als gesonderter Freigabeschritt zu dokumentieren und, sofern noch nicht vollständig umgesetzt, als offener Implementierungspunkt zu kennzeichnen. Ein Trainingslauf gilt nur dann als reproduzierbar, wenn neben den Ergebnisdateien auch Datenstand, Splitkonfiguration, Featuredefinition, Labelpolicy, Modellparameter, Seed und Softwareumgebung eindeutig festgehalten werden. \subsection{Zusammenhang mit der bisherigen Arbeit} \label{subsec:zusammenhang_reproduzierbarkeit} Die Reproduzierbarkeit setzt die in den vorherigen Kapiteln formulierten Anforderungen an Nachvollziehbarkeit, Erweiterbarkeit und Aktualisierbarkeit technisch um. Artefakt A1 wird durch versionierte Datenstände und Datensplits abgesichert. Die Konfiguration von A2 wird durch die gespeicherten LambdaMART-Parameter dokumentiert. Die in A3 definierte Pipeline wird durch Statusdateien, Checkpoints und Laufmanifeste ausführbar und überprüfbar gemacht. Die in Abschnitt~\ref{sec:erzeugung_des_artefakts_a4} beschriebenen Bestandteile von A4 werden einem eindeutigen Trainingslauf zugeordnet. Dadurch bleiben Modell, Kandidatenprofile, Featuredefinition und Nachweispaket konsistent. Die gespeicherten Metriken und Vorhersagen bilden die technische Grundlage für Kapitel~8. Das Ressourcenmonitoring und die Pipelineprotokolle dienen dagegen der technischen Bewertung der Ausführung und nicht der fachlichen Beurteilung der Empfehlungsqualität. \subsection{Darstellung der reproduzierbaren Pipelineausführung} \label{subsec:darstellung_reproduzierbare_pipelineausfuehrung} Abbildung~\ref{fig:technische_reproduzierbarkeit} zeigt den Zusammenhang zwischen neuen Daten, Pipelineausführung, Laufüberwachung, Artefakterzeugung und produktiver Freigabe. Der gestrichelte Freigabepfad kennzeichnet, dass ein erzeugtes Bundle erst nach erfolgreicher Prüfung aktiviert werden darf. 
% \begin{figure}[H] \centering \resizebox{\textwidth}{!}{ \begin{tikzpicture}[ >=Stealth, font=\small, input/.style={ draw=blue!65!black, fill=blue!8, rounded corners=2mm, text width=4.5cm, minimum height=1.4cm, align=center, inner sep=7pt }, process/.style={ draw=green!55!black, fill=green!9, rounded corners=2mm, text width=4.8cm, minimum height=1.45cm, align=center, inner sep=7pt }, monitoring/.style={ draw=gray!70!black, fill=gray!8, rounded corners=2mm, text width=4.7cm, minimum height=1.4cm, align=center, inner sep=7pt }, checkpoint/.style={ draw=violet!65!black, fill=violet!7, rounded corners=2mm, text width=4.7cm, minimum height=1.4cm, align=center, inner sep=7pt }, artifact/.style={ draw=orange!75!black, fill=orange!12, rounded corners=2mm, text width=5.1cm, minimum height=1.5cm, align=center, inner sep=7pt }, control/.style={ draw=red!65!black, fill=red!4, dashed, rounded corners=2mm, text width=5.0cm, minimum height=1.4cm, align=center, inner sep=7pt }, active/.style={ draw=teal!65!black, fill=teal!8, rounded corners=2mm, text width=4.8cm, minimum height=1.4cm, align=center, inner sep=7pt }, flow/.style={ ->, thick }, monitorflow/.style={ ->, thick, dashed, gray!70!black }, releaseflow/.style={ ->, very thick, dashed, red!65!black } ] 
%     % Eingangsdaten 
%     \node[input] (newData) at (-4.2,6.6) { \textbf{Aktueller Datenstand}\\ vorhandene sowie neue Artikel und Makros }; \node[input] (configuration) at (4.2,6.6) { \textbf{Laufkonfiguration}\\ Splits, Features, Labels, Hyperparameter und Seed }; 
%     % Pipeline 
%     \node[process] (pipeline) at (0,4.2) { \textbf{Versionierter Pipeline-Lauf}\\ Datenaufbereitung, Split, Profilbildung, Training und Evaluation }; \draw[flow] (newData.south) |- (pipeline.west); \draw[flow] (configuration.south) |- (pipeline.east); 
%     % Monitoring und Checkpoints 
%     \node[monitoring] (monitoring) at (-4.2,1.8) { \textbf{Protokollierung und Monitoring}\\ \texttt{live\_status.jsonl}, Logs sowie CPU-, RAM- und GPU-Werte }; \node[checkpoint] (checkpoints) at (4.2,1.8) { \textbf{Checkpoints}\\ Wiederanlauf nur bei identischem Daten- und Konfigurationsstand }; \draw[monitorflow] (pipeline.south) |- (monitoring.east); \draw[monitorflow] (pipeline.south) |- (checkpoints.west); 
%     % Neues A4 
%     \node[artifact] (newBundle) at (0,-0.7) { \textbf{Neues Artefakt A4}\\ versioniertes Runtime-Bundle und zugehöriges Nachweispaket }; \draw[flow] (pipeline.south) -- (newBundle.north); \draw[monitorflow] (monitoring.south) |- (newBundle.west); \draw[monitorflow] (checkpoints.south) |- (newBundle.east); 
%     % Prüfung 
%     \node[control] (releaseCheck) at (0,-3.0) { \textbf{Freigabeprüfung}\\ Vollständigkeit, Konsistenz, Leakage-Sicherheit und Qualitätsanforderungen }; \draw[flow] (newBundle.south) -- (releaseCheck.north); 
%     % Aktives Bundle 
%     \node[active] (activeBundle) at (0,-5.4) { \textbf{Aktives Runtime-Bundle}\\ atomare Aktivierung der vollständig geprüften Version }; \draw[releaseflow] (releaseCheck.south) -- (activeBundle.north); 
%     % Rückfallpfad 
%     \node[control] (fallback) at (6.2,-4.2) { Bei fehlgeschlagener Prüfung bleibt das bisherige Bundle aktiv. }; \draw[releaseflow] (releaseCheck.east) -| (fallback.north); \end{tikzpicture} } \caption{ Technische Sicherstellung der Reproduzierbarkeit durch versionierte Pipeline-Läufe, Statusprotokollierung, Ressourcenmonitoring, Checkpoints und die kontrollierte Freigabe eines vollständigen Runtime-Bundles. } \label{fig:technische_reproduzierbarkeit} \end{figure}
%     % \section{Bildung der Vergleichsmerkmale} 

% \section{Zusammenfassung} \label{sec:zusammenfassung_implementierung} Kapitel~7 hat die technische Umsetzung der zuvor entwickelten Konzeption dargestellt. Hierzu wurden die aufbereitete Datenbasis A1, der ausgewählte Modellansatz A2 und die Systemarchitektur A3 in einer modularen Python-Pipeline zusammengeführt. Das Ergebnis ist das Artefakt A4 als trainiertes und reproduzierbar anwendbares Empfehlungssystem. \subsection{Technische Schwerpunkte} \label{subsec:technische_schwerpunkte_zusammenfassung} Die Implementierung umfasst: \begin{itemize} \item die modulare Gliederung der Software in Konfigurations-, Daten-, Pipeline-, Anwendungs- und Werkzeugmodule, \item die strukturierte Speicherung von Makrostruktur-, Artikel- und Lexikalmerkmalen, \item die Bildung begrenzter Kandidatenprofile ausschließlich aus dem Trainingsanteil, \item die Erzeugung numerischer Anfrage-Kandidat-Merkmale mit einer festen Feature-Reihenfolge, \item die Zuordnung konfigurierbarer Relevanzlabels, \item den Ausschluss von Identitäts-, Ziel-, Prozess- und redundanten Merkmalsinformationen, \item das Training und die Validierung des LambdaMART-Modells mit LightGBM, \item die Zusammenführung von Modell, Kandidatenprofilen, Featuredefinition und Konfiguration zum Runtime-Bundle, \item die Speicherung eines ergänzenden Nachweispakets, \item die Protokollierung, Überwachung und Wiederholung der Pipelineausführung. \end{itemize} \subsection{Ergebnis der Implementierung} \label{subsec:ergebnis_der_implementierung} Die Implementierung bildet die konzeptionell festgelegten Verarbeitungsschritte durch die Pipeline-Stufen 910 bis 940 ab. Die aufbereiteten Artikelinformationen werden zunächst in getrennte Makrostruktur-, Artikel- und Lexikalrepräsentationen überführt. Diese Repräsentationen bilden die Grundlage für die Aggregation der Kandidatenprofile. Die Kandidatenprofile werden nach Zielmakro gruppiert und ausschließlich aus den Trainingsartikeln erzeugt. Die aggregierten Merkmalswerte werden entsprechend ihrer Häufigkeit und der konfigurierten dimensionsbezogenen Obergrenzen in das Profil übernommen. Ergänzende Support-, Familien- und Fingerprintinformationen unterstützen die technische Verwaltung und die spätere Merkmalsberechnung. Für jede Anfrage wird eine Rankinggruppe aus den verfügbaren Kandidaten gebildet. Die Anfrage- und Kandidatenrepräsentationen werden nicht direkt an LambdaMART übergeben. Aus jeder Anfrage-Kandidat-Kombination entsteht zunächst ein numerischer Paarmerkmalsvektor. Die resultierende Featurematrix wird gemeinsam mit den Relevanzlabels und Gruppengrößen an den LightGBM-Ranker übergeben. Die Modellmerkmale werden vor dem Training auf unzulässige Ziel-, Identitäts- und Prozessinformationen geprüft. Ergänzend verhindern Rollenbereinigung und Redundanzfilter, dass dieselbe Information mehrfach oder über eine unzulässige Repräsentation in die Featurematrix eingeht. Die Leave-one-out-Diagnostik prüft die Eigenevidenz eines Trainingsartikels im Kandidatenprofil seines Zielmakros. Das LambdaMART-Modell wird auf dem Trainingsanteil gelernt. Die Auswahl der Hyperparameter erfolgt anhand des Validierungsanteils. Der Testanteil wird erst nach Festlegung der Modellkonfiguration ausgewertet und besitzt keine Rückwirkung auf Featureauswahl, Labelpolitik, Profilbildung oder Hyperparameteroptimierung. Das erzeugte Artefakt A4 besteht aus zwei Teilen: \begin{equation} A4 = A4_{\mathrm{runtime}} \cup A4_{\mathrm{evidence}}. \label{eq:zusammensetzung_artefakt_a4} \end{equation} Das Runtime-Bundle enthält: \begin{itemize} \item das trainierte LambdaMART-Modell, \item die Kandidatenprofile, \item die verbindliche Featuredefinition, \item die aktive Feature- und Laufzeitkonfiguration. \end{itemize} Das Nachweispaket enthält: \begin{itemize} \item das Run-Manifest, \item die Trainingszusammenfassung, \item die Validierungs- und Testmetriken, \item die gespeicherten Vorhersagen, \item die Leakage-, Split- und Qualitätsberichte, \item die Status- und Ressourcenprotokolle. \end{itemize} Die Run-Kennung verbindet beide Teile mit dem verwendeten Daten-, Split- und Konfigurationsstand. Ein Runtime-Bundle ist nur gültig, wenn Modell, Kandidatenprofile, Featuredefinition und Konfiguration demselben Trainingslauf zugeordnet sind. \subsection{Abgleich zwischen bisheriger Arbeit und Code} \label{subsec:abgleich_zusammenfassung_implementierung} \subsubsection*{Übereinstimmungen} Die grundlegende Architektur des Codes entspricht der bisherigen Konzeption. Dies betrifft insbesondere: \begin{itemize} \item die Trennung von Datenaufbereitung, Datenaufteilung, Training und Evaluation, \item die getrennte Verarbeitung von Artikel-, Makrostruktur- und Lexikalmerkmalen, \item die Bildung der Kandidatenprofile aus dem Trainingsanteil, \item die Erzeugung von Anfrage-Kandidat-Paaren, \item die Verwendung von LambdaMART als Rankingverfahren, \item die Trennung von Training, Validierung und Test, \item den Ausschluss direkter Makroidentitäten aus den trainierbaren Merkmalen. \end{itemize} Die Implementierung setzt damit die in Kapitel~6 festgelegte Gesamtarchitektur um. Die Unterschiede betreffen überwiegend die technische Konkretisierung und nicht den grundlegenden Lösungsansatz. \subsubsection*{Technische Konkretisierungen} Der Code konkretisiert die bisherige Arbeit in mehreren Punkten: \begin{itemize} \item Das fachliche Makroprofil wird als erweitertes Kandidatenprofil realisiert. \item Die Profilwerte werden häufigkeitsbasiert und dimensionsabhängig begrenzt. \item Support-, Familien-, Singleton- und Fingerprintinformationen ergänzen die aggregierten Profile. \item Anfrage- und Kandidateninformationen werden durch spezialisierte Vergleichsfunktionen in numerische Paarmerkmale überführt. \item Die Relevanzlabel-Policy ist konfigurierbar und mit den LightGBM-Gain-Werten gekoppelt. \item Leakage-Sicherung, Rollenbereinigung und Leave-one-out-Prüfung ergänzen den Trainingsprozess. \item Run-Manifeste, Checkpoints, Statusdateien und Ressourcenmonitoring sichern die Pipelineausführung ab. \end{itemize} Diese Bestandteile erweitern die konzeptionelle Beschreibung um die für eine ausführbare und reproduzierbare Softwarelösung erforderlichen Implementierungsdetails. \subsubsection*{Abweichungen} Zwischen der bisherigen Formulierung und dem Code bestehen folgende wesentliche Abweichungen: \begin{enumerate} \item Die Artikel-, Makrostruktur- und Lexikalrepräsentationen sind keine direkten LambdaMART-Eingaben. Das Modell verarbeitet ausschließlich numerische Anfrage-Kandidat-Merkmale. \item Die Kandidatenprofile enthalten nicht unbegrenzt alle beobachteten Merkmalswerte. Die gespeicherten Werte werden häufigkeitsbasiert und dimensionsabhängig begrenzt. \item Anfrageprofile werden überwiegend während der Matrixerzeugung aus den splitbezogenen Artikeldaten abgeleitet. Eine separate persistierte Ablage ist für die Trainingsimplementierung nicht zwingend erforderlich. \item Die Labelimplementierung unterstützt sowohl ein binäres als auch ein abgestuftes Relevanzschema. Maßgeblich ist die im finalen Trainingslauf aktivierte Konfiguration. \item Artefakt A4 besteht nicht nur aus der trainierten Modelldatei. Für die Anwendung werden zusätzlich Kandidatenprofile, Featuredefinition und Laufzeitkonfiguration benötigt. \item Dashboard, Ressourcenmonitoring, Checkpoints und Importwerkzeuge sind im Code umfangreicher ausgeprägt als in der bisherigen konzeptionellen Beschreibung. \item Der erneute Pipeline-Lauf und die versionierte Speicherung sind implementiert. Die atomare produktive Aktivierung eines vollständigen Runtime-Bundles ist davon zu unterscheiden und gegebenenfalls als offener Implementierungspunkt zu behandeln. \end{enumerate} Ein zusätzlicher Konsistenzpunkt betrifft die Makrostrukturmerkmale auf der Anfrageseite. Diese sind nur zulässig, wenn sie ohne Kenntnis des Zielmakros aus dem Artikel oder dem zugehörigen Datenblatt abgeleitet werden. Eine Übernahme aus der bekannten Artikel-Makro-Zuordnung würde eine Zielfleckage verursachen. \subsubsection*{Konsequenz für die Dokumentation} Die technische Dokumentation verwendet folgende Begriffsabgrenzung: \begin{description} \item[Makroprofil] bezeichnet das fachliche Konzept der aggregierten Kandidatenrepräsentation. \item[Kandidatenprofil] bezeichnet die technische Umsetzung des Makroprofils einschließlich Support-, Familien-, Fingerprint- und Laufzeitinformationen. \item[Rollenrepräsentation] bezeichnet die strukturierten Makrostruktur-, Artikel- oder Lexikalinformationen. \item[Paarmerkmalsvektor] bezeichnet die numerischen Vergleichsmerkmale einer Anfrage-Kandidat-Kombination. \item[Runtime-Bundle] bezeichnet die für die Modellanwendung erforderliche Einheit aus Modell, Kandidatenprofilen, Featuredefinition und Konfiguration. \item[Nachweispaket] bezeichnet die ergänzenden Dateien zur Reproduktion und Evaluation des Trainingslaufs. \end{description} Nicht verwendete Legacy- und Alternativkonfigurationen werden nicht als Bestandteil des finalen Modells dargestellt. Für die Beschreibung des Artefakts A4 sind ausschließlich die im finalen Trainingslauf tatsächlich aktiven Komponenten maßgeblich. \subsection{Zusammenhang mit der bisherigen Arbeit} \label{subsec:zusammenhang_zusammenfassung_implementierung} Kapitel~7 führt die Artefakte A1, A2 und A3 technisch zusammen. Artefakt A1 stellt die aufbereiteten Daten, Gruppierungen und Splits bereit. Artefakt A2 legt LambdaMART als Rankingverfahren fest. Artefakt A3 definiert die Datenrepräsentationen, Kandidatenprofile, Vergleichsmerkmale und Verarbeitungsschritte. Die technische Umsetzung ordnet diese Vorgaben konkreten Python-Modulen, Funktionen, Konfigurationen und Dateien zu. Das Ergebnis ist das Artefakt A4 als versioniertes Runtime-Bundle mit zugehörigem Nachweispaket. Damit beantwortet Kapitel~7 die Frage, wie die zuvor entwickelte Architektur technisch umgesetzt wurde. Die fachliche Eignung des Modellansatzes wurde bereits in den vorherigen Kapiteln begründet. Die mit der Implementierung erreichte Empfehlungsqualität ist nicht Gegenstand dieses Kapitels. Kapitel~8 verwendet die in A4 gespeicherten Testvorhersagen, Metriken und Prüfberichte zur Demonstration und Evaluation des Empfehlungssystems. Dort werden insbesondere die Rankingqualität, die Generalisierungsfähigkeit, Fehlerfälle und die Wirkung der implementierten Merkmale untersucht. \subsection{Darstellung der Gesamtimplementierung} \label{subsec:darstellung_gesamtimplementierung} Abbildung~\ref{fig:gesamtimplementierung_kapitel_sieben} fasst die technische Verarbeitung von den Eingangsartefakten bis zum Artefakt A4 zusammen. Die Evaluation in Kapitel~8 ist als nachgelagerter Schritt dargestellt und besitzt keinen Rückfluss in den bereits abgeschlossenen Testprozess. 
% \begin{figure}[H] \centering \resizebox{\textwidth}{!}{ \begin{tikzpicture}[ >=Stealth, font=\small, input/.style={ draw=blue!65!black, fill=blue!8, rounded corners=2mm, text width=3.8cm, minimum height=1.35cm, align=center, inner sep=6pt }, process/.style={ draw=green!55!black, fill=green!9, rounded corners=2mm, text width=4.5cm, minimum height=1.4cm, align=center, inner sep=7pt }, control/.style={ draw=red!65!black, fill=red!4, dashed, rounded corners=2mm, text width=4.5cm, minimum height=1.3cm, align=center, inner sep=6pt }, bundle/.style={ draw=orange!75!black, fill=orange!12, rounded corners=2mm, text width=5.1cm, minimum height=1.5cm, align=center, inner sep=7pt }, evaluation/.style={ draw=violet!65!black, fill=violet!7, rounded corners=2mm, text width=5.0cm, minimum height=1.4cm, align=center, inner sep=7pt }, flow/.style={ ->, thick }, controlflow/.style={ ->, thick, dashed, red!65!black }, evaluationflow/.style={ ->, thick, dashed, violet!65!black } ] 
%     % Eingangsartefakte 
%     \node[input] (a1) at (-5.2,7.0) { \textbf{A1}\\ Datenbasis und Splits }; \node[input] (a2) at (0,7.0) { \textbf{A2}\\ LambdaMART }; \node[input] (a3) at (5.2,7.0) { \textbf{A3}\\ Systemarchitektur }; \coordinate (inputJoin) at (0,5.8); \draw[thick] (a1.south) |- (inputJoin); \draw[thick] (a2.south) -- (inputJoin); \draw[thick] (a3.south) |- (inputJoin); 
%     % Datenrepräsentationen 
%     \node[process] (representations) at (0,4.6) { \textbf{Datenrepräsentationen}\\ Makrostruktur-, Artikel- und Lexikalinformationen }; \draw[flow] (inputJoin) -- (representations.north); 
%     % Kandidatenprofile 
%     \node[process] (profiles) at (-3.5,2.4) { \textbf{Kandidatenprofile}\\ begrenzte Aggregation der Trainingsartikel je Zielmakro }; 
%     % Anfrage-Kandidat-Matrix 
%     \node[process] (matrix) at (3.5,2.4) { \textbf{Anfrage-Kandidat-Matrix}\\ numerische Paarmerkmale, Labels und Anfragegruppen }; \draw[flow] (representations.south) |- (profiles.north); \draw[flow] (representations.south) |- (matrix.north); \draw[flow] (profiles.east) -- (matrix.west); 
%     % Leakage-Kontrolle 
%     \node[control] (leakage) at (7.4,2.4) { \textbf{Leakage-Sicherung}\\ Ziel-, Identitäts-, Prozess- und redundante Merkmale ausschließen }; \draw[controlflow] (leakage.west) -- (matrix.east); 
%     % Training 
%     \node[process] (training) at (0,0.0) { \textbf{LambdaMART-Training}\\ Training, Validierungsauswahl und abschließende Testbewertung }; \draw[flow] (matrix.south) |- (training.east); 
%     % A4 
%     \node[bundle] (runtime) at (-3.2,-2.5) { \textbf{Runtime-Bundle}\\ Modell, Kandidatenprofile, Featuredefinition und Konfiguration }; \node[bundle] (evidence) at (3.2,-2.5) { \textbf{Nachweispaket}\\ Manifest, Zusammenfassung, Metriken, Vorhersagen und Prüfberichte }; \draw[flow] (training.south) |- (runtime.north); \draw[flow] (training.south) |- (evidence.north); \node[ draw=orange!75!black, rounded corners=3mm, fit=(runtime)(evidence), inner sep=11pt, label={ [font=\bfseries, text=orange!75!black] below:Artefakt A4 } ] (a4) {}; 
%     % Kapitel 8 
%     \node[evaluation] (chapter8) at (0,-5.7) { \textbf{Kapitel 8}\\ Demonstration, Evaluation und Interpretation der Ergebnisse }; \draw[evaluationflow] (a4.south) -- (chapter8.north); \end{tikzpicture} } \caption{ Zusammenfassung der technischen Implementierung. Die Artefakte A1, A2 und A3 werden über Datenrepräsentationen, Kandidatenprofile, Anfrage-Kandidat-Matrizen, Leakage-Sicherung und LambdaMART-Training in das Artefakt A4 überführt. Die wissenschaftliche Evaluation erfolgt anschließend in Kapitel~8. } \label{fig:gesamtimplementierung_kapitel_sieben} \end{figure}

% % Zur Veranschaulichung wird das Motorbeispiel aus Abschnitt \ref{sec:herleitung} fortgeführt. 
% % Das Anfrageprofil enthält unter anderem den Hersteller Siemens, eine Spannung von \(400\,V\), einen Strom von \(10\,A\), den Produkttyp Elektromotor, 
% % den \textit{technischen Hinweis} ``Wasserdichte Ausführung'' sowie die Anschlussstruktur U1, V1 und W1. 
% % Das betrachtete Makroprofil enthält die Hersteller Siemens und SEW, die Spannungen \(380\,V\), \(400\,V\) und \(420\,V\), einen Strom von \(10\,A\), den Makrotyp Elektromotor, 
% % den \textit{technischen Hinweis} ``Wasserdichte Ausführung'' sowie ebenfalls die Anschlussstruktur U1, V1 und W1. Die Makrobezeichnung \texttt{M1\_F1\_V1} ist dem Makroprofil zwar zugeordnet, 
% % wird jedoch nicht zur Bildung der Vergleichsmerkmale verwendet. 
% % Sie beschreibt die Identität des Makros, während im Anfrageprofil kein zulässiger Vergleichswert vorhanden ist. 
% % Ihre Verwendung würde daher unmittelbar auf die bekannte Artikel-Makro-Zuordnung verweisen und nicht die technische Eignung des Makrokandidaten beschreiben. 

% % \paragraph{Merkmalsübereinstimmungen} 

% % Merkmalsübereinstimmungen prüfen, ob eine einzelne Eigenschaft des Anfrageprofils auch im Makroprofil enthalten ist. 
% % Das Ergebnis kann als binärer Wert dargestellt werden: \[ \operatorname{Match}(a,B) = \begin{cases} 1, & \text{wenn } a \in B,\\ 0, & \text{wenn } a \notin B. \end{cases} \] 
% % Für den Hersteller ergibt sich im Motorbeispiel: \[ a=\text{Siemens}, \qquad B=\{\text{Siemens},\text{SEW}\}. \] 
% % Da Siemens im Makroprofil enthalten ist, gilt: \[ \operatorname{Match}(a,B)=1. \] 
% % Entsprechend können unter anderem Übereinstimmungen des Produkttyps, der Klassifikation, der Spannungsklasse oder der Stromklasse ermittelt werden. 
% % Der Wert \(1\) kennzeichnet dabei lediglich eine Übereinstimmung. 
% % Welche Bedeutung diese Übereinstimmung für die Rangfolge besitzt, wird während des Trainings durch LambdaMART gelernt. 

% % \paragraph{Mengenbasierte Ähnlichkeiten} 

% % Da ein Makroprofil aus mehreren zugeordneten Artikeln gebildet wird, kann eine Eigenschaft mehrere Werte enthalten. Mengenbasierte Ähnlichkeiten beschreiben, 
% % wie stark sich solche Merkmalsmengen des Anfrageprofils und des Makroprofils überschneiden. 
% % Für die Spannung werden die folgenden Mengen betrachtet: \[ A=\{400\,V\} \] und \[ B=\{380\,V,\;400\,V,\;420\,V\}. \] 
% % Ihre Schnittmenge und ihre Vereinigungsmenge lauten: \[ A\cap B=\{400\,V\} \] und \[ A\cup B=\{380\,V,\;400\,V,\;420\,V\}. \] 
% % Die Jaccard-Ähnlichkeit setzt die Anzahl gemeinsamer Werte ins Verhältnis zur Anzahl aller unterschiedlichen Werte: \[ J(A,B) = \frac{|A\cap B|}{|A\cup B|}. \] 
% % Für das Motorbeispiel ergibt sich: \[ J(A,B) = \frac{1}{3} \approx 0{,}33. \] 
% % Der Wert zeigt, dass die angefragte Spannung im Makroprofil enthalten ist, das Makroprofil jedoch zusätzlich zwei weitere Spannungswerte besitzt. 
% % Ein Wert von \(0\) bedeutet, dass keine gemeinsamen Werte vorhanden sind. 
% % Ein Wert von \(1\) bedeutet, dass beide Mengen identisch sind. 
% % Der Overlap-Koeffizient setzt die Anzahl gemeinsamer Werte dagegen ins Verhältnis zur Größe der kleineren Menge: \[ O(A,B) = \frac{|A\cap B|} {\min(|A|,|B|)}. \] 
% % Damit ergibt sich: \[ O(A,B) = \frac{1}{\min(1,3)} = 1. \] 
% % Der Overlap-Koeffizient beträgt \(1\), weil die kleinere Menge \(\{400\,V\}\) vollständig im Makroprofil enthalten ist. 
% % Im Gegensatz zur Jaccard-Ähnlichkeit wird der Wert nicht dadurch verringert, dass das Makroprofil weitere Spannungswerte enthält. 

% % \paragraph{Anfrage- und Kandidatenabdeckung} 

% % Die Anfrageabdeckung beschreibt, welcher Anteil der im Anfrageprofil vorhandenen Werte durch das Makroprofil abgedeckt wird: \[ C_{\mathrm{Anfrage}}(A,B) = \frac{|A\cap B|}{|A|}. \] 
% % Für die betrachteten Spannungsmengen ergibt sich: \[ C_{\mathrm{Anfrage}}(A,B) = \frac{1}{1} = 1. \] Die angefragte Spannung wird somit vollständig durch das Makroprofil abgedeckt. 
% % Die Kandidatenabdeckung beschreibt dagegen, welcher Anteil der im Makroprofil vorhandenen Werte ebenfalls im Anfrageprofil vorkommt: \[ C_{\mathrm{Kandidat}}(A,B) = \frac{|A\cap B|}{|B|}. \] 
% % Für das Beispiel ergibt sich: \[ C_{\mathrm{Kandidat}}(A,B) = \frac{1}{3} \approx 0{,}33. \] 
% % Anfrageabdeckung und Kandidatenabdeckung können somit unterschiedliche Werte annehmen. 
% % Dies ist erforderlich, da Anfrageprofil und Makroprofil unterschiedliche Rollen besitzen. 
% % Eine Anfrage kann vollständig durch ein Makroprofil abgedeckt werden, obwohl das Makroprofil zusätzliche technische Ausprägungen enthält. 

% % \paragraph{Numerische Vergleichsmerkmale} 

% % Numerische Vergleichsmerkmale beschreiben die Distanz zwischen technischen Zahlenwerten, beispielsweise zwischen Spannungen, Strömen, Leistungen oder Frequenzen. 
% % Enthält das Makroprofil mehrere mögliche Werte, kann die geringste absolute Differenz zum Anfragewert verwendet werden: \[ D(a,B) = \min_{b\in B}|a-b|. \] 
% % Für den Anfragewert \[ a=400\,V \] und die Spannungswerte des Makroprofils \[ B=\{380\,V,\;400\,V,\;420\,V\} \] 
% % ergibt sich: \[ D(400\,V,B) = \min \left( |400-380|, |400-400|, |400-420| \right) = 0\,V. \] 
% % Die Differenz von \(0\,V\) kennzeichnet eine exakte numerische Übereinstimmung. 
% % Würde das Makroprofil lediglich die Spannungen \(380\,V\) und \(420\,V\) enthalten, ergäbe sich: \[ D(400\,V,\{380\,V,420\,V\}) = \min(20\,V,20\,V) = 20\,V. \] 
% % Das Differenzmerkmal zeigt damit nicht nur, ob ein Wert exakt enthalten ist, sondern auch, wie weit der nächstgelegene Wert des Makroprofils vom Anfragewert entfernt liegt. 

% % \paragraph{Lexikalische Vergleichsmerkmale} 

% % Lexikalische Vergleichsmerkmale berücksichtigen Übereinstimmungen zwischen Typbezeichnungen, Produktserien, Beschreibungen und \textit{technischen Hinweisen}. 
% % Im Motorbeispiel ist der \textit{technische Hinweis} ``Wasserdichte Ausführung'' sowohl im Anfrageprofil als auch im Makroprofil enthalten. 
% % Daraus kann ein binäres Vergleichsmerkmal mit dem Wert \(1\) gebildet werden. Bei mehreren textuellen Merkmalen können zusätzlich die zuvor beschriebenen mengenbasierten Maße eingesetzt werden. 
% % Dadurch lässt sich beispielsweise bestimmen, welcher Anteil der \textit{technischen Hinweise} einer Anfrage auch im Makroprofil enthalten ist. 

% % \paragraph{Makrostrukturelle Vergleichsmerkmale} 

% % Makrostrukturelle Vergleichsmerkmale beschreiben die Übereinstimmung von Anschlussstrukturen, Makrokategorien, Familien und Varianten. 
% % Im Motorbeispiel enthalten Anfrageprofil und Makroprofil die Anschlussstruktur \[ \{\text{U1},\text{V1},\text{W1}\}. \] 
% % Da beide Anschlussmengen identisch sind, ergeben sowohl die Merkmalsübereinstimmung als auch die Jaccard-Ähnlichkeit, der Overlap-Koeffizient und beide Abdeckungsmaße den Wert \(1\). 
% % Die Makrobezeichnung wird hiervon abgegrenzt. Sie kennzeichnet den Makrokandidaten, stellt jedoch kein zulässiges Vergleichsmerkmal dar, wenn im Anfrageprofil kein entsprechender Vergleichspartner vorhanden ist. 

% % \paragraph{Aggregierte Vergleichsmerkmale} 

% % Zusätzlich können mehrere Einzelvergleiche zu aggregierten Vergleichsmerkmalen zusammengefasst werden. Hierzu zählen beispielsweise die Anzahl gemeinsamer Merkmale, 
% % der Anteil übereinstimmender Eigenschaften oder die gesamte Anfrageabdeckung. 
% % Die Anzahl der Merkmalsbereiche mit mindestens einer Übereinstimmung kann wie folgt beschrieben werden: \[ N_{\mathrm{gemeinsam}} = \sum_{i=1}^{d} \mathbb{1}(A_i\cap B_i\neq\emptyset). \] 
% % Die Funktion \(\mathbb{1}\) liefert den Wert \(1\), wenn im jeweiligen Merkmalsbereich mindestens eine Übereinstimmung vorhanden ist, und andernfalls den Wert \(0\). 
% % Stimmen im Motorbeispiel Hersteller, Produkttyp, Spannung, Strom, Anschlussstruktur und \textit{technischer Hinweis} überein, ergibt sich: \[ N_{\mathrm{gemeinsam}}=6. \] 
% % Aggregierte Vergleichsmerkmale ersetzen die einzelnen Vergleichsmerkmale nicht. Sie ergänzen diese um eine zusammenfassende Beschreibung der technischen Gesamtübereinstimmung. 

% % Tabelle \ref{tab:vergleichsmerkmale_beispiel} fasst die am Motorbeispiel gebildeten Vergleichsmerkmale zusammen. 

% % \begin{table}[H] 
% %     \centering 
% %     \caption{Beispielhafte Bildung von Vergleichsmerkmalen zwischen Anfrageprofil und Makroprofil} 
% %     \label{tab:vergleichsmerkmale_beispiel} 
% %     \renewcommand{\arraystretch}{1.15} 
% %     \begin{tabular}{p{4.0cm}p{6.8cm}p{1.5cm}} 
% %         \hline \textbf{Vergleichsmerkmal} & \textbf{Bedeutung im Motorbeispiel} & \textbf{Wert} \\ 
% %         \hline 
% %         Hersteller-Match & Siemens ist im Makroprofil enthalten. & \(1\) \\ 
% %         Produkttyp-Match & Der Produkttyp Elektromotor stimmt überein. & \(1\) \\ 
% %         Jaccard-Ähnlichkeit der Spannungen & Ein gemeinsamer Wert bei insgesamt drei unterschiedlichen Spannungswerten. & \(0{,}33\) \\ 
% %         Overlap-Koeffizient der Spannungen & Die Spannungsmenge der Anfrage ist vollständig im Makroprofil enthalten. & \(1{,}00\) \\ 
% %         Anfrageabdeckung der Spannungen & Der Anfragewert \(400\,V\) ist vollständig abgedeckt. & \(1{,}00\) \\ 
% %         Kandidatenabdeckung der Spannungen & Einer von drei Spannungswerten des Makroprofils kommt in der Anfrage vor. & \(0{,}33\) \\ 
% %         Geringste Spannungsdifferenz & \(400\,V\) ist im Makroprofil exakt enthalten. & \(0\,V\) \\ 
% %         Lexikalischer Match & Der \textit{technische Hinweis} ``Wasserdichte Ausführung'' stimmt überein. & \(1\) \\ 
% %         Anschlussstruktur-Match & U1, V1 und W1 stimmen vollständig überein. & \(1\) \\ 
% %         Anzahl gemeinsamer Merkmale & Sechs betrachtete Merkmalsbereiche besitzen mindestens eine Übereinstimmung. & \(6\) \\ 
% %         \hline 
% %     \end{tabular} 
% % \end{table} 

% % Die Gesamtheit der für ein Anfrageprofil und ein Makroprofil berechneten Vergleichsmerkmale wird in einer festen Reihenfolge zu einem Merkmalsvektor zusammengeführt. 
% % Für einen Ausschnitt der im Motorbeispiel berechneten Vergleichsmerkmale kann dieser vereinfacht wie folgt dargestellt werden: 
% % \[ \mathbf{x}_{q,m} = \begin{bmatrix} 1 & 1 & 0{,}33 & 1{,}00 & 1{,}00 & 0{,}33 & 0 & 1 & 1 & 6 \end{bmatrix}. \] 

% % Dabei bezeichnet \(q\) das Anfrageprofil und \(m\) das Makroprofil eines Makrokandidaten. Jede Position besitzt bei allen Anfrage-Kandidat-Paaren dieselbe Bedeutung. 
% % Der erste Wert beschreibt beispielsweise stets den Hersteller-Match, während die weiteren Positionen die jeweils festgelegten mengenbasierten, numerischen, lexikalischen, 
% % makrostrukturellen und aggregierten Vergleichsmerkmale enthalten. Der Merkmalsvektor bildet die Grundlage für das zugehörige Anfrage-Kandidat-Paar. 
% % LambdaMART verarbeitet für eine Anfrage die Merkmalsvektoren aller verfügbaren Makrokandidaten und lernt anhand der bekannten Artikel-Makro-Zuordnungen, welche Übereinstimmungen, 
% % Unterschiede und Kombinationen von Vergleichsmerkmalen für die Rangfolge der Makrokandidaten relevant sind.


% % \section{Trainingsmatrix} 

% % Training: Training mit n Artikeln zu 1nem Makro (n:1-Beziehung)

% % Nach der Erstellung der Kandidatenprofile werden für jede Anfrage die verfügbaren Makrokandidaten bewertet. 
% % Für jede Anfrage werden Anfrage-Kandidat-Paare erzeugt und anhand der Pairwise-Merkmale aus Kapitel \ref{chap:analyse_der_datenstrukturen} beschrieben. 
% % Die Referenzkonfiguration A08 verwendet insgesamt 375 Kandidatenprofile und 2.381 Merkmale. 
% % Für jede Anfrage werden alle 375 Kandidatenprofile vollständig bewertet. 

% % Tabelle \ref{tab:matrixdimensionen} fasst die Größenordnung der erzeugten Rankingmatrizen zusammen.

% % \begin{table}[H]
% %     \centering
% %     \setlength{\tabcolsep}{5pt}
% %     \renewcommand{\arraystretch}{1.08}
% %     \caption{Dimensionen der erzeugten Rankingmatrizen}
% %     \label{tab:matrixdimensionen}
% %     \begin{tabular}{lrrrr}
% %         \toprule
% %         \textbf{Split} & \textbf{Anfragen} & \textbf{Kandidaten je Anfrage} & \textbf{Matrixzeilen} & \textbf{Merkmale} \\
% %         \midrule
% %         Training & 984 & 375 & 369.000 & 2.381 \\
% %         Validation & 322 & 375 & 120.750 & 2.381 \\
% %         Test & 253 & 375 & 94.875 & 2.381 \\
% %         \bottomrule
% %     \end{tabular}
% % \end{table}

% % Der erzeugte Artikel-Kandidat-Merkmalsraum ist sehr groß und viele Merkmalswerte besitzen den Wert null. Aus diesem Grund werden Sparse-Matrizen verwendet. Bei einer Sparse-Matrix werden nicht belegte Werte nicht vollständig gespeichert, wodurch der Speicherbedarf reduziert wird. Die numerischen Werte werden im 32-Bit-Gleitkommaformat \texttt{float32} verarbeitet. 


% % \section{Hyperparameteroptimierung} 

% % \section{Hyperparameteroptimierung}

% % Für die Hyperparameteroptimierung wird eine Grid-Search mit insgesamt 15 definierten Parameterkonfigurationen verwendet. Ein Grid-Search ist die systematische Erprobung von zuvor festgelegten Kombinationen an Modellparametern. Jede Kombination wird dabei trainiert und anhand des Validierungsbestands bewertet; anschließend wird anhand definierter Rankingmetriken die beste Konfiguration ausgewählt. 
% % Variiert werden dabei folgende Parameter: 

% % \begin{itemize} 
% %     \item \texttt{n\_estimators}, 
% %     \item \texttt{learning\_rate}, 
% %     \item \texttt{num\_leaves}, 
% %     \item \texttt{min\_child\_samples}, 
% %     \item \texttt{lambdarank\_truncation\_level}. 
% % \end{itemize} 


% % Die Tabelle \ref{tab:hyperparameter} zeigt die untersuchten Hyperparameterkombinationen.

% % \begin{table}[H] 
% %     \centering 
% %     \setlength{\tabcolsep}{8pt}
% %     \renewcommand{\arraystretch}{1.1}
% %     \caption{Variierte Hyperparameter der Grid-Search \parencite{srcStep931}}
% %     \label{tab:hyperparameter} 
% %     \begin{tabularx}{0.85\textwidth}{>{\raggedright\arraybackslash}p{5.5cm} >{\raggedright\arraybackslash}X} 
% %         \toprule 
% %         \textbf{Parameter} & \textbf{Wertebereich} \\ 
% %         \midrule 
% %         \texttt{n\_estimators} & 300 bis 900 \\ 
% %         \texttt{learning\_rate} & 0.025 bis 0.05 \\ 
% %         \texttt{num\_leaves} & 15 bis 127 \\ 
% %         \texttt{min\_child\_samples} & 5 bis 80 \\ 
% %         \texttt{lambdarank\_truncation\_level} & 5 bis 20 \\
% %         \bottomrule 
% %     \end{tabularx} 
% % \end{table}

% % Dabei sind die Hyperparameter wie folgt zu verstehen:

% % \begin{itemize}
% %     \item \textbf{\texttt{n\_estimators}:} Anzahl der aufgebauten Entscheidungsbäume im Modell. Mehr Bäume = komplexere Zusammenhänge, aber mehr Rechenaufwand.
% %     \item \textbf{\texttt{learning\_rate}:} Lernrate pro Trainingsschritt. Kleinere Werte führen zu vorsichtigerem, aber stabilerem Lernen.
% %     \item \textbf{\texttt{num\_leaves}:} Maximale Anzahl der Blätter pro Baum. Mehr Blätter führen zu komplexeren Bäumen und erhöhen dadurch das Risiko von Overfitting. 
% %     \item \textbf{\texttt{min\_child\_samples}:} Minimale Anzahl an Trainingsbeispielen in einem Blatt. Kleinere Werte = mehr Detailtreue pro Blatt, größere Werte = bessere Regularisierung (verhindert Overfitting).
% %     \item \textbf{\texttt{lambdarank\_\allowbreak truncation\_level}:} Begrenzung des Ranglistenbereichs, der durch LambdaRank optimiert werden soll (Top 10, Top 20).
% % \end{itemize}



% % \section{Training} 


% % Das zentrale Element dieser Architektur basiert auf LightGBM. LightGBM ist eine Bibliothek für Gradient Boosting mithilfe von Entscheidungsbäumen. Dabei werden sequentiell Entscheidungsbäume erzeugt, welche die Fehler des bisherigen Modells schrittweise korrigieren. Für Rankingaufgaben ist der \texttt{LGBMRanker} besonders geeignet. Dieser ist auf Rankinggruppen sowie relative Ergebnisreihenfolgen spezialisiert und verwendet in diesem Projekt das LambdaRank-Verfahren \parencite{LightGBM}.
% % Dieser basiert auf dem in Kapitel \ref{chap:stand_der_technik} beschriebenen LambdaRank-/LambdaMART-Prinzip.

% % Das Modell wird mit den Parametern 

% % \begin{align} 
% %     \text{objective} = \text{lambdarank} \qquad \text{und} \qquad \text{metric} = \text{ndcg}
% % \end{align} 

% % trainiert \parencite{srcStep931}.

% % Der Parameter \texttt{objective} legt das Optimierungsziel des Modells fest. Mit \texttt{lambdarank} wird die Optimierung der Kandidatenreihenfolge durch LambdaRank aktiviert. Der Parameter \texttt{metric} bestimmt die verwendete Rankingmetrik der Bewertung; \texttt{ndcg} steht für Normalized Discounted Cumulative Gain und wird im Kapitel \ref{chap:stand_der_technik} näher erläutert.

% % Die für die spätere Referenzkonfiguration A08 ausgewählten Modellparameter lauten: 

% % \begin{align} 
% %     n\_estimators &= 900 \\ 
% %     learning\_rate &= 0.025 \\ 
% %     num\_leaves &= 15 \\ 
% %     min\_child\_samples &= 50 \\ 
% %     lambdarank\_truncation\_level &= 5 
% % \end{align}

% % Zusätzlich wird eine Kandidatenprofilbreite von 16 Artikeln pro Makro verwendet. 

% % Wie in Kapitel \ref{chap:analyse_der_datenstrukturen} beschrieben, werden Kandidaten einer Anfrage als gemeinsame Rankinggruppe betrachtet.
% % Entsprechend entsteht für jede Anfrage eine Gruppe:

% % \begin{align} 
% %     G_q = \{(q,c_1),(q,c_2),\dots,(q,c_n)\} 
% % \end{align} 

% % mit der Anfrage $q$ und den zugehörigen Makrokandidaten $c_i$.

% % Dadurch erlernt das Modell nicht die allgemeine Zugehörigkeit eines einzelnen Kandidaten zu einer Klasse. Stattdessen wird die optimale Kandidatenreihenfolge innerhalb derselben Anfrage gelernt.
% % Ziel dabei ist es, dass der relevanteste Kandidat möglichst weit oben in der Rangliste erscheint.

% % Die Eingaben für das Modell sind die im Kapitel \ref{chap:analyse_der_datenstrukturen} beschriebenen Pairwise-Merkmale.
% % Das Training erfolgt dabei ausschließlich auf den Trainingsdaten, während die Optimierung der Hyperparameter über die Validierungsdaten erfolgt \parencite{srcStep931}.



% % \section{Kapitelergebnis} 

% % Artefakt A4:
% % Trainiertes Empfehlungssystem
